%% ============================================================================
%% Shared preamble. Governed by /home/renaud/Models/NOTATION_CONTRACT.md.
%% Supplies the document class, packages, theorem environments, and the
%% mandated notation macros; raw retypings of contract symbols are contract
%% violations.
%% ============================================================================

\documentclass[11pt]{book}

%% ---------------------------------------------------------------- packages
\usepackage[T1]{fontenc}
\usepackage[utf8]{inputenc}
\usepackage{amsmath}
\usepackage{amssymb}          % \varkappa, blackboard bold
\usepackage{amsthm}
\usepackage{mathtools}        % \DeclarePairedDelimiter, \coloneqq
\usepackage[margin=1in]{geometry}
\usepackage{graphicx}
\usepackage{xcolor}
\usepackage{tikz}
\usepackage{pgfplots}
\pgfplotsset{compat=newest}
\usepackage{booktabs}
\usepackage{enumitem}

%% ------------------------------------------------- theorem-like environments
%% Numbered within chapter; lemma/proposition/etc. share the theorem counter.
\theoremstyle{plain}
\newtheorem{theorem}{Theorem}[chapter]
\newtheorem{lemma}[theorem]{Lemma}
\newtheorem{proposition}[theorem]{Proposition}
\newtheorem{corollary}[theorem]{Corollary}
\theoremstyle{definition}
\newtheorem{definition}[theorem]{Definition}
\newtheorem{example}[theorem]{Example}
\newtheorem{assumption}[theorem]{Assumption}
\theoremstyle{remark}
\newtheorem{remark}[theorem]{Remark}

%% -------------------------------------------------------- bra-ket (mandated)
%% Hand-rolled, braket-package syntax: \braket{\phi(m)|\beta} takes ONE
%% argument with a literal | inside. Deliberately NOT the physics package
%% (whose \braket takes two arguments and which redefines \div, \Re, ...).
\newcommand{\ket}[1]{\mathinner{\lvert #1 \rangle}}
\newcommand{\bra}[1]{\mathinner{\langle #1 \rvert}}
\newcommand{\braket}[1]{\mathinner{\langle #1 \rangle}}
\newcommand{\ketbra}[2]{\ket{#1}\!\bra{#2}}
\newcommand{\featket}{\ket{\phi(\mln)}}          % feature ket |phi(m)>
\newcommand{\parket}{\ket{\beta}}                % parameter ket |beta>

%% ------------------------------------- mandated volatility decomposition
%% m = ln(K/F);  sigma_smile(m,T);  sigma_ATM(l_F,T),  l_F = ln(F/F_ref).
\newcommand{\mln}{m}                             % strike log-moneyness ln(K/F)
\newcommand{\lF}{\ell_F}                         % forward log-level ln(F/F_ref)
\newcommand{\Fref}{F_{\mathrm{ref}}}             % reference forward (anchor)
\newcommand{\ssm}{\sigma_{\mathrm{smile}}}       % smile component, args (m,T)
\newcommand{\satm}{\sigma_{\mathrm{ATM}}}        % ATM component, args (l_F,T)
\newcommand{\satmi}[1]{\sigma_{\mathrm{ATM},#1}} % per-name ATM vol (basket, C1)
\newcommand{\lFi}[1]{\ell_{F,#1}}                % per-name forward log-level ln(F_i/F_{ref,i}) (C1 note)
\newcommand{\Frefi}[1]{F_{\mathrm{ref},#1}}      % per-name reference forward (C1 note)
\newcommand{\volslope}{\mathrm{vol.slope}}       % := d(satm)/d(lF)  (C2)
\newcommand{\sloc}{\sigma_{\mathrm{loc}}}        % Dupire local vol
\newcommand{\sref}{\sigma_{\mathrm{ref}}}        % kernel reference vol

%% ------------------------------------------------------------ normal glyphs
\newcommand{\npdf}{\varphi}                      % standard normal pdf (R5)
\newcommand{\ncdf}{N}                            % standard normal CDF (R5)
\newcommand{\Nrm}{\mathcal{N}}                   % Gaussian LAW N(mu,Sigma); \ncdf reserved for the scalar CDF (R5)

%% --------------------------------------- Vol I Ch 4 (Calibration) local block
\newcommand{\parvec}{\vartheta}                  % calibration parameter vector (I.4); NOT BS theta (Greeks table)
%% \lFi is defined once in the volatility-decomposition block above (C1 note); shared by I.4 and II.1.
\newcommand{\gDurr}{\mathfrak{g}}                % Durrleman butterfly function (I.4); NOT the metadata slot g
\newcommand{\Ninstr}{N_{\mathrm{ins}}}           % instrument count (I.4); NOT the normal CDF N(x)

%% ------------------------------------------------------------------- Greeks
\newcommand{\vega}{\nu}                          % BS vega dC/dsigma (R2)
\newcommand{\Vanna}{\mathrm{Vanna}}
\newcommand{\Volga}{\mathrm{Volga}}
\newcommand{\Gadj}{\Gamma_{\mathrm{adj}}}        % total (smile-adjusted) gamma

%% ------------------------------------------------- multi-name (Vol II) block
\newcommand{\sB}{\sigma_B}                       % basket vol
\newcommand{\sbmean}{\bar\sigma}                 % (sigma1+sigma2)/2
\newcommand{\Cega}{\mathrm{Cega}}                % correlation sensitivity
\newcommand{\VB}{\mathcal{V}_B}                  % basket vega per vol point
\newcommand{\shvega}{\mathrm{ShadowVega}}
\newcommand{\shdelta}{\mathrm{ShadowDelta}}
\newcommand{\stickyrho}{\textnormal{sticky-}\ensuremath{\rho}} % d(rho)=0 convention (R4); works in text and math
\newcommand{\stickycvc}{\textnormal{sticky-CVC}}               % d(xi)=0 convention (R4); works in text and math
\newcommand{\SigmaM}{\Sigma_M}                   % Margrabe volatility
\newcommand{\omM}{\omega_M}                      % Sigma_M*sqrt(T), replaces nu (R2)
\newcommand{\symk}{\varkappa}                    % Margrabe symmetry constant, replaces k (R3.3)
\newcommand{\SchurC}[2]{\mathrm{Schur}_{#1}\!\left(#2\right)} % Schur complement

%% -------------------------------------------------------- calculus & algebra
\newcommand{\dd}{\mathrm{d}}                                  % upright differential
\newcommand{\pd}[2]{\frac{\partial #1}{\partial #2}}          % partial derivative
\newcommand{\pdd}[2]{\frac{\partial^2 #1}{\partial #2^2}}     % second partial
\newcommand{\pdc}[3]{\frac{\partial^2 #1}{\partial #2\,\partial #3}} % cross partial
\newcommand{\td}[2]{\frac{\dd #1}{\dd #2}}                    % total derivative
\newcommand{\atfixed}[2]{\left.#1\right|_{#2}}                % held-fixed bar, e.g. \atfixed{\pd{\SigmaM}{\sigma_1}}{\xi}
\newcommand{\E}{\mathbb{E}}
\newcommand{\Var}{\operatorname{Var}}
\newcommand{\Cov}{\operatorname{Cov}}
\DeclarePairedDelimiter{\abs}{\lvert}{\rvert}
\DeclarePairedDelimiter{\norm}{\lVert}{\rVert}

%% ============================================================================

%% ============================================================================
%% Title and document body
%% ============================================================================
\title{Quantitative Volatility Corpus\\[0.5em]\Large Volume~I: Single-Name Foundations}
\author{Quant Corpus Consolidation Committee}
\date{}

\begin{document}

\frontmatter
\maketitle
\tableofcontents

\mainmatter

%% ============================================================================
%% Vol I, Chapter 1 -- Foundations: gamma P&L, theta--gamma, options as
%% volatility bets. Conforms to NOTATION_CONTRACT.md; macros from
%% build/preamble.tex only. All worked numbers computed under the standing
%% zero-rate assumption (Assumption 1.1).
%% ============================================================================

\chapter{Gamma P\&L and the Volatility Bet}
\label{ch:foundations}

The value of a European option is a convex function of the underlying price. Convexity makes the position more than a linear exposure: its profit and loss over a price move carries a quadratic term, and holding that quadratic term costs time decay. The balance between the two turns an option into a bet on realized volatility. This chapter derives the balance exactly.

\section{The Black--Scholes frame}
\label{sec:bs-frame}

\begin{assumption}[Standing assumption]
\label{ass:zero-rates}
Rates and dividends are zero. The forward price $F$ is used generally; every collapse $F = S$ is flagged at its point of use.
\end{assumption}

The market variables are the spot price $S$, the forward $F$, the strike $K$, and the time to expiry $T$. The strike log-moneyness is
\begin{equation}
\mln = \ln(K/F).
\label{eq:moneyness}
\end{equation}
A European call with implied volatility $\sigma$ has the Black--Scholes price
\begin{equation}
C(K,T) = F\,\ncdf(d_1) - K\,\ncdf(d_2),
\qquad
d_1 = -\frac{\mln}{\sigma\sqrt{T}} + \frac{\sigma\sqrt{T}}{2},
\qquad
d_2 = d_1 - \sigma\sqrt{T},
\label{eq:bs-price}
\end{equation}
where $\ncdf$ is the standard normal CDF and $\npdf$ the standard normal pdf. Under Assumption~\ref{ass:zero-rates}, $F = S$, and \eqref{eq:bs-price} reads $C = S\,\ncdf(d_1) - K\,\ncdf(d_2)$.

\begin{assumption}[Fixed implied volatility]
\label{ass:fixed-sigma}
Throughout this chapter $\sigma$ is the fixed Black--Scholes input for a single option: it moves with neither spot nor time. Its dependence on $\mln$, on the forward level, and on $T$, and the P\&L that its movement generates, are deferred to Chapters~2--3.
\end{assumption}

\begin{definition}[Delta, gamma, theta]
\label{def:greeks}
\begin{equation}
\Delta = \pd{C}{S}, \qquad
\Gamma = \pdd{C}{S} = \pd{\Delta}{S}, \qquad
\theta = -\pd{C}{T}.
\end{equation}
$\Delta$ is the sensitivity of the option price to the underlying, $\Gamma$ the sensitivity of $\Delta$ itself, and $\theta$ the drift of the price as calendar time advances ($T$ decreasing). Since \eqref{eq:bs-price} expresses $C$ through $F$ alone, the $S$-derivatives carry the chain rule through $F(S)$: $\Delta = (\partial C/\partial F)(\dd F/\dd S)$, and correspondingly for $\Gamma$. Under Assumption~\ref{ass:zero-rates}, $\dd F/\dd S = 1$ ($F = S$, flagged), so the $S$-partials coincide with the $F$-partials.
\end{definition}

\begin{lemma}[Closed forms at zero rates]
\label{lem:closed-forms}
Under Assumption~\ref{ass:zero-rates} ($F = S$),
\begin{equation}
\Delta = \ncdf(d_1), \qquad
\Gamma = \frac{\npdf(d_1)}{S\sigma\sqrt{T}}, \qquad
\theta = -\frac{S\,\npdf(d_1)\,\sigma}{2\sqrt{T}}.
\label{eq:closed-forms}
\end{equation}
\end{lemma}

\begin{proof}
Direct differentiation of \eqref{eq:bs-price}; the terms in $\partial d_1/\partial S$ and $\partial d_2/\partial S$ cancel by the identity $F\,\npdf(d_1) = K\,\npdf(d_2)$. The identity is one line: $d_1^2 - d_2^2 = (d_1 + d_2)(d_1 - d_2) = -2\mln$, so $\npdf(d_1)/\npdf(d_2) = e^{\mln} = K/F$. Under Assumption~\ref{ass:zero-rates} ($F = S$) it reads $S\,\npdf(d_1) = K\,\npdf(d_2)$.
\end{proof}

Since $\npdf > 0$, $\Gamma > 0$: the call price is strictly convex in $S$, and $\theta < 0$: a long option loses value as time passes. These two signs are the entire subject of the chapter. Figure~\ref{fig:convexity} shows the geometry: the tangent line at the current spot is the delta approximation, and the gap between curve and tangent grows quadratically with the size of the move.

\begin{figure}[t]
\centering
\begin{tikzpicture}
\begin{axis}[
    width=0.72\textwidth, height=0.48\textwidth,
    axis lines=left,
    xlabel={$S$}, ylabel={$C(S)$},
    xtick={100,110}, xticklabels={$S_1$,$S_2$},
    ytick=\empty,
    xmin=87, xmax=117, ymin=0, ymax=18,
    clip=false,
]
\addplot[domain=88:115, samples=120, thick] {0.5*((x-100)+sqrt((x-100)^2+120))};
\addplot[domain=90:115, dashed] {5.4772+0.5*(x-100)};
\addplot[only marks, mark=*, mark size=1.6pt] coordinates {(100,5.4772) (110,12.4164) (110,10.4772)};
\draw[dotted] (axis cs:110,0) -- (axis cs:110,12.4164);
\draw[<->] (axis cs:110.9,10.4772) -- (axis cs:110.9,12.4164)
    node[midway, right] {$\tfrac12\Gamma\,(\delta S)^2$};
\node[below left] at (axis cs:107.5,8.6) {slope $\Delta$};
\end{axis}
\end{tikzpicture}
\caption{Convexity of the call price (schematic). The dashed tangent at $S_1$ is the delta approximation; over the move $\delta S = S_2 - S_1$ the option gains the quadratic gap above the tangent. For a long position the gap is favorable in both directions of the move.}
\label{fig:convexity}
\end{figure}

\section{The P\&L of a price move}
\label{sec:gamma-pnl}

At fixed time, a move of the underlying from $S_1$ to $S_2 = S_1 + \delta S$ changes the option price by
\begin{equation}
C(S_2) - C(S_1) = \int_{S_1}^{S_2} \Delta(u)\,\dd u,
\label{eq:exact-integral}
\end{equation}
the cumulative effect of delta along the path.

\begin{theorem}[Gamma P\&L formula]
\label{thm:gamma-pnl}
For a move $\delta S$ at fixed time,
\begin{equation}
\boxed{\;\mathrm{P\&L} \;=\; \Delta\cdot\delta S \;+\; \tfrac12\,\Gamma\,(\delta S)^2\;}
\label{eq:gamma-pnl}
\end{equation}
with error $O\!\left((\delta S)^3\right)$; $\Delta$ and $\Gamma$ are evaluated at $S_1$.
\end{theorem}

\begin{proof}
Second-order Taylor expansion of $C$ in $S$ at fixed time, with Definition~\ref{def:greeks}.
\end{proof}

The first term is the delta P\&L, linear in the move; the second is the gamma P\&L, quadratic in the move and positive for a long position in either direction. Two further readings of \eqref{eq:gamma-pnl} give the same expression and locate its accuracy.

\begin{proposition}[Average delta and the trapezoidal rule]
\label{prop:three-readings}
Under the linearization $\Delta(S_2) \approx \Delta(S_1) + \Gamma\,\delta S$, the right-hand side of \eqref{eq:gamma-pnl} equals
\begin{enumerate}[label=(\roman*)]
\item the average-delta P\&L $\bar\Delta\cdot\delta S$, where $\bar\Delta = \tfrac12\left[\Delta(S_1)+\Delta(S_2)\right]$;
\item the trapezoidal-rule approximation of the exact integral \eqref{eq:exact-integral}.
\end{enumerate}
\end{proposition}

\begin{proof}
The trapezoidal rule applied to \eqref{eq:exact-integral} is exactly $\bar\Delta\cdot\delta S$, so (i) and (ii) are the same quantity. Substituting the linearization,
\begin{equation}
\bar\Delta\cdot\delta S
= \frac{\Delta + (\Delta + \Gamma\,\delta S)}{2}\,\delta S
= \left(\Delta + \frac{\Gamma\,\delta S}{2}\right)\delta S
= \Delta\cdot\delta S + \tfrac12\,\Gamma\,(\delta S)^2 .
\qedhere
\end{equation}
\end{proof}

\begin{figure}[t]
\centering
\begin{tikzpicture}
\begin{axis}[
    width=0.72\textwidth, height=0.48\textwidth,
    axis lines=left,
    xlabel={$u$}, ylabel={$\Delta(u)$},
    xtick={100,108}, xticklabels={$S_1$,$S_2$},
    ytick={0.5,0.73106}, yticklabels={$\Delta(S_1)$,$\Delta(S_2)$},
    xmin=90, xmax=114, ymin=0, ymax=1,
]
\fill[gray!25] (axis cs:100,0) -- (axis cs:100,0.5) -- (axis cs:108,0.73106) -- (axis cs:108,0) -- cycle;
\addplot[domain=90:114, samples=90, thick] {1/(1+exp(-(x-100)/8))};
\draw[dashed] (axis cs:100,0.5) -- (axis cs:108,0.73106);
\addplot[only marks, mark=*, mark size=1.6pt] coordinates {(100,0.5) (108,0.73106)};
\end{axis}
\end{tikzpicture}
\caption{The trapezoidal reading of the gamma P\&L formula (schematic). The exact P\&L \eqref{eq:exact-integral} is the area under the delta curve between $S_1$ and $S_2$; the shaded trapezoid, with base $\delta S$ and average height $\bar\Delta$, approximates it and equals $\Delta\cdot\delta S + \tfrac12\Gamma(\delta S)^2$. The error is the sliver between the chord and the curve.}
\label{fig:trapezoid}
\end{figure}

\section{A worked example}
\label{sec:worked-example}

\begin{example}[At-the-money call, zero rates]
\label{ex:atm-call}
An at-the-money call under Assumption~\ref{ass:zero-rates} ($F = S$, so $K = S$ gives $\mln = 0$):
\begin{itemize}
\item spot $S_1 = \$100$, strike $K = \$100$;
\item time to expiry $T = 0.25$ (three months);
\item implied volatility $\sigma = 30\%$;
\item move $\delta S = \$2$, to $S_2 = \$102$.
\end{itemize}
Evaluating \eqref{eq:bs-price} and \eqref{eq:closed-forms} ($d_1 = 0.075$, $d_2 = -0.075$ at $S_1$):
\begin{equation}
C(100) = 5.9785, \quad
C(102) = 7.0908, \quad
\Delta(100) = 0.5299, \quad
\Gamma(100) = 0.0265 .
\end{equation}
\end{example}

Table~\ref{tab:pnl-analysis} decomposes the P\&L of the move.

\begin{table}[t]
\centering
\begin{tabular}{@{}lccc@{}}
\toprule
\textbf{Component} & \textbf{Formula} & \textbf{Calculation} & \textbf{Value (\$)} \\
\midrule
Repricing at constant $\sigma$ & $C(102) - C(100)$ & $7.0908 - 5.9785$ & $1.1122$ \\
Delta P\&L & $\Delta\cdot\delta S$ & $0.5299 \times 2$ & $1.0598$ \\
Gamma P\&L & $\tfrac12\Gamma\,(\delta S)^2$ & $\tfrac12 \times 0.0265 \times 4$ & $0.0530$ \\
\midrule
Approximation \eqref{eq:gamma-pnl} & delta $+$ gamma & $1.0598 + 0.0530$ & $1.1128$ \\
Error & & $1.1128 - 1.1122$ & $0.0006$ \\
Error (relative) & & $0.0006 / 1.1122$ & $0.05\%$ \\
\bottomrule
\end{tabular}
\caption{P\&L decomposition for the \$2 move of Example~\ref{ex:atm-call}. All values are computed under Assumption~\ref{ass:zero-rates} and at fixed $\sigma$ (Assumption~\ref{ass:fixed-sigma}): the first row is the single-option repricing at unchanged implied volatility, not the book P\&L. The surface P\&L generated by $\sigma$ moving with spot is excluded here and enters in Chapters~2--3.}
\label{tab:pnl-analysis}
\end{table}

The delta-only estimate understates the repricing by the gamma term: convexity works for the long position. The quadratic term, \$0.0530, is already material on a 2\% move, and the two-term formula reproduces the constant-$\sigma$ repricing to 0.05\%. That figure belongs to the 2\% move, not to the formula: the error term of \eqref{eq:gamma-pnl} is $O\!\left((\delta S)^3\right)$, so the relative error grows like $(\delta S/S)^2$, and on a large or gapping move the estimate degrades on two counts --- the cubic Taylor error and, chiefly, the excluded repricing of $\sigma$ itself.

\section{Theta and the breakeven move}
\label{sec:theta-gamma}

Time decay is quoted per trading day, with the convention of $252$ trading days per year:
\begin{equation}
\theta_{\mathrm d} = \frac{\theta}{252}.
\label{eq:theta-daily}
\end{equation}
Consider a long option, delta-hedged: the hedge removes the first-order term of \eqref{eq:gamma-pnl}, and under Assumption~\ref{ass:zero-rates} the hedge carries no financing cost. Over one trading day with underlying move $\delta S$, the position's P\&L is, to second order in $\delta S$ and at fixed $\sigma$ (Assumption~\ref{ass:fixed-sigma}, which suppresses the day's $\delta\sigma$ terms),
\begin{equation}
\mathrm{P\&L}_{\mathrm{day}} = \tfrac12\,\Gamma\,(\delta S)^2 + \theta_{\mathrm d}.
\label{eq:daily-pnl}
\end{equation}
The gamma term is a sure gain; the theta term is a sure loss. Their balance point is the breakeven move.

\begin{theorem}[Theta--gamma breakeven]
\label{thm:breakeven}
The daily P\&L \eqref{eq:daily-pnl} vanishes exactly when $\abs{\delta S}$ equals the breakeven move
\begin{equation}
\delta S_{\mathrm{BE}} = \sqrt{\frac{2\,\abs{\theta_{\mathrm d}}}{\Gamma}},
\qquad\text{equivalently}\qquad
\abs{\theta_{\mathrm d}} = \tfrac12\,\Gamma\,(\delta S_{\mathrm{BE}})^2,
\label{eq:breakeven}
\end{equation}
and the daily P\&L takes the form
\begin{equation}
\boxed{\;\mathrm{P\&L}_{\mathrm{day}} \;=\; \tfrac12\,\Gamma\left[(\delta S)^2 - (\delta S_{\mathrm{BE}})^2\right].\;}
\label{eq:gamma-theta-pnl}
\end{equation}
\end{theorem}

\begin{proof}
Set \eqref{eq:daily-pnl} to zero and solve for $\abs{\delta S}$, using $\theta_{\mathrm d} < 0$ for the long option; \eqref{eq:gamma-theta-pnl} is \eqref{eq:daily-pnl} with $\theta_{\mathrm d} = -\tfrac12\Gamma(\delta S_{\mathrm{BE}})^2$ substituted.
\end{proof}

Equation \eqref{eq:gamma-theta-pnl} contains the whole trade: the day is profitable precisely when the underlying moves more than the breakeven, in either direction.

\begin{corollary}[Closed form of the breakeven]
\label{cor:breakeven-closed}
Under Assumption~\ref{ass:zero-rates} ($F = S$), the ratio of the closed forms \eqref{eq:closed-forms} gives $\abs{\theta}/\Gamma = S^2\sigma^2/2$ exactly, so
\begin{equation}
\delta S_{\mathrm{BE}} = \sqrt{\frac{2\abs{\theta}}{252\,\Gamma}} = \frac{\sigma S}{\sqrt{252}},
\label{eq:breakeven-closed}
\end{equation}
independent of strike and of time to expiry: the factor $\npdf(d_1)$ cancels in the ratio.
\end{corollary}

For Example~\ref{ex:atm-call}, $\theta = -11.9347$ per year, so $\theta_{\mathrm d} = -0.0474$ per trading day, and
\begin{equation}
\delta S_{\mathrm{BE}} = \frac{\sigma S}{\sqrt{252}} = \frac{0.30 \times 100}{\sqrt{252}} = \$1.89 .
\end{equation}
Table~\ref{tab:breakeven-analysis} and Figure~\ref{fig:breakeven} evaluate \eqref{eq:gamma-theta-pnl} across daily move sizes.

\begin{table}[t]
\centering
\begin{tabular}{@{}cccc@{}}
\toprule
\textbf{Daily move (\$)} & \textbf{Gamma P\&L (\$)} & \textbf{Theta P\&L (\$)} & \textbf{Total P\&L (\$)} \\
\midrule
$0.50$ & $0.0033$ & $-0.0474$ & $-0.0440$ \\
$1.00$ & $0.0133$ & $-0.0474$ & $-0.0341$ \\
$1.89$ & $0.0474$ & $-0.0474$ & $0.0000$ \\
$2.00$ & $0.0530$ & $-0.0474$ & $0.0057$ \\
$3.00$ & $0.1193$ & $-0.0474$ & $0.0720$ \\
\bottomrule
\end{tabular}
\caption{Daily P\&L of the delta-hedged position of Example~\ref{ex:atm-call} across move sizes, from \eqref{eq:daily-pnl} with $\Gamma = 0.0265$, $\theta_{\mathrm d} = -0.0474$. The breakeven row is $\delta S_{\mathrm{BE}} = \$1.89$. All values are computed under Assumption~\ref{ass:zero-rates} and at fixed $\sigma$ (Assumption~\ref{ass:fixed-sigma}).}
\label{tab:breakeven-analysis}
\end{table}

\begin{figure}[t]
\centering
\begin{tikzpicture}
\begin{axis}[
    width=0.72\textwidth, height=0.48\textwidth,
    axis lines=left,
    xlabel={daily move $\delta S$ (\$)},
    ylabel={daily P\&L (\$)},
    xmin=-3.2, xmax=3.2, ymin=-0.065, ymax=0.09,
    xtick={-3,-1.8898,0,1.8898,3},
    xticklabels={$-3$,$-\delta S_{\mathrm{BE}}$,$0$,$\delta S_{\mathrm{BE}}$,$3$},
    ytick={-0.0474,0,0.05},
    yticklabels={$\theta_{\mathrm d}$,$0$,$0.05$},
]
\addplot[domain=-3.1:3.1, samples=100, thick] {0.5*0.026521*x^2 - 0.047360};
\draw[dashed] (axis cs:-3.2,0) -- (axis cs:3.2,0);
\draw[dotted] (axis cs:-1.8898,-0.065) -- (axis cs:-1.8898,0);
\draw[dotted] (axis cs:1.8898,-0.065) -- (axis cs:1.8898,0);
\addplot[only marks, mark=*, mark size=1.6pt] coordinates {(-1.8898,0) (1.8898,0) (0,-0.047360)};
\end{axis}
\end{tikzpicture}
\caption{Daily P\&L \eqref{eq:gamma-theta-pnl} of the delta-hedged position of Example~\ref{ex:atm-call} as a function of the day's move. The parabola crosses zero at $\pm\delta S_{\mathrm{BE}} = \pm\$1.89$; days inside the breakeven band lose theta faster than gamma earns, days outside earn more. The intercept at $\delta S = 0$ is the full day's decay $\theta_{\mathrm d}$.}
\label{fig:breakeven}
\end{figure}

\section{Options as volatility bets}
\label{sec:vol-bets}

Corollary~\ref{cor:breakeven-closed} identifies the breakeven move: $\sigma S / \sqrt{252}$ is the one-standard-deviation daily move of the underlying under its implied volatility, since for the diffusion $\dd S = \sigma S\,\dd W$ the standard deviation of $\delta S$ over one trading day is $\sigma S \sqrt{1/252}$ to first order. The delta-hedged long option is therefore a bet that the underlying's realized daily moves exceed the one-standard-deviation move priced by the implied volatility. Aggregated over the option's life, the bet is on realized versus implied variance.

\begin{proposition}[Expected P\&L of the volatility bet]
\label{prop:vol-bet}
Denote by $\sigma_{\mathrm{real}}$ the realized volatility of the underlying over the option's life, treated as deterministic (equivalently, every expectation conditions on the realized variance), and by $\sigma_{\mathrm{impl}}$ the implied volatility at which the option was bought. With $\Gamma$ and $S$ held at their initial values, the expected P\&L of the delta-hedged position over the option's life is
\begin{equation}
\E[\mathrm{P\&L}] \approx \tfrac12\,\Gamma\,S^2\left(\sigma_{\mathrm{real}}^2 - \sigma_{\mathrm{impl}}^2\right) T .
\label{eq:vol-bet}
\end{equation}
\end{proposition}

\begin{proof}
Summing the daily P\&L \eqref{eq:gamma-theta-pnl} over the $252\,T$ trading days of the option's life, with $(\delta S_{\mathrm{BE}})^2 = \sigma_{\mathrm{impl}}^2 S^2/252$ from Corollary~\ref{cor:breakeven-closed} and $\E[(\delta S)^2] = \sigma_{\mathrm{real}}^2 S^2/252$ per day,
\begin{equation}
\E[\mathrm{P\&L}] \approx 252\,T \cdot \tfrac12\,\Gamma\,\frac{S^2}{252}\left(\sigma_{\mathrm{real}}^2 - \sigma_{\mathrm{impl}}^2\right)
= \tfrac12\,\Gamma\,S^2\left(\sigma_{\mathrm{real}}^2 - \sigma_{\mathrm{impl}}^2\right) T .
\qedhere
\end{equation}
\end{proof}

The approximation freezes the dollar gamma $\Gamma S^2$ --- and with it the breakeven $\delta S_{\mathrm{BE}}$, which depends on $\Gamma$, $S$, and $\sigma_{\mathrm{impl}}$ through Corollary~\ref{cor:breakeven-closed} --- at initial values; along an actual path both vary with spot and with the passage of time. The exact statement carries the running weight.

\begin{proposition}[Pathwise P\&L of the hedged position]
\label{prop:pathwise-pnl}
Let the underlying follow $\dd S_t = \sigma_t S_t\,\dd W_t$ with $\sigma_t$ the instantaneous realized volatility, and let the long option be delta-hedged continuously at the constant implied volatility $\sigma_{\mathrm{impl}}$. Then, on every path,
\begin{equation}
\mathrm{P\&L} = \tfrac12 \int_0^T \Gamma_t\,S_t^2 \left(\sigma_t^2 - \sigma_{\mathrm{impl}}^2\right) \dd t,
\label{eq:pathwise-pnl}
\end{equation}
where $\Gamma_t$ is the Black--Scholes gamma at $\sigma_{\mathrm{impl}}$ evaluated along the path.
\end{proposition}

\begin{proof}
It\^o's formula on the Black--Scholes value $C(t, S_t; \sigma_{\mathrm{impl}})$ gives $\dd C = \partial_t C\,\dd t + \Delta\,\dd S_t + \tfrac12\,\Gamma_t S_t^2 \sigma_t^2\,\dd t$; the Black--Scholes equation at zero rates gives $\partial_t C = -\tfrac12\,\Gamma_t S_t^2 \sigma_{\mathrm{impl}}^2$. The increment of the hedged position, $\dd C - \Delta\,\dd S_t$, is therefore $\tfrac12\,\Gamma_t S_t^2 (\sigma_t^2 - \sigma_{\mathrm{impl}}^2)\,\dd t$; integrating over $[0,T]$ gives \eqref{eq:pathwise-pnl}. (N.~El Karoui, M.~Jeanblanc-Picqu\'e, and S.~E.~Shreve, ``Robustness of the Black and Scholes formula,'' \emph{Mathematical Finance} 8(2), 1998, 93--126.)
\end{proof}

Since $\Gamma_t > 0$, the sign of the P\&L is guaranteed pathwise whenever $\sigma_t^2 - \sigma_{\mathrm{impl}}^2$ keeps one sign over the option's life. No such guarantee attaches to the aggregate spread $\sigma_{\mathrm{real}}^2 - \sigma_{\mathrm{impl}}^2$: once the dollar-gamma weight in \eqref{eq:pathwise-pnl} is unfrozen, realized variance is collected only where the position has gamma, and the expected P\&L can take the opposite sign to the spread --- as when volatility spikes while spot sits far from the strike and $\Gamma_t$ is negligible. Only the frozen form \eqref{eq:vol-bet} carries the sign of the spread, and over a three-month life the freeze is a leading-order distortion, not a small correction: a gap away from the strike early in the life removes the gamma through which the remaining realized variance would have been collected. Two further omissions of \eqref{eq:vol-bet}: the expectation conceals wide dispersion across paths, so a position right on average --- realized above implied --- still loses on paths that concentrate the variance where its gamma is not; and the formula is gross of hedging frictions, the discrete-hedge error and the transaction costs of rebalancing. A long option position is long realized variance against implied, weighted by where its gamma lives; a short position is the opposite bet.

Everything above prices and differentiates a single option at a single fixed $\sigma$. On a real surface the implied volatility is a function $\sigma(\mln, \lF, T)$ of strike log-moneyness, forward level, and expiry, and its movement generates P\&L terms of its own; that surface, and the sensitivities routed through it, occupy the rest of this volume.
%% ============================================================================
%% Vol I, Chapter 2 -- The Implied and Local Volatility Surface
%% Curvature-kernel model, time dependence, calibration, sensitivities,
%% Dupire local volatility, no-arbitrage / admissibility ladder.
%% Conforms to NOTATION_CONTRACT.md; uses only build/preamble.tex macros.
%% ============================================================================

\chapter{The Implied and Local Volatility Surface}
\label{ch:surface}

The implied-volatility surface is not a collection of independent quotes: it is
a single coherent object constrained, at every strike and every expiry, by the
requirement that no static arbitrage be admitted. This chapter builds a
parsimonious representation of that object, parameterising the smile through
its \emph{curvature}---its second derivative in log-moneyness---rather than
directly, so that calibration collapses to a linear regression on
compactly supported basis functions, the Dupire local volatility follows in
closed form from the same coefficients, and the no-arbitrage conditions become
low-degree polynomial inequalities on those coefficients. The construction
opens the single-name admissibility ladder (wing and convexity bounds) that
carries forward into the multi-name volume.

\medskip
\noindent\textbf{Standing assumption.} The risk-free rate and the dividend
yield are zero, so the forward and spot coincide, $F=S$; the collapse is
flagged at each point of use. The forward $F$ nonetheless remains the primary
variable, so that reintroducing non-zero carry amounts only to restoring
$F = S\,e^{(r-q)T}$ in the moneyness definition.

%% ======================================================================
\section{The Volatility Decomposition and Its Curvature Representation}
\label{sec:implied}
%% ======================================================================

\subsection{The ATM/Smile Decomposition}
\label{sec:decomposition}

The Black--Scholes implied volatility splits into an at-the-money
\emph{level} and a \emph{smile shape},
\begin{equation}\label{eq:decomp}
\sigma(\mln,\lF,T) \;=\; \satm(\lF,T) \;+\; \ssm(\mln,T),
\qquad \ssm(0,T) = 0,
\end{equation}
where the strike log-moneyness is $\mln = \ln(K/F)$ and the forward log-level
is $\lF = \ln(F/\Fref)$, with $\Fref$ a fixed reference forward (the surface
anchor). The ATM component $\satm(\lF,T)$ carries all dependence on where the
forward sits; the smile component $\ssm(\mln,T)$ carries the shape across
strikes at a fixed forward. The normalisation $\ssm(0,T)=0$ is what makes
$\satm$ the genuine at-the-money volatility and makes the split unique: setting
$\mln = 0$ in \eqref{eq:decomp} gives
\begin{equation}\label{eq:atm-recovery}
\sigma(0,\lF,T) = \satm(\lF,T)
\end{equation}
exactly. The whole of the present chapter constructs $\ssm(\mln,T)$; the ATM
level $\satm(\lF,T)$ and its slope in $\lF$ (the $\volslope$) are the objects
that transmit surface dynamics to the routed sensitivities of
Section~\ref{sec:sensitivity} and to the multi-name book. This chapter does
\emph{not} calibrate that $\lF$-dependence: it fits a single smile at the
prevailing forward, so $\volslope$ and $\partial_{\lF}^2\satm$ are exogenous
inputs owned by the surface-dynamics chapter, defaulting to zero (the C4
$\satm=\mathrm{const}$ case) until supplied---a default flagged at every
spot-routed use below.

\begin{remark}[Two slopes, never to be confused]\label{rem:two-slopes}
Two distinct first-order objects live on the surface. The \emph{smile skew}
$s(\mln) = \pd{\ssm}{\mln}$ is the slope of the shape \emph{across strikes} at a
fixed forward. The \emph{$\volslope := \pd{\satm}{\lF}$} is the slope of the
single ATM point \emph{as the forward itself moves}. For equity indices
$\volslope < 0$. These are separate derivatives of separate arguments of
\eqref{eq:decomp}; the routed Greeks of Section~\ref{sec:sensitivity} keep
them apart by construction.
\end{remark}

\subsection{From Curvature to Smile}
\label{sec:curv-to-smile}

Fix the expiry $T$ and the forward. Define the smile \emph{skew} and
\emph{curvature},
\[
s(\mln) = \td{\ssm}{\mln}(\mln), \qquad
c(\mln) = \frac{\dd^2 \ssm}{\dd \mln^2}(\mln).
\]
With the anchoring data $\ssm(0,T) = 0$ (normalisation) and $s(0) = s_0$ (the
ATM smile skew), two integrations recover the shape from its curvature:
\begin{align}
s(\mln)   &= s_0 + \int_0^{\mln} c(u)\,\dd u, \label{eq:skew-integral}\\
\ssm(\mln,T) &= s_0\,\mln + \int_0^{\mln}\!\!\int_0^{u} c(v)\,\dd v\,\dd u.
\label{eq:sigma-integral}
\end{align}
This is the Green's-function solution of $\ssm'' = c(\mln)$ under the boundary
data $(\,\ssm(0,T),\,s(0)\,) = (0,\,s_0)$. The ATM level $\satm(\lF,T)$ enters
\eqref{eq:decomp} additively and is invisible to \eqref{eq:sigma-integral};
this is precisely why calibrating the shape is a linear problem.

\begin{example}[Constant curvature]\label{ex:const-curv}
If $c(\mln) = c_0$ for all $\mln$, then
$\ssm(\mln,T) = s_0\,\mln + \tfrac{1}{2}c_0\,\mln^2$, a parabolic smile with
$\ssm(0,T)=0$. This quadratic smile is revisited as a fully worked local-volatility
example in Section~\ref{sec:quadratic-example}.
\end{example}

\subsection{Kernel Representation of the Curvature}
\label{sec:kernel-rep}

The curvature is a superposition of $n$ compactly supported
Epanechnikov kernels,
\begin{equation}\label{eq:curvature-kernel}
c(\mln) = \sum_{i=1}^n c_i \,\frac{1}{h_i}\,
\kappa\!\left(\frac{\mln - \mu_i}{h_i}\right),
\qquad
\kappa(u) = \begin{cases}
\tfrac{3}{4}(1-u^2), & \abs{u}\le 1,\\[2pt]
0, & \text{otherwise},
\end{cases}
\end{equation}
where $c_i$ are curvature coefficients, $\mu_i$ the kernel centres, and
$h_i>0$ the bandwidths. Because $\int_{-1}^{1}\kappa = 1$, each $c_i$ is the
\emph{total skew change} contributed by kernel $i$; the total skew swing from
$\mln = -\infty$ to $\mln = +\infty$ is $\sum_i c_i$. The choice of the
Epanechnikov kernel is not incidental: it vanishes at its endpoints,
$\kappa(\pm 1) = 0$, which (as Section~\ref{sec:closed} shows) makes the
reconstructed smile globally $C^2$ and uniquely yields closed-form polynomial
basis functions.

\begin{remark}[The flexible-ruler picture]\label{rem:ruler}
A flexible ruler laid along the moneyness axis carries the picture. Two
conditions pin it at the origin: its height (here $0$, by the normalisation)
and its slope $s_0$.
Each curvature kernel $c_i$ is a localised bump placed under the ruler, bending
it within the support $[\mu_i - h_i,\;\mu_i + h_i]$ and leaving it straight
elsewhere. Beyond the last bump on each side, the ruler continues as a straight
line---it has no further reason to bend. The subsections that follow make this
picture precise, and the wing constraints of
Section~\ref{sec:no-arb} are exactly the requirement that these straight tails
be flat rather than sloped.
\end{remark}

\subsubsection*{Function-space structure of the kernel basis}

Let $I_i(\mln)$ denote the double integral of the $i$-th kernel from the origin
(defined precisely in \eqref{eq:Ii} below). The model space
\[
\mathcal{V} = \mathrm{span}\{\,\mln,\,I_1(\mln),\ldots,I_n(\mln)\,\}
\]
of smile shapes decomposes canonically as
\[
\mathcal{V} = \mathcal{V}_0 \;\oplus\; \mathcal{V}_c,
\qquad
\mathcal{V}_0 = \ker \partial_\mln^2 \cap \mathcal{V} = \mathrm{span}\{\mln\},
\qquad
\mathcal{V}_c = \mathrm{span}\{I_1,\ldots,I_n\}.
\]
Here $\mathcal{V}_0$ carries the linear (zero-curvature) part---the ruler with
no bumps, a straight line of slope $s_0$ through the origin---while each basis
function $I_i \in \mathcal{V}_c$ records the deformation produced by one bump.
The Green's-function structure~\eqref{eq:sigma-integral} gives every $I_i$ two
roles: a \emph{local} role (inside its support it contributes genuine
curvature) and a \emph{propagating} role (beyond its support it becomes affine,
collapsing into $\mathcal{V}_0$ and contributing only a permanent change of
slope).

Let
\begin{equation}\label{eq:alpha-def}
\alpha_i = P_\kappa\!\left(-\mu_i/h_i\right)
\end{equation}
denote the fraction of kernel $i$ lying to the left of the origin, where
$P_\kappa$ is the kernel's cumulative (single-integral) function, defined in
\eqref{eq:Pkappa}. Then the single integral satisfies $J_i(\mln)\to 1-\alpha_i$
as $\mln\to+\infty$ and $J_i(\mln)\to-\alpha_i$ as $\mln\to-\infty$, so the skew
change $c_i$ partitions into $c_i(1-\alpha_i)$ propagating rightward and
$-c_i\,\alpha_i$ propagating leftward. Because both $J_i$ and $I_i$ integrate
from $\mln = 0$, one has $J_i(0) = I_i(0) = 0$ for \emph{every} kernel, so
$\ssm(0,T) = 0$ and $s(0) = s_0$ regardless of kernel placement: the anchor at
the origin is immovable. This is exactly the contract normalisation $\ssm(0,T)=0$
holding by construction.

\medskip
\noindent\textit{(i) Central kernel} ($\mu_i = 0$, hence $\alpha_i = \tfrac12$).
The propagation is symmetric: $+c_i/2$ rightward, $-c_i/2$ leftward, and by the
symmetry of $\kappa$ about its centre $I_i(-\mln) = I_i(\mln)$. The central
kernel adds a symmetric dome (or trough, if $c_i < 0$) about the money.

\medskip
\noindent\textit{(ii) Off-centre and outermost kernels}
($\alpha_i \neq \tfrac12$). The skew offset propagates asymmetrically, tilting
one wing more than the other---in the ruler analogy, a bump away from the pin
bends the longer lever arm more. Beyond the last kernel support, $\ssm$ lies
entirely in $\mathcal{V}_0$: the smile is affine, with right-wing slope
$s_0 + \sum_i c_i(1-\alpha_i)$ and left-wing slope $s_0 - \sum_i c_i\alpha_i$
(each divided by $\sqrt{\tilde T}$ after the time scaling of
Section~\ref{sec:time}). Imposing the wing constraints of
Section~\ref{sec:no-arb} forces both slopes to vanish, flattening the smile to
a constant beyond the grid on either side.

\begin{remark}[Grid design]\label{rem:grid}
The grid need not be symmetric: a steep put skew calls for denser kernel
placement on the put side than the call side, with one central kernel at
$\mu_i = 0$. Non-overlapping or minimally overlapping supports give each $c_i$
control of a distinct moneyness region and keep the design matrix $\Phi$ of
Section~\ref{sec:calibration} well conditioned. Since the two wing constraints
consume two degrees of freedom among the $c_i$, a \emph{buffer kernel}
on each side beyond the last liquid strike prevents the constraint solver from
distorting the fit in the data-rich region.
\end{remark}

\subsection{Integral Basis Functions}
\label{sec:integral-basis}

Define the cumulative integrals of the $i$-th kernel from the origin,
\begin{align}
J_i(\mln) &= \int_0^{\mln} \frac{1}{h_i}\,
\kappa\!\left(\frac{u - \mu_i}{h_i}\right)\dd u, \label{eq:Ji}\\
I_i(\mln) &= \int_0^{\mln} J_i(u)\,\dd u. \label{eq:Ii}
\end{align}
The function $J_i$ is a $C^1$ sigmoid transitioning across the support
$[\mu_i - h_i,\;\mu_i + h_i]$, and $I_i$ is its antiderivative; both vanish at
$\mln=0$. (The kernel cumulative $J_i$, \emph{always} carrying a kernel
subscript, is distinct from the bare observation count $J$ of
Section~\ref{sec:calibration}; the two never collide in an unindexed form.)
Their closed forms are collected in Section~\ref{sec:closed}. Unlike
spline interpolation, which imposes junction conditions at knots, the
piecewise-polynomial structure here \emph{arises} from double integration of
the compactly supported $\kappa$; the $C^2$ junctions at $\mu_i \pm h_i$ are
automatic consequences of $\kappa(\pm 1) = 0$, not imposed constraints.

\subsection{The Linear Model and the Bra-Ket Form}
\label{sec:linear-model}

Substituting \eqref{eq:curvature-kernel} into \eqref{eq:sigma-integral} and
restoring the ATM level gives the full implied volatility as a single inner
product. Writing the intercept parameter $\sigma_0$ for the calibrated ATM
level,
\begin{equation}\label{eq:sigma-linear}
\sigma(\mln,\lF,T)
= \sigma_0 + s_0\,\mln + \sum_{i=1}^n c_i\, I_i(\mln)
= \braket{\phi(\mln) | \beta},
\end{equation}
with the feature ket and parameter ket
\[
\ket{\phi(\mln)} = \begin{pmatrix} 1 \\ \mln \\ I_1(\mln) \\ \vdots \\ I_n(\mln)
\end{pmatrix} \in \mathbb{R}^{n+2},
\qquad
\parket = \begin{pmatrix} \sigma_0 \\ s_0 \\ c_1 \\ \vdots \\ c_n \end{pmatrix}
\in \mathbb{R}^{n+2}.
\]
Because $I_i(0)=0$ for every kernel (the \emph{origin-adjusted} convention of
Section~\ref{sec:closed}), the intercept is the ATM level itself,
$\sigma_0 = \satm(\lF,T)$, and the slope is the ATM smile skew, $s_0 = s(0)$;
the smile part $\ssm(\mln,T) = s_0\mln + \sum_i c_i I_i(\mln)$ then satisfies
$\ssm(0,T) = 0$ exactly, as required by \eqref{eq:decomp}.

The calibration is performed at the prevailing forward, with $\lF$
frozen at its current market value, so the right-hand side of
\eqref{eq:sigma-linear} carries no explicit $\lF$ argument and $\sigma_0$
estimates $\satm(\lF,T)$ at that single point; the $\lF$-dependence itself
remains the exogenous input of Section~\ref{sec:decomposition}, consumed by
the routed Greeks of Section~\ref{sec:total-greeks}.

\begin{remark}[Bra-ket convention]\label{rem:braket}
Dirac notation is used with the standard Euclidean inner product on
$\mathbb{R}^{n+2}$: $\ket{\phi(\mln)}$ is the feature ket (column),
$\bra{\phi(\mln)}$ its transpose (row), and $\braket{\phi(\mln)|\beta}$ their dot
product. The parameter ket $\parket$ is the state of the calibrated market;
$\bra{\phi(\mln)}$ is the measurement functional that reads off the implied
volatility at moneyness $\mln$. This bra-ket machinery is reserved throughout
both volumes for the feature/parameter linear algebra; generic inner products
elsewhere are written $x^\top y$.
\end{remark}

\begin{remark}[Derivative kets]\label{rem:derivative-kets}
Differentiation in $\mln$ acts on the feature ket while $\parket$ stays fixed:
\[
\ket{\phi'(\mln)} = \begin{pmatrix} 0 \\ 1 \\ J_1(\mln) \\ \vdots \\ J_n(\mln)
\end{pmatrix},
\qquad
\ket{\phi''(\mln)} = \begin{pmatrix} 0 \\ 0 \\ \kappa_1(\mln)/h_1 \\ \vdots \\
\kappa_n(\mln)/h_n \end{pmatrix},
\qquad \kappa_i(\mln) = \kappa\!\left(\tfrac{\mln-\mu_i}{h_i}\right),
\]
so that $s(\mln) = \braket{\phi'(\mln)|\beta}$ and
$c(\mln) = \braket{\phi''(\mln)|\beta}$. The Dupire denominator of
Section~\ref{sec:localvol} depends on the $2$-jet
$(\ssm,\partial_\mln\ssm,\partial_\mln^2\ssm)
= (\braket{\phi|\beta},\braket{\phi'|\beta},\braket{\phi''|\beta})$, the
numerator on the $1$-jet in $T$. Differentiation is thus an operator on the ket
space that preserves linear dependence on $\parket$: the formalism is
structural, not merely notational.
\end{remark}

\begin{remark}[Origin-adjusted versus cumulative convention]\label{rem:convention}
If one instead builds the basis from the bare kernel cumulative $P_\kappa$
without subtracting its value at the origin (the \emph{cumulative} convention),
kernels whose support straddles or lies left of $\mln=0$ contribute non-zero
offsets $J_i(0), I_i(0)$. The intercept $\sigma_0$ is then a bare regression
coefficient, \emph{not} the ATM level; the ATM level is recovered as
$\satm(\lF,T) = \sigma(0,\lF,T) = \sigma_0 + \sum_i c_i I_i(0)$ and the ATM skew
as $s(0) = s_0 + \sum_i c_i J_i(0)$. The two conventions differ only by
components in $\mathrm{span}\{1,\mln\}$, so ordinary least squares compensates
automatically and the calibrated curve $\sigma(\mln,\lF,T)$---hence the
decomposition $\ssm(\mln,T) = \sigma(\mln,\lF,T) - \satm(\lF,T)$---is identical
either way. The origin-adjusted convention is canonical here precisely
because it makes the contract normalisation $\ssm(0,T)=0$ hold at the level of
the intercept parameter.
\end{remark}

\begin{remark}[Regularisation implication]\label{rem:reg}
Since $\mathcal{V}_0$ carries no curvature and needs no smoothing, a natural
ridge penalty acts only on the curvature coefficients:
$\lambda \sum_{i=1}^n c_i^2$, leaving $\sigma_0$ and $s_0$ unpenalised. The
rationale is again the ruler---the pins at the origin are set by the data; one
wishes only to prevent the bumps from becoming unreasonably large.
\end{remark}

\begin{example}[Three-kernel equity smile]\label{ex:three-kernel}
Take an intercept $\sigma_0 = 0.20$ and slope $s_0 = -0.16$ (both in the
cumulative convention of Remark~\ref{rem:convention}), with three kernels at a
single expiry:

\medskip
\begin{center}
\begin{tabular}{lccc}
\toprule
Kernel & $\mu_i$ & $w_i$ & $c_i$ \\
\midrule
Put wing  & $-1.5\,\sref\sqrt{T}$ & $0.8$ & $+0.60$ \\
ATM       & $0$                    & $0.8$ & $+0.20$ \\
Call wing & $+1.0\,\sref\sqrt{T}$  & $0.8$ & $+0.10$ \\
\bottomrule
\end{tabular}
\end{center}
\medskip

\noindent With $\sref = 0.20$ and $T = 0.25$ the centres are $\mu_1 = -0.15$,
$\mu_2 = 0$, $\mu_3 = +0.10$, all bandwidths $h_i = 0.08$, and the left-of-origin
fractions are $\alpha_1 = 1$, $\alpha_2 = \tfrac12$, $\alpha_3 = 0$. The put-side
kernel bends the skew from $s_0 = -0.16$ up to $s_0 + c_1 = +0.44$ across its
support, flattening the put wing.

\medskip
\noindent\textbf{Reading off the decomposition.} In the cumulative convention
$J_1(0)=1$, $I_1(0)=0.15$, $J_2(0)=0.5$, $I_2(0)=0.015$, $J_3(0)=I_3(0)=0$, so
$\sigma_0$ is \emph{not} the ATM level. The ATM component is recovered as
\[
\satm(\lF,T) = \sigma(0,\lF,T)
= \sigma_0 + \sum_i c_i I_i(0)
= 0.20 + 0.60(0.15) + 0.20(0.015) + 0.10(0) = 0.293,
\]
and the ATM smile skew is
$s(0) = s_0 + \sum_i c_i J_i(0) = -0.16 + 0.60 + 0.10 = 0.54$. The smile
component is then $\ssm(\mln,T) = \sigma(\mln,\lF,T) - 0.293$, which vanishes at
$\mln=0$ by construction. Representative values of the full implied vol:

\medskip
\begin{center}
\begin{tabular}{cccccccc}
\toprule
$\mln$ & $-0.30$ & $-0.20$ & $-0.10$ & $0$ & $+0.10$ & $+0.20$ & $+0.30$ \\
\midrule
$\sigma(\mln,\lF,T)$ & $0.248$ & $0.233$ & $0.247$ & $0.293$ & $0.356$ & $0.428$ & $0.502$ \\
$\ssm(\mln,T)$ & $-0.045$ & $-0.060$ & $-0.046$ & $0$ & $+0.063$ & $+0.135$ & $+0.209$ \\
\bottomrule
\end{tabular}
\end{center}
\medskip

\noindent The smile minimum lies near $\mln = -0.20$, reflecting the interplay
between the negative slope $s_0$ and the positive curvature contributions;
$\ssm$ vanishes at the money, as \eqref{eq:decomp} demands.

\begin{center}
\begin{tikzpicture}
\begin{axis}[
  width=0.86\textwidth, height=6.4cm,
  xlabel={$\mln = \ln(K/F)$}, ylabel={implied volatility},
  xmin=-0.34, xmax=0.34, ymin=0.20, ymax=0.53,
  grid=both, grid style={gray!18},
  legend pos=north west, legend cell align=left,
  tick label style={font=\small}, label style={font=\small},
]
\addplot[thick, mark=*, mark size=1.4pt] coordinates {
  (-0.30,0.248) (-0.20,0.233) (-0.10,0.247) (0,0.293)
  (0.10,0.356) (0.20,0.428) (0.30,0.502) };
\addlegendentry{$\sigma(\mln,\lF,T)$}
\addplot[dashed, thick] coordinates { (-0.34,0.293) (0.34,0.293) };
\addlegendentry{$\satm(\lF,T)=0.293$}
\end{axis}
\end{tikzpicture}
\end{center}
\end{example}

%% ======================================================================
\section{Time Dependence}
\label{sec:time}
%% ======================================================================

\subsection{Square-Root Scaling of the Feature Ket}
\label{sec:sqrt-scaling}

\begin{assumption}[Leading-order skew scaling]\label{ass:scaling}
To leading order the ATM smile skew scales as $1/\sqrt{T}$. This is exact for
surfaces that are functions of the standardised moneyness $\mln/(\sigma\sqrt{T})$
(the sticky-delta self-similar case). It is \emph{not} exact for the general diffusive
stochastic-volatility class: Heston, for instance, has an ATM skew that tends to
a finite constant as $T\to0$ and decays like $1/T$ at long maturity. In
rough-volatility models the empirical scaling is closer to $T^{H-1/2}$ with
Hurst exponent $H\approx 0.1$ (Fukasawa 2011; Bayer, Friz \& Gatheral 2016),
which the framework accommodates by replacing
$\sqrt{\tilde T}=\tilde T^{1/2}$ with $\tilde T^{1/2-H}$ in the feature scaling
below. Dividing the moneyness features by $\tilde T^{1/2-H}$ makes the skew
scale as $\tilde T^{-(1/2-H)}=\tilde T^{H-1/2}$, as required (since $1/2-H>0$ for
$H\approx0.1$, this is a milder division than $1/\sqrt{\tilde T}$, giving a
\emph{flatter} short-dated skew---the opposite of $\tilde T^{H+1/2}$, which
would steepen it), while leaving the regression linear.
\end{assumption}

Under Assumption~\ref{ass:scaling} every moneyness-dependent
component of the feature ket is divided by $\sqrt{\tilde T}$:
\begin{equation}\label{eq:phi-time}
\ket{\phi(\mln,T)} = \begin{pmatrix} 1 \\ \mln/\sqrt{\tilde T} \\
I_1(\mln)/\sqrt{\tilde T} \\ \vdots \\ I_n(\mln)/\sqrt{\tilde T}
\end{pmatrix},
\qquad
\sigma(\mln,\lF,T) = \braket{\phi(\mln,T)|\beta}.
\end{equation}
The intercept row is left untouched: the ATM level carries its own term
structure through $\satm(\lF,T)$, developed in Section~\ref{sec:time-beta}. The
kernel centres and bandwidths scale with $\sqrt{\tilde T}$,
\begin{equation}\label{eq:kernel-scaling}
\mu_i(T) = w_i^{\mathrm{ctr}}\,\sref\sqrt{\tilde T},
\qquad
h_i(T)   = w_i\,\sref\sqrt{\tilde T},
\end{equation}
with $w_i^{\mathrm{ctr}}, w_i$ fixed constants. The reference volatility $\sref$
is held \emph{fixed} during calibration---chosen close to recent ATM vol but
never set to the ATM parameter itself---so that the design matrix depends only
on observable moneyness and time, and the model stays linear in $\parket$.

\subsection{Modified Time to Maturity}
\label{sec:modified-time}

To prevent singular behaviour as $T\to 0$ and excessive flattening for large
$T$, define the modified time $\tilde T$ by
\begin{equation}\label{eq:T-tilde}
T' = \frac{T_{\mathrm{d}} + z}{252 + z},
\qquad
\tilde T = \frac{T'}{1 + T'/T_{\max}} + T_{\min},
\end{equation}
where $T_{\mathrm{d}}$ is time to expiry in \emph{trading days} and
$T' \approx T_{\mathrm{d}}/252$ converts to years (with shift $z$ for
smoothing). Everywhere else in this chapter $T$ denotes time to expiry in
years; \eqref{eq:T-tilde} is the sole appearance of trading-day units. The
constants $T_{\min}, T_{\max}, z$ floor and cap $\tilde T$.

\begin{proposition}[Bounds on $\tilde T$]\label{prop:T-bounds}
For $T\ge 0$, $z>0$, $T_{\min}>0$, $T_{\max}>0$,
\[
T_{\min} < \tilde T < T_{\max} + T_{\min}.
\]
\end{proposition}
\begin{proof}
The map $g(x) = x/(1+x/T_{\max})$ sends $[0,\infty)$ onto $[0,T_{\max})$
strictly monotonically, with supremum $T_{\max}$ never attained. Since
$T' \ge z/(252+z) > 0$ and $T' < \infty$ for finite $T$,
$g(T') \in (0,T_{\max})$; adding $T_{\min}$ gives the claim.
\end{proof}

\begin{remark}[Reference constants]\label{rem:prod-values}
\begin{center}
\begin{tabular}{lcl}
\toprule
Parameter & Value & Rationale \\
\midrule
$T_{\min}$ & $2/252 \approx 0.008$ & floor at $2$ trading days \\
$T_{\max}$ & $5.0$                 & cap at longest liquid maturity \\
$z$        & $2$                   & smooths the $T=0$ singularity \\
\bottomrule
\end{tabular}
\end{center}
\end{remark}

\subsection{Time-Dependent Parameters}
\label{sec:time-beta}

Each component of $\parket$---including the ATM level---is expanded over
time-decay kernels,
\[
\beta_i(T) = \sum_{j} b_{i,j}\,\psi(T;\tau_j,p),
\qquad
\psi(T;\tau,p) = 2^{-(T/\tau)^p} = \exp\!\bigl(-(T/\tau)^p\ln 2\bigr).
\]
At $T=\tau$ one has $\psi = \tfrac12$ regardless of $p$, justifying the name
half-life; here $\tau$ is a locally scoped decay parameter, not a second time
variable. If each of the $n+2$ parameters is expanded over $L$ decay kernels,
the block-diagonal matrix $\Lambda(T) \in \mathbb{R}^{(n+2)\times N_b}$, with
$N_b = (n+2)L$, collects the $\psi$-values and
\begin{equation}\label{eq:sigma-full}
\ket{\beta(T)} = \Lambda(T)\,\ket{\mathbf b},
\qquad
\sigma(\mln,\lF,T) = \braket{\tilde\phi(\mln,T)|\mathbf b},
\qquad
\ket{\tilde\phi(\mln,T)} = \Lambda(T)^\top \ket{\phi(\mln,T)}.
\end{equation}
The model remains linear in the extended parameter ket $\ket{\mathbf b}$. The
$T$-derivative of $\psi$, needed for the local-volatility numerator, is
\begin{equation}\label{eq:psi-prime}
\psi'(T;\tau,p) = -\frac{p\ln 2}{\tau}\!\left(\frac{T}{\tau}\right)^{p-1}
\psi(T;\tau,p).
\end{equation}

%% ======================================================================
\section{Calibration}
\label{sec:calibration}
%% ======================================================================

\subsection{Ordinary Least Squares}
\label{sec:ols}

Given $J$ market observations $\{(\hat m_j, \hat\sigma_j)\}_{j=1}^{J}$, form the
design matrix $\Phi \in \mathbb{R}^{J\times(n+2)}$ whose $j$-th row is
$\bra{\phi(\hat m_j)}$, and the observation ket
$\ket{\hat\sigma} = (\hat\sigma_1,\ldots,\hat\sigma_J)^\top$. Since
$\sigma(\hat m_j) = \braket{\phi(\hat m_j)|\beta}$, the model-implied vector is
$\Phi\parket$ and the OLS estimator is
\begin{equation}\label{eq:ols}
\ket{\hat\beta} = (\Phi^\top\Phi)^{-1}\Phi^\top\ket{\hat\sigma},
\end{equation}
the orthogonal projection of $\ket{\hat\sigma}$ onto $\mathrm{Im}(\Phi)$.

\begin{proposition}[Regularity]\label{prop:ols-regularity}
The estimator~\eqref{eq:ols} is well defined if and only if $\Phi$ has full
column rank $n+2$. Sufficient: $J \ge n+2$ and at least one observed moneyness
inside each kernel support.
\end{proposition}

\begin{remark}[Conditioning]\label{rem:conditioning}
When supports overlap heavily or data is sparse, $\Phi^\top\Phi$ can be
ill-conditioned; ridge regression
$\ket{\hat\beta} = (\Phi^\top\Phi + \lambda I)^{-1}\Phi^\top\ket{\hat\sigma}$
with small $\lambda > 0$ restores stability. Conversely, when the kernels have
disjoint (or nearly disjoint) supports and each support holds at least one
observation, the curvature columns of $\Phi$ have disjoint non-affine parts,
$\Phi^\top W\Phi$ is nearly block-diagonal, and the problem is well
conditioned.
\end{remark}

\subsection{Weighted and Constrained Regression}
\label{sec:wls}

Observations are weighted by inverse bid--ask spread or by vega:
\begin{equation}\label{eq:wls}
\ket{\hat\beta} = (\Phi^\top W \Phi + \lambda I)^{-1}\Phi^\top W\ket{\hat\sigma},
\qquad
W = \mathrm{diag}(\varpi_1,\ldots,\varpi_J),
\quad \varpi_j \propto \vega_j^2/\mathrm{spread}_j^2,
\end{equation}
where the locally scoped weight symbol $\varpi_j$ is distinct both from the
kernel width constants $w_i$ and from the contract-reserved $\omM$ (Margrabe,
R2), and $\vega_j$ is the Black--Scholes vega of quote $j$. To
enforce the arbitrage constraints of Section~\ref{sec:no-arb}, replace
unconstrained OLS by constrained least squares,
\[
\min_{\parket} \bigl\| \ket{\hat\sigma} - \Phi\parket \bigr\|_W^2
\quad\text{subject to}\quad g_\ell(\parket)\ge 0,\ \ell=1,\dots,L,
\]
where the $g_\ell$ are the admissibility constraints of
Section~\ref{sec:no-arb} evaluated on a moneyness grid. These are \emph{not}
uniformly linear in $\parket$, in keeping with the chapter's opening claim that
the no-arbitrage conditions are low-degree \emph{polynomial} inequalities. Only
some are affine: positivity (A1) $\sigma>0$, the calendar condition (A3)
$\mathcal N>0$ (equivalently $\sigma+2T\,\partial_T\sigma>0$, affine in $\parket$
once $\sigma$ is), and the flat-wing constraints
\eqref{eq:wing-right}--\eqref{eq:wing-left}. The butterfly condition (A2) $D>0$
is genuinely quadratic-to-quartic in $\parket$---it contains $\sigma c$, $s^2$
and $\sigma^2 T^2 s^2$ with $\sigma,s,c$ all linear in $\parket$---and even the
sufficient surrogate (S1) $c>(\sigma T/4)s^2$ is quadratic. The programme is
solved either in its true polynomially-constrained form (SOCP/SDP or
sequential quadratic programming), or by successive linearisation of $D$
about the current fit---a linear surrogate that is only \emph{sufficient} and
may reject some admissible surfaces---followed by a \emph{mandatory} final pass
verifying $D>0$ and $\mathcal N>0$ exactly on the grid.

\subsection{Convexity Bias Between Mid-Price and Mid-Vol}
\label{sec:convexity}

Let $p(\sigma)$ be the Black--Scholes put price as a function of implied
volatility, and set $p_{\mathrm{mid}} = \tfrac12(p(\sigma_{\mathrm{bid}}) +
p(\sigma_{\mathrm{ask}}))$.

\begin{proposition}[Convexity bias for OTM options]\label{prop:convexity}
Suppose $d_1 d_2 > 0$ for \emph{every} $\sigma\in[\sigma_{\mathrm{bid}},
\sigma_{\mathrm{ask}}]$ (under the standing zero rates $F=S$ this is
$\abs{\mln} > \sigma_{\mathrm{ask}}^2 T/2$, i.e.\ the whole quoted band is
sufficiently out-of-the-money). Then the Volga
$\partial^2 p/\partial\sigma^2 = \vega\,d_1 d_2/\sigma
= F\sqrt{T}\,\npdf(d_1)\,d_1 d_2/\sigma > 0$ throughout the band,
$p(\sigma)$ is convex there, and the implied vol $\sigma_{\mathrm{mid}}$ of
$p_{\mathrm{mid}}$ satisfies
\[
\sigma_{\mathrm{mid}} \ge \tfrac12(\sigma_{\mathrm{bid}} + \sigma_{\mathrm{ask}}).
\]
\end{proposition}
\begin{proof}
Write $\bar\sigma = \tfrac12(\sigma_{\mathrm{bid}}+\sigma_{\mathrm{ask}})$.
Convexity of $p$ on the band gives, by Jensen,
$p_{\mathrm{mid}} = \tfrac12\bigl(p(\sigma_{\mathrm{bid}})+p(\sigma_{\mathrm{ask}})\bigr)
\ge p(\bar\sigma)$. Since $p$ is strictly increasing in $\sigma$ (vega $>0$), its
inverse is increasing, whence
$\sigma_{\mathrm{mid}} = p^{-1}(p_{\mathrm{mid}}) \ge p^{-1}\!\bigl(p(\bar\sigma)\bigr)
= \bar\sigma$. The sign hypothesis is required across the \emph{whole} band
because $d_1 d_2$ changes sign at $\abs{\mln}=\sigma^2 T/2$; a wide quoted
spread straddling that point breaks the convexity and the conclusion.
\end{proof}

\begin{remark}\label{rem:convexity-atm}
Near the money $d_1 d_2 < 0$, the Volga is negative, $p(\sigma)$ is concave, and
the inequality reverses. The effect is small for tight spreads but material for
illiquid deep-OTM options; it biases naive mid-vol fits and is one reason to
weight by vega rather than to average vols.
\end{remark}

\begin{example}[Convexity bias]\label{ex:convexity}
Under the standing zero-rate assumption ($F = S$), take $S = 100$, $K = 85$
(an OTM put), $T = 0.25$, $\sigma_{\mathrm{bid}} = 0.25$,
$\sigma_{\mathrm{ask}} = 0.35$. Then $p(0.25) \approx 0.52$,
$p(0.35) \approx 1.54$, so $p_{\mathrm{mid}} \approx 1.03$, whose implied vol is
$\sigma_{\mathrm{mid}} \approx 0.307$---above the arithmetic mid $0.300$.
\end{example}

%% ======================================================================
\section{Sensitivities and Total Greeks}
\label{sec:sensitivity}
%% ======================================================================

\subsection{Vega and Volga Decomposed per Kernel}
\label{sec:vega-decomp}

For the operative time-scaled model \eqref{eq:phi-time},
$\sigma(\mln,\lF,T) = \braket{\phi(\mln,T)|\beta}$, the first-order
sensitivity of an option price $C$ to each parameter factorises through the
Black--Scholes vega $\vega = \partial C/\partial\sigma$:
\begin{equation}\label{eq:vega-decomp}
\frac{\partial C}{\partial\beta_i} = \vega\cdot\phi_i(\mln,T),
\end{equation}
where $\phi_i(\mln,T)$ is the $i$-th \emph{time-scaled} feature component,
carrying for the moneyness rows the $1/\sqrt{\tilde T}$ factor of
\eqref{eq:phi-time} (using the un-scaled $\phi_i(\mln)$ mis-scales per-kernel
risk attribution by $\sqrt{\tilde T}$). The second-order sensitivity factorises
through the Volga:
\begin{equation}\label{eq:volga-decomp}
\frac{\partial^2 C}{\partial\beta_i\,\partial\beta_j}
= \Volga\cdot\phi_i(\mln,T)\,\phi_j(\mln,T),
\qquad
\Volga = \frac{\partial^2 C}{\partial\sigma^2}
= \vega\,\frac{d_1 d_2}{\sigma}.
\end{equation}

\begin{remark}[Curvature orthogonality is not feature orthogonality]
\label{rem:orthogonality}
For non-overlapping kernels the \emph{curvature} functions satisfy
$\kappa_i(\mln)\kappa_j(\mln) = 0$ for $i\neq j$, so curvature contributions are
independent and risk attribution at the level of $c(\mln)$ is diagonal. The
\emph{feature} functions do not inherit this: beyond its support $I_i(\mln)$ is
affine (not zero), so $I_i(\mln) I_j(\mln) \neq 0$ in general and the Hessian
\eqref{eq:volga-decomp} is not diagonal in the $\beta$ indices for a single
instrument. Diagonality holds in curvature space, not feature space.
\end{remark}

\subsection{Routing Through the Forward: Total Greeks}
\label{sec:total-greeks}

The implied-vol model supplies the derivatives that convert flat Black--Scholes
Greeks into smile-adjusted (total) Greeks. Here the decomposition
\eqref{eq:decomp} matters decisively: as spot moves, \emph{both} the ATM level
(through $\lF$) and the smile shape (through $\mln$) respond. Carrying the chain
rule through $F(S)$,
\[
\lF = \ln(F/\Fref), \quad \pd{\lF}{S} = \frac{1}{F}\td{F}{S};
\qquad
\mln = \ln(K/F), \quad \pd{\mln}{S} = -\frac{1}{F}\td{F}{S}.
\]
Under the standing zero-rate assumption $F = S$, $\dd F/\dd S = 1$,
so $\partial\lF/\partial S = 1/S$ and $\partial\mln/\partial S = -1/S$.
Hence the total first and second spot-derivatives of implied vol are
\begin{align}
\td{\sigma}{S}
&= \pd{\satm}{\lF}\pd{\lF}{S} + \pd{\ssm}{\mln}\pd{\mln}{S}
= \frac{\volslope - s(\mln)}{S}, \label{eq:dsig-dS}\\[4pt]
\frac{\dd^2\sigma}{\dd S^2}
&= \frac{\bigl(\partial_{\lF}^2\satm - \volslope\bigr) + \bigl(c(\mln) + s(\mln)\bigr)}{S^2}.
\label{eq:d2sig-dS2}
\end{align}
The routed spot-derivatives on the left are \emph{total} ($\dd/\dd S$), since
they run through both arguments $\mln(S)$ and $\lF(S)$; the components
$\partial\satm/\partial\lF$ and $\partial\ssm/\partial\mln$ on the right remain
genuine partials. These feed the total (smile-adjusted) gamma,
\begin{equation}\label{eq:Gadj}
\Gadj = \Gamma + 2\,\Vanna\cdot\td{\sigma}{S}
+ \Volga\cdot\!\left(\td{\sigma}{S}\right)^{\!2}
+ \vega\cdot\frac{\dd^2\sigma}{\dd S^2}.
\end{equation}
Equation \eqref{eq:Gadj} is the second total derivative $\dd^2/\dd S^2$ of
$C(S,\sigma(S))$: expanding
$\dd C/\dd S = \partial_S C + \partial_\sigma C\cdot\dd\sigma/\dd S$ once more,
and using $C\in C^2$ so that mixed partials commute
($\partial_S\partial_\sigma C = \partial_\sigma\partial_S C = \Vanna$), gives the
four terms, with $\Gamma=\partial_S^2 C$, $\Vanna=\partial_S\partial_\sigma C$,
$\Volga=\partial_\sigma^2 C$, $\vega=\partial_\sigma C$ the frozen-vol
Black--Scholes Greeks.

\begin{remark}[The $\satm=\text{const}$ collapse, flagged]\label{rem:const-flag}
The C4 convention sets $\satm = \mathrm{const}$, i.e.\
$\volslope = 0$ and $\partial_{\lF}^2\satm = 0$. Under that convention
\eqref{eq:dsig-dS}--\eqref{eq:d2sig-dS2} reduce to
$\dd\sigma/\dd S = -s(\mln)/S$ and
$\dd^2\sigma/\dd S^2 = (c(\mln)+s(\mln))/S^2$---the pure-smile terms
only, discarding exactly the $\volslope$ contribution of
Remark~\ref{rem:two-slopes}. The full
forms \eqref{eq:dsig-dS}--\eqref{eq:d2sig-dS2} are the ones that must enter any
book carrying genuine ATM dynamics, and it is through $\volslope$ that the
surface term structure enters the multi-name shadow Greeks of the second
volume.
\end{remark}

%% ======================================================================
\section{Local Volatility}
\label{sec:localvol}
%% ======================================================================

\subsection{Dupire's Formula and the Black--Scholes Partials}
\label{sec:dupire}

Under zero rates and dividends, Dupire's formula (Dupire 1994) reads
\begin{equation}\label{eq:dupire}
\sloc^2(K,T) = \frac{\partial_T C}{\tfrac12 K^2\,\partial_K^2 C}.
\end{equation}
The derivatives are those of the Black--Scholes call
$C = F\,\ncdf(d_1) - K\,\ncdf(d_2)$, in which
$d_1 = -\mln/(\sigma\sqrt{T}) + \sigma\sqrt{T}/2$ and $d_2 = d_1 - \sigma\sqrt{T}$;
$\sigma = \satm(\lF,T) + \ssm(\mln,T)$ depends on $K$ only through
$\mln = \ln(K/F)$ (the ATM component is constant in $K$ at fixed forward).

\paragraph{Temporal derivative.} Holding $K$ fixed---and, under the standing
zero-rate assumption $F = S$, the forward is constant across maturities so
$\mln$ is fixed as well---
\begin{equation}\label{eq:dT-C}
\partial_T C = \frac{F\,\npdf(d_1)}{\sqrt{T}}
\left(\frac{\sigma}{2} + T\,\pd{\sigma}{T}\right).
\end{equation}

\paragraph{Spatial second derivative.} With $\sigma = \sigma(\mln(K),T)$ the
chain rule gives
\begin{equation}\label{eq:d2K-C}
\frac{\dd^2 C}{\dd K^2}
= \atfixed{\pdd{C}{K}}{\sigma}
+ 2\,\pdc{C}{K}{\sigma}\pd{\sigma}{K}
+ \pdd{C}{\sigma}\left(\pd{\sigma}{K}\right)^{\!2}
+ \pd{C}{\sigma}\,\pdd{\sigma}{K},
\end{equation}
where the frozen-vol Black--Scholes partials are
\begin{equation}\label{eq:bs-partials}
\atfixed{\pdd{C}{K}}{\sigma} = \frac{\npdf(d_2)}{K\sigma\sqrt{T}},
\qquad
\pdc{C}{K}{\sigma} = \frac{\npdf(d_2)\,d_1}{\sigma},
\qquad
\pdd{C}{\sigma} = F\sqrt{T}\,\npdf(d_1)\,\frac{d_1 d_2}{\sigma}.
\end{equation}
These are verified in Section~\ref{sec:bs-verify}; the key identity is
$F\,\npdf(d_1) = K\,\npdf(d_2)$, exact in the forward measure since
$d_1^2 - d_2^2 = -2\mln$ and $K/F = e^{\mln}$.

\subsection{The Local-Volatility Master Formula}
\label{sec:master-formula}

Converting $K$-derivatives to $\mln$-derivatives via
$\partial\sigma/\partial K = (1/K)\,\partial\sigma/\partial\mln$ and
$K^2\,\partial_K^2\sigma = \partial_\mln^2\sigma - \partial_\mln\sigma$, writing
$s = \partial_\mln\sigma = s(\mln)$ and $c = \partial_\mln^2\sigma = c(\mln)$ for
the smile skew and curvature, and factoring out
$F\,\npdf(d_1)/(2\sigma\sqrt{T})$, the denominator becomes
\begin{equation}\label{eq:denom-K}
\tfrac12 K^2\,\partial_K^2 C = \frac{F\,\npdf(d_1)}{2\sigma\sqrt{T}}\,D(\mln,T),
\end{equation}
with
\begin{equation}\label{eq:D-def}
D(\mln,T) = 1 - \frac{2\mln}{\sigma}\,\pd{\sigma}{\mln}
+ \left(\frac{\mln^2}{\sigma^2} - \frac{\sigma^2 T^2}{4}\right)\!
\left(\pd{\sigma}{\mln}\right)^{\!2}
+ \sigma T\,\pdd{\sigma}{\mln}.
\end{equation}
The clean form \eqref{eq:D-def} rests on one cancellation:
the substitution $K^2\partial_K^2\sigma = \partial_\mln^2\sigma -
\partial_\mln\sigma$ injects a term $-\sigma T\,s$, while the expansion
$2 d_1\sqrt{T} = -2\mln/\sigma + \sigma T$ injects an equal and opposite
$+\sigma T\,s$. Combining \eqref{eq:dT-C} and
\eqref{eq:denom-K},
\begin{equation}\label{eq:local-vol}
\boxed{\;\sloc^2(\mln,T) = \sigma^2\,\frac{\mathcal N(\mln,T)}{D(\mln,T)},\;}
\qquad
\mathcal N(\mln,T) = 1 + \frac{2T}{\sigma}\,\pd{\sigma}{T},
\end{equation}
and $D$ as in \eqref{eq:D-def}. The numerator $\mathcal N$ carries the
calendar-spread condition; the denominator $D$ carries the butterfly condition
(Section~\ref{sec:no-arb}).

\subsection{Application to the Kernel Model}
\label{sec:kernel-localvol}

From the time-scaled model \eqref{eq:phi-time},
\begin{align}
\pd{\sigma}{\mln} &= s(\mln)
= \frac{1}{\sqrt{\tilde T}}\Bigl(s_0 + \sum_i c_i J_i(\mln)\Bigr),
\label{eq:dsig-dm}\\
\pdd{\sigma}{\mln} &= c(\mln)
= \frac{1}{\sqrt{\tilde T}}\sum_i \frac{c_i}{h_i}\,
\kappa\!\left(\frac{\mln-\mu_i}{h_i}\right).
\label{eq:d2sig-dm2}
\end{align}
The temporal derivative $\partial\sigma/\partial T$ needed for $\mathcal N$
follows from \eqref{eq:sigma-full} by the chain rule, with two sources: the
time scaling of the moneyness features and the decay functions. For the first,
the kernel scaling \eqref{eq:kernel-scaling} makes the centres $\mu_i$ and
widths $h_i$ themselves scale with $\sqrt{\tilde T}$, so each moneyness feature
$I_i(\mln)/\sqrt{\tilde T}$ is a \emph{pure function} $f_i(u)$ of the
standardised moneyness $u=\mln/\sqrt{\tilde T}$, with $f_i'(u)=J_i(\mln)$ (and
the linear feature is $f_0(u)=u$). Differentiating $f_i(\mln/\sqrt{\tilde T})$
in $T$ therefore routes entirely through $u$:
\[
\td{\tilde T}{T} = \frac{252}{252+z}\cdot\frac{1}{(1+T'/T_{\max})^2},
\qquad
\pd{}{T}\ket{\phi(\mln,T)}
= -\frac{1}{2\tilde T}\td{\tilde T}{T}\!
\begin{pmatrix} 0 \\ \mln/\sqrt{\tilde T} \\ \mln\,J_1(\mln)/\sqrt{\tilde T} \\ \vdots
\\ \mln\,J_n(\mln)/\sqrt{\tilde T} \end{pmatrix},
\]
each curvature row being
$-\tfrac{1}{2\tilde T}(\dd\tilde T/\dd T)\,u\,f_i'(u)
= -\tfrac{1}{2\tilde T}(\dd\tilde T/\dd T)\,\mln J_i(\mln)/\sqrt{\tilde T}$---the
$\mln J_i(\mln)$, \emph{not} $I_i(\mln)$, being the correct chain-rule factor
once $\mu_i,h_i$ scale with $\sqrt{\tilde T}$; the two coincide only for the
linear row, where $u f_0'(u)=f_0(u)$ preserves the form $\mln/\sqrt{\tilde T}$.
The second source, the decay functions, enters through \eqref{eq:psi-prime}. In the single-expiry
case (fixed $T$, no decay) $\partial\sigma/\partial T = 0$ and $\mathcal N = 1$.
Inside the support of kernel $i$ the curvature peaks at the centre,
\begin{equation}\label{eq:max-curvature}
c(\mu_i) = \frac{3\,c_i}{4\,w_i\,\sref\,\tilde T},
\end{equation}
the single most useful quantity for the admissibility checks below.

\subsection{Worked Example: The Quadratic Smile}
\label{sec:quadratic-example}

Example~\ref{ex:const-curv} is now worked completely: a smile quadratic in
log-moneyness, i.e.\ constant curvature $c_0$, produced by one dominant central
kernel. Write
\[
\sigma(\mln,\lF,T) = \satm(\lF,T) + s_0\,\mln + \tfrac12 c_0\,\mln^2,
\qquad
\pd{\sigma}{\mln} = s_0 + c_0\,\mln,
\qquad
\pdd{\sigma}{\mln} = c_0.
\]
Substituting into the master formula \eqref{eq:local-vol} gives the local
variance at every $(\mln,T)$; two regimes follow.

\medskip
\noindent\textbf{At the money ($\mln = 0$).} The terms in $\mln$ drop and
$\sigma = \satm(\lF,T)$, so, using the completed-square form
\eqref{eq:D-completed} below,
\begin{equation}\label{eq:quad-atm}
\sloc^2(0,T)
= \satm^2\;
\frac{\,1 + \dfrac{2T}{\satm}\,\partial_T\satm\,}
{\,1 + \satm T\,c_0 - \dfrac{\satm^2 T^2}{4}\,s_0^2\,}.
\end{equation}
Positivity of the denominator is a lower bound on the ATM curvature,
\begin{equation}\label{eq:quad-constraint}
c_0 \;>\; \frac{\satm T}{4}\,s_0^2 \;-\; \frac{1}{\satm T},
\end{equation}
which tightens with maturity: a long-dated smile with appreciable ATM skew
$s_0$ \emph{must} carry positive curvature to remain free of butterfly
arbitrage. This is the exact quadratic-smile instance of the general
sufficient condition of Proposition~\ref{prop:D-sufficient}.

\medskip
\noindent\textbf{The purely linear (zero-curvature) sub-case.} Setting
$c_0 = 0$ leaves $\sigma$ affine in $\mln$. The denominator at the money is then
$1 - \tfrac14\satm^2 T^2 s_0^2$, which turns negative once
$\satm T\,\abs{s_0} > 2$: a straight-line smile with non-zero skew inevitably
violates the butterfly bound at sufficiently long maturity. Curvature is not a
cosmetic refinement---it is what rescues admissibility. This is precisely why
the wing tails of Section~\ref{sec:no-arb} are flattened to constants rather
than left sloped, and why every kernel grid must retain genuine curvature in
the fitted region.

%% ======================================================================
\section{The No-Arbitrage Admissibility Ladder}
\label{sec:no-arb}
%% ======================================================================

A positive, well-defined local variance requires both $\mathcal N > 0$ and
$D > 0$ everywhere; together with positivity of the vol itself these are the
first rungs of the admissibility ladder that continues, in the multi-name
volume, into correlation and basket constraints.

\begin{definition}[Admissible parameter set]\label{def:admissible}
A parameter ket $\ket{\mathbf b}$ is \emph{admissible} if, for all $(\mln,T)$ in
the model domain:
\begin{enumerate}[label=\textup{(A\arabic*)}]
\item \textbf{Positivity:} $\sigma(\mln,\lF,T) > 0$.
\item \textbf{Butterfly:} $D(\mln,T) > 0$ (strictly positive risk-neutral
density).
\item \textbf{Calendar spread:} $\mathcal N(\mln,T) > 0$, i.e.\ total variance
$w(\mln,T) = \sigma^2 T$ is strictly increasing in $T$.
\end{enumerate}
Write $\mathcal A$ for the admissible set. $\mathcal A$ is taken \emph{open}
(strict inequalities): the strict interior is what keeps the division in the
local-variance formula \eqref{eq:local-vol} well posed and gives the constrained
optimiser of Section~\ref{sec:wls} a safe interior to sit in. The closed
relaxation ($D\ge0 \iff$ non-negative density, $\mathcal N\ge0 \iff w$
non-decreasing) is arbitrage-free but sits on the boundary, where $\sloc$ may
degenerate---the open-versus-closed distinction is operationally load-bearing
because the optimiser will activate exactly the flat-calendar and zero-density
faces.
\end{definition}

\begin{proposition}[Necessity for arbitrage freedom]\label{prop:no-arb}
If $\ket{\mathbf b}\notin\mathcal A$, the surface admits static arbitrage.
Precisely: if (A2) fails the risk-neutral density is negative (a butterfly
arbitrage); if (A3) fails total variance decreases in $T$ at fixed strike, a
calendar-spread arbitrage; if (A1) fails the Black--Scholes map is ill-posed
($\sloc^2$ undefined), and a negative total variance $w=\sigma^2 T$ prices
options below intrinsic value, breaking monotonicity of $C$ in $\sigma$.
\end{proposition}
\begin{proof}
(A2): by \eqref{eq:denom-K}, $\tfrac12 K^2\partial_K^2 C
= F\npdf(d_1)\,D/(2\sigma\sqrt T)$, so the Breeden--Litzenberger density
$q(K)=\partial_K^2 C = F\npdf(d_1)\,D/(K^2\sigma\sqrt T)$ is a \emph{positive}
multiple of $D$; hence $q>0 \iff D>0$. (A3): $\partial_T w
= \partial_T(\sigma^2 T) = \sigma^2 + 2\sigma T\,\partial_T\sigma
= \sigma^2\bigl(1+(2T/\sigma)\partial_T\sigma\bigr) = \sigma^2\,\mathcal N$, so
$\mathcal N>0 \iff w$ strictly increasing in $T$; a strike at which $w$ decreases
furnishes a calendar-spread arbitrage (under the standing zero rates $F=S$,
fixed $K$ is fixed $\mln$---flagged). (A1): $\sigma<0$ does not give literally
negative prices (a deep-ITM call stays above zero) but a price below intrinsic
value, violating monotonicity of $C$ in $\sigma$, and $\sloc^2$ in
\eqref{eq:local-vol} is undefined. The calendar-spread argument for (A3)
follows Lee (2004) and Dupire (1994).
\end{proof}

\subsection{Wing Behaviour and Lee's Moment Formula}
\label{sec:wings}

Beyond all kernel supports $\sigma$ is affine in $\mln$, so $w = \sigma^2 T$
grows quadratically in $\abs{\mln}$ unless the wing slope vanishes. Lee's moment
formula (Lee 2004) bounds the asymptotic slope of total variance,
\[
\limsup_{\mln\to+\infty}\frac{w(\mln,T)}{\mln}\le 2,
\qquad
\limsup_{\mln\to-\infty}\frac{w(\mln,T)}{\abs{\mln}}\le 2,
\]
which an unconstrained affine wing violates. The remedy is to force each wing
flat. With $\alpha_i = P_\kappa(-\mu_i/h_i)$ as in \eqref{eq:alpha-def}, the
right- and left-wing slopes are proportional to $s_0 + \sum_i c_i(1-\alpha_i)$
and $s_0 + \sum_i c_i J_i^{-}$ respectively, where
\[
J_i^{-} = \lim_{\mln\to-\infty}J_i(\mln) = -\alpha_i =
\begin{cases}
0, & \mu_i - h_i > 0,\\
-1, & \mu_i + h_i < 0,\\
-P_\kappa(-\mu_i/h_i), & \abs{\mu_i} < h_i,
\end{cases}
\]
giving the two \emph{wing constraints}
\begin{align}
s_0 + \sum_{i=1}^n c_i(1-\alpha_i) &= 0 \quad\text{(right wing)},
\label{eq:wing-right}\\
s_0 + \sum_{i=1}^n c_i\,J_i^{-} &= 0 \quad\text{(left wing)}.
\label{eq:wing-left}
\end{align}
When all supports lie right of the origin ($\alpha_i = 0$), these reduce to
$s_0 + \sum_i c_i = 0$ and $s_0 = 0$.

\begin{proposition}[Compliance with Lee's bound]\label{prop:lee}
Under \eqref{eq:wing-right}--\eqref{eq:wing-left}, $\partial_\mln\sigma = 0$
beyond all supports, so $\sigma\to\sigma_\infty^{\pm}$ (finite positive
constants), $w = (\sigma_\infty^{\pm})^2 T$ is bounded, and
$w(\mln,T)/\abs{\mln}\to 0 < 2$ as $\abs{\mln}\to\infty$. Lee's moment formula is
satisfied with room to spare.
\end{proposition}

\begin{remark}[Flat wings under-parameterise the tails]\label{rem:flat-wing-cost}
The room to spare is not free. Flat wings are strictly stronger than Lee's
bound, which permits a total-variance wing slope up to $2$. Pinning
$\partial_\mln\sigma=0$ beyond the supports fixes deep-OTM implied vol at the
constants $\sigma_\infty^{\pm}$, so the model is structurally blind to
legitimate wing \emph{steepening}: it cannot represent linear wing growth up to
Lee's boundary and therefore systematically under-prices the tails and
tail-hedge cost relative to a surface sitting near that boundary. This is a
deliberate ``precise in the centre, silent in the tails'' trade; a book with
material tail exposure requires a bounded non-zero wing slope within
Lee's bound rather than a forced zero.
\end{remark}

\subsection{Sufficient Conditions for the Butterfly Bound}
\label{sec:butterfly}

\begin{proposition}[Completing the square in $D$]\label{prop:D-sufficient}
Writing $s = \partial_\mln\sigma$, $c = \partial_\mln^2\sigma$, the denominator
\eqref{eq:D-def} admits the identity
\begin{equation}\label{eq:D-completed}
D(\mln,T) = \left(1 - \frac{\mln\,s}{\sigma}\right)^{\!2}
+ \sigma T\,c - \frac{\sigma^2 T^2}{4}\,s^2.
\end{equation}
Consequently $D(\mln,T) > 0$ if either holds:
\begin{enumerate}[label=\textup{(S\arabic*)}]
\item \textbf{Curvature dominates skew (pointwise, $\mln$-free):}
$\;\sigma T\,c > \dfrac{\sigma^2 T^2}{4}\,s^2$, i.e.\
$c > \dfrac{\sigma T}{4}\,s^2$.
\item \textbf{Joint bound (exact):}
$\;\sigma T\,c > -\Bigl(1 - \dfrac{\mln s}{\sigma}\Bigr)^{\!2}
+ \dfrac{\sigma^2 T^2}{4}\,s^2.$
\end{enumerate}
\end{proposition}
\begin{proof}
Expanding $(1 - \mln s/\sigma)^2 = 1 - 2\mln s/\sigma + \mln^2 s^2/\sigma^2$ and
comparing with \eqref{eq:D-def} establishes \eqref{eq:D-completed}. Since the
squared term is non-negative, $D \ge \sigma T c - \tfrac14\sigma^2 T^2 s^2$,
strict under (S1). Condition (S2) is the exact rearrangement $D>0$, and (S1)
implies (S2) for all $\mln$ because dropping the non-negative squared term only
weakens the requirement.
\end{proof}

\begin{remark}[Kernel-parameter estimate]\label{rem:kernel-estimate}
At a kernel centre $\mu_i$ the curvature attains its peak
\eqref{eq:max-curvature}, $c(\mu_i)=3c_i/(4w_i\sref\tilde T)$, so, carrying
$\tilde T$ explicitly,
\[
\sigma T\,c(\mu_i) = \frac{3\,\sigma\,c_i\,T}{4\,w_i\,\sref\,\tilde T}
\approx \frac{3\,c_i\,\sigma}{4\,w_i\,\sref}
\qquad (T\approx\tilde T),
\]
a genuinely dimensionless quantity, but not because $c_i$ and $\sigma$ share
units. In the time-scaled feature model $c_i$ is itself \emph{dimensionless}:
the $1/\sqrt{\tilde T}$ carried in the moneyness rows of the feature map absorbs
the volatility dimension, so the fitted coefficients $c_i$ are pure numbers. The
product is then dimensionless term by term---$c_i$ dimensionless, $\sigma/\sref$
a dimensionless volatility ratio, and $T/\tilde T\approx1$---rather than through
any cancellation of a shared vol unit against $\sref$. (The earlier
$\approx\frac{3c_i\sigma}{4w_i\sref}$ is the same quantity at $T\approx\tilde T$;
it is dimensionless for the identical reason, the $\sref$ in the denominator
scaling $\sigma$, not $c_i$, to a ratio.) Two facts follow, and \emph{neither
is a cap} on $\abs{c_i}$. \emph{(i)} Where the curvature is negative
($c_i<0$), the term $\sigma T c$ pulls $D$ down, and $D>0$ near the centre
requires roughly $3\abs{c_i}\sigma/(4w_i\sref) < 1+(1-\mln s/\sigma)^2$ by the
completed-square form \eqref{eq:D-completed}. \emph{(ii)} Condition (S1) is a
\emph{lower} bound, $c(\mln)>(\sigma T/4)\,s(\mln)^2$: butterfly no-arbitrage
demands \emph{enough} positive curvature, not a bounded amount---the threat is
curvature too negative, not too large. Because $c$ is largest at the centre and
vanishes between and beyond the kernels---where $s^2$ is nonetheless
accumulated---(S1) is easiest to satisfy at the very point evaluated here; the
binding check therefore lies between and beyond the supports and at the
$\mu_i\pm h_i$ kinks, where the exact denominator
\eqref{eq:D-def} or the joint bound (S2) must be used.
\end{remark}

\subsection{Worked Admissibility and Local-Volatility Checks}
\label{sec:admissibility-examples}

\begin{example}[Positivity at a kernel centre]\label{ex:positivity}
These worked checks use the single-expiry, \emph{un-time-scaled} model
\eqref{eq:sigma-linear} (equivalently, $\tilde T$ folded into the bandwidth
$h_1$), so the centre curvature is read directly as
$c(\mu_1)=\kappa(0)\,c_1/h_1 = 3c_1/(4h_1)$, \emph{not} from the time-scaled
\eqref{eq:max-curvature}. A single put-side kernel: $c_1 = 2.0$,
$\mu_1 = -0.15$, $h_1 = 0.08$, $\sigma_0 = 0.20$, $s_0 = -0.10$, $T = 0.50$. The
support $[-0.23,-0.07]$ is entirely left of the origin, so $\alpha_1 = 1$. At
$\mln = \mu_1 = -0.15$:
\begin{itemize}
\item $\sigma \approx 0.245$;
\item $s = s_0 + c_1 J_1(\mu_1) = -0.10 + 2.0(-0.50) = -1.10$
(at the centre $J_1 = P_\kappa(0) - \alpha_1 = 0.5 - 1 = -0.5$);
\item $c = 3(2.0)/(4\cdot 0.08) = 18.75$;
\item $D = 1 - 2\dfrac{-0.15}{0.245}(-1.10)
+ \Bigl(\dfrac{0.15^2}{0.245^2} - \dfrac{0.245^2\cdot 0.5^2}{4}\Bigr)(1.21)
+ 0.245\cdot 0.5\cdot 18.75
= 1 - 1.35 + 0.45 + 2.30 = 2.40 > 0.$
\end{itemize}
Flipping $c_1$ to $-3.0$ (a concave smile) and reading a butterfly
violation off the curvature term alone is an invalid \emph{ceteris
paribus}: $\sigma$, $s$ \emph{and} $c$ all depend on $c_1$. Recomputing
honestly at the same point $\mln=\mu_1=-0.15$:
$\sigma = 0.20 + 0.015 + (-3.0)(0.015) = 0.17$,
$s = -0.10 + (-3.0)(-0.50) = +1.40$ (the skew flips sign),
$c = 3(-3.0)/(4\cdot0.08) = -28.125$, giving
\[
D = 1 - 2\dfrac{-0.15}{0.17}(1.40)
+ \Bigl(\dfrac{0.15^2}{0.17^2} - \dfrac{0.17^2\cdot 0.5^2}{4}\Bigr)(1.96)
+ 0.17\cdot 0.5\cdot(-28.125)
= 1 + 2.47 + 1.52 - 2.39 = +2.60 > 0.
\]
The centre is \emph{rescued}: the sign flip of $s$ makes the $-2\mln s/\sigma$
term strongly positive, more than offsetting the negative $\sigma T c$. The
lesson is precisely that negative curvature lowers $D$ through $\sigma T c$ but
simultaneously moves $s$, so a butterfly violation cannot be diagnosed from one
term at one point: it must be \emph{located} by evaluating $D$ over the whole
support and its edges. Indeed $c(\mln)$ is continuous but kinks at
$\mu_i\pm h_i$ (extra extremum candidates), and with $c_1=-3.0$ the vol itself
turns negative in the far left tail---an (A1) violation---so both positivity and
butterfly must be scanned across moneyness, never sampled at the centre alone.
\end{example}

\begin{example}[Local volatility at the kernel centre]\label{ex:localvol}
Continuing Example~\ref{ex:positivity} ($\sigma = 0.245$, $s = -1.10$,
$c = 18.75$, $T = 0.50$), suppose a mild term-structure decay
$\partial\sigma/\partial T \approx -0.02$. Then
\[
\mathcal N = 1 + \frac{2(0.50)}{0.245}(-0.02) = 0.918,
\qquad
D = 2.40,
\qquad
\sloc^2 = \sigma^2\,\frac{\mathcal N}{D}
= 0.245^2\cdot\frac{0.918}{2.40} = 0.0230,
\]
so $\sloc \approx 0.152$, about $15.2\%$---well below the implied $24.5\%$. The
strong positive curvature inflates $D$; intuitively, a hedging portfolio
captures the pronounced gamma of the convex smile, lowering the local
volatility the model needs to reproduce the same prices.
\end{example}

\begin{remark}[The ladder continues]\label{rem:ladder-forward}
Definition~\ref{def:admissible} and Propositions~\ref{prop:no-arb}--\ref{prop:D-sufficient}
close the single-name rungs: positivity, butterfly, calendar spread, and Lee
wing control. In the multi-name volume the same ladder acquires correlation
positivity and basket-convexity rungs, with the constituent vols entering as
the ATM components $\satmi{i}(\lFi{i},T)$ of this chapter's decomposition, each
carrying its own forward log-level $\lFi{i}$ (contract bridging identity C1).
\end{remark}

%% ======================================================================
\section{Summary}
\label{sec:summary}
%% ======================================================================

The curvature-kernel construction gives the implied surface as a single inner
product, $\sigma(\mln,\lF,T) = \braket{\tilde\phi(\mln,T)|\mathbf b}$, split by
the contract decomposition into an ATM level $\satm(\lF,T)$ and a smile shape
$\ssm(\mln,T)$ with $\ssm(0,T)=0$. Four properties make the representation
worth its parameters. \emph{Linearity}: calibration is (weighted) least
squares, closed-form and non-iterative. \emph{Locality}: compactly supported
Epanechnikov kernels give regional control of curvature, and non-overlapping
supports yield independent contributions. \emph{Dupire in closed form}: the
local variance is a rational function of the same coefficients, its denominator
encoding the butterfly bound and its numerator the calendar-spread bound.
\emph{Routed sensitivities}: vega decomposes per kernel, and---carrying the
chain rule through $F(S)$ and keeping $\volslope$ alive rather than collapsing
$\satm$ to a constant---the model feeds the total-gamma and shadow-Greek spine.
The one structural liability, the affine wing, is disciplined by the flat-wing
constraints \eqref{eq:wing-right}--\eqref{eq:wing-left}; the unconstrained OLS
is replaced by the vega-weighted constrained regression of
Section~\ref{sec:wls}, which enforces the invariants of
Definition~\ref{def:admissible}.

%% ======================================================================
\section{Appendix: Closed-Form Basis Functions}
\label{sec:closed}
%% ======================================================================

Standardise via $\mln^* = (\mln - \mu_i)/h_i$ and define the kernel's cumulative
(single-integral) function $P_\kappa$, normalised to transition from $0$ to $1$:
\begin{equation}\label{eq:Pkappa}
P_\kappa(\mln^*) = \begin{cases}
0, & \mln^* < -1,\\[2pt]
\dfrac34\!\left(\mln^* - \dfrac{(\mln^*)^3}{3} + \dfrac23\right),
& -1\le\mln^*\le 1,\\[8pt]
1, & \mln^* > 1,
\end{cases}
\qquad P_\kappa(\mp 1) = \{0,1\}.
\end{equation}
Since \eqref{eq:Ji}--\eqref{eq:Ii} integrate from $\mln = 0$, subtract the
origin value $P_\kappa(\mln^*_0)$ with $\mln^*_0 = -\mu_i/h_i$, and write
$\alpha = P_\kappa(\mln^*_0)$:
\begin{equation}\label{eq:Ji-closed}
J_i(\mln) = \begin{cases}
-\alpha, & \mln < \mu_i - h_i,\\[2pt]
P_\kappa(\mln^*) - \alpha, & \mu_i - h_i \le \mln \le \mu_i + h_i,\\[2pt]
1 - \alpha, & \mln > \mu_i + h_i.
\end{cases}
\end{equation}
Thus $J_i$ transitions $0\to 1$ when the kernel is right of the origin
($\alpha=0$), $-1\to 0$ when left ($\alpha=1$), and $-\tfrac12\to+\tfrac12$ when
centred ($\alpha=\tfrac12$). The double-integral function is
\[
Q_\kappa(\mln^*) = \begin{cases}
0, & \mln^* < -1,\\[2pt]
h_i\!\left(\dfrac{3}{16} + \dfrac{\mln^*}{2} + \dfrac{3(\mln^*)^2}{8}
- \dfrac{(\mln^*)^4}{16}\right), & -1\le\mln^*\le 1,\\[8pt]
h_i\,\mln^*, & \mln^* > 1,
\end{cases}
\]
and the origin-adjusted double integral is
\begin{equation}\label{eq:Ii-closed}
I_i(\mln) = Q_\kappa(\mln^*) - Q_\kappa(\mln^*_0) - \alpha\,\mln,
\end{equation}
the final $-\alpha\,\mln$ absorbing the linear offset in $J_i$. Boundary checks
confirm $Q_\kappa(-1) = 0$ and $Q_\kappa(+1) = h_i$; for a kernel left of the
origin ($\alpha = 1$, $Q_\kappa(\mln^*_0) = -\mu_i$) one finds $I_i(\mln) = 0$
for $\mln > \mu_i + h_i$ (all mass traversed, the contribution absorbed into
slope) and $I_i(\mln) = \mu_i - \mln$ as $\mln\to-\infty$ (affine, consistent
with the linear wing).

\begin{proposition}[Derivatives and regularity]\label{prop:derivatives}
On the interior of each region,
$\dd I_i/\dd\mln = J_i(\mln)$ and
$\dd J_i/\dd\mln = h_i^{-1}\kappa((\mln-\mu_i)/h_i)$, with the one-sided
derivatives agreeing at $\mu_i\pm h_i$. Hence $I_i \in C^2$ and $J_i \in C^1$
globally, and the reconstructed smile $\ssm(\mln,T)$ is globally $C^2$; the
extra degree of smoothness is exactly the payoff of $\kappa(\pm 1) = 0$.
\end{proposition}
\begin{proof}
In the support region, using $\dd\mln^*/\dd\mln = 1/h_i$,
\[
\td{I_i}{\mln} = h_i\!\left(\frac12 + \frac{3\mln^*}{4} - \frac{(\mln^*)^3}{4}
\right)\frac{1}{h_i} - \alpha = P_\kappa(\mln^*) - \alpha = J_i(\mln),
\qquad
\td{J_i}{\mln} = \frac34\bigl(1-(\mln^*)^2\bigr)\frac{1}{h_i}
= \frac{\kappa(\mln^*)}{h_i}.
\]
In the exterior both are affine, matching the vanishing kernel; the boundary
checks give one-sided agreement, so $c(\mln) = \partial_\mln^2\ssm$ is
continuous (with kinks in $c'$, hence jumps only in the third derivative). For
Dupire's formula, which needs $\partial_\mln^2\sigma$, the piecewise evaluation
is well defined everywhere.
\end{proof}

%% ======================================================================
\section{Appendix: Verification of the Black--Scholes Cross-Derivatives}
\label{sec:bs-verify}
%% ======================================================================

Under $r = 0$ and, at the collapse points below, the standing $F = S$, the call
is $C = F\,\ncdf(d_1) - K\,\ncdf(d_2)$, and $F\,\npdf(d_1) = K\,\npdf(d_2)$
(exact in the forward measure).

\paragraph{Strike gamma.} From $\partial C/\partial K|_\sigma = -\ncdf(d_2)$
(the vega-weighted $d$-terms cancel by the density identity),
\[
\atfixed{\pdd{C}{K}}{\sigma} = -\npdf(d_2)\,\pd{d_2}{K}
= -\npdf(d_2)\!\left(-\frac{1}{K\sigma\sqrt{T}}\right)
= \frac{\npdf(d_2)}{K\sigma\sqrt{T}}.
\]

\paragraph{Cross-derivative.} Still from
$\partial C/\partial K|_\sigma = -\ncdf(d_2)$,
\[
\pdc{C}{K}{\sigma} = -\npdf(d_2)\,\pd{d_2}{\sigma}
= -\npdf(d_2)\!\left(-\frac{d_1}{\sigma}\right)
= \frac{\npdf(d_2)\,d_1}{\sigma}.
\]
In the assembled formula \eqref{eq:d2K-C} this multiplies
$\partial\sigma/\partial K = (1/K)\partial\sigma/\partial\mln$, so the factor
$1/K$ enters through the chain rule, not the cross-derivative.

\paragraph{Volga.} From vega $= F\sqrt{T}\,\npdf(d_1)$ and
$\partial d_1/\partial\sigma = -d_2/\sigma$,
\[
\pdd{C}{\sigma} = F\sqrt{T}\,\npdf(d_1)(-d_1)\pd{d_1}{\sigma}
= F\sqrt{T}\,\npdf(d_1)\,\frac{d_1 d_2}{\sigma}.
\]

\paragraph{Vanna.} This is the one cross-derivative entering the total gamma
\eqref{eq:Gadj}. Under $r=0$ (and $F=S$, flagged), from
$\Delta = \partial C/\partial S = \ncdf(d_1)$ and
$\partial d_1/\partial\sigma = -d_2/\sigma$,
\[
\Vanna = \pdc{C}{S}{\sigma} = \npdf(d_1)\,\pd{d_1}{\sigma}
= -\frac{\npdf(d_1)\,d_2}{\sigma},
\]
equal by the commuting of mixed partials ($C\in C^2$) to
$\partial^2 C/\partial\sigma\,\partial S$.

\paragraph{The denominator cancellation.} The $K$-to-$\mln$ transition uses
$K\,\partial_K\sigma = \partial_\mln\sigma$ and
$K^2\,\partial_K^2\sigma = \partial_\mln^2\sigma - \partial_\mln\sigma$; the
$-\sigma T\,\partial_\mln\sigma$ from the curvature substitution cancels the
$+\sigma T\,\partial_\mln\sigma$ from $2 d_1\sqrt{T} = -2\mln/\sigma + \sigma T$,
yielding the clean $D(\mln,T)$ of \eqref{eq:D-def}.
%% ============================================================================
%% Vol I, Chapter 3: Routed Sensitivities I (Total Greeks & P&L)
%% Governed by NOTATION_CONTRACT.md; macros from the shared preamble only.
%% Label namespaces: *:spine:* (framework, reserved) and *:ch3:* (this body).
%% ============================================================================

%% Chapter-local evaluation notation: \evalat{f}{a} typesets (f)(a) ---
%% point evaluation at a state. Deliberately distinct from the held-fixed
%% constraint bar \atfixed, which the contract (S5.4) and this chapter's
%% spine reserve for genuine constraints (e.g. |_xi in def:spine:hook);
%% the bar is never used for point evaluation in this chapter.
\newcommand{\evalat}[2]{\bigl(#1\bigr)(#2)}

\chapter{Routed Sensitivities I: Total Greeks and Profit-and-Loss}
\label{ch:routed-one}

A book is long or short three things through a single vanilla position: spot,
volatility, and the \emph{coupling} between them that the surface enforces.
The position's profit and loss over a day is one number; the desk's job is to
say, before the day, what that number will be for a given move, and, after the
day, why it was what it was. Both tasks need the same two ingredients, and
this chapter builds each and then composes them. The first ingredient is the
correct sensitivity: when the spot moves, the forward moves, the surface
coordinates move, and the implied volatility of every fixed strike moves with
them --- so the derivative that predicts P\&L is a \emph{routed} total, not a
frozen partial. The second ingredient is the correct bookkeeping: an exact
algebraic identity that splits a price difference between two risk states into
first-, second-, and third-order contributions with nothing left over. The
routed calculus supplies the \emph{content} of the coefficients; the identity
supplies the \emph{form} of the decomposition. Their composition is the
attribution a risk system should print every night.

%% ============================================================================
%% Opening framework sections: the routed-derivative operator, the
%% spot-volatility specialization, and the Vol II specialization hook.
%% Vol II Ch 2 specializes the same operator through the hook marked in
%% Section \ref{sec:spine:hook}.
%% ============================================================================

\section{The Routed-Derivative Framework}\label{sec:spine:framework}

A Greek is a derivative of a value function. What decides every
derivation in this chapter --- and every derivation in the multi-name theory
of Volume~II --- is the path along which the derivative is taken. The
classical Greeks are partial derivatives: they bump one argument of the value
function while holding every other argument fixed. But the arguments of the
value function are not independent in the market. When the spot moves, the
forward moves with it; when the forward moves, the implied volatility read off
the surface moves, because both the strike log-moneyness $\mln = \ln(K/F)$ and
the forward log-level $\lF = \ln(F/\Fref)$ move; and, one link deeper, when
the constituent volatilities of a basket move, a quoting convention may force
the implied correlation to move as well. A sensitivity that is to predict
profit and loss must differentiate \emph{along these routes}, not across them.

This section states the routing calculus once, at the level of elementary
multivariable analysis: gradients, Hessians, and the chain rule, applied to a
finite chain of intermediate variables. Everything is proved. The two
specializations of the two volumes --- the spot--volatility route
$S \to F \to \sigma \to V$ of this chapter, and the deeper route
$S \to F \to \sigma \to \rho \to V$ carrying the shadow Greeks of
Volume~II, Chapter~2 --- are instances of a single operator. The first is
derived in Section~\ref{sec:spine:spot-vol}; the insertion point of the second
is marked explicitly in Section~\ref{sec:spine:hook}.

\subsection{Standing assumptions and derivative notation}\label{sec:spine:assumptions}

\begin{assumption}[Standing assumptions]\label{ass:spine:standing}
Throughout this chapter:
\begin{enumerate}
  \item \emph{Rates and dividends.} Interest rates and dividend yields are
        zero. Consequently the forward map is the identity, $F(S) = S$.
        Nevertheless, every derivation below carries a general forward map
        $F \colon \mathbb{R}_+ \to \mathbb{R}_+$, assumed of class $C^2$ and
        strictly increasing, and invokes the collapse $F = S$ only at
        explicitly flagged points. The routed structure --- not the special
        case --- is the object of the theory.
  \item \emph{Surface decomposition.} The implied volatility of the option
        with strike $K$ and expiry $T$ is
        \begin{equation}\label{eq:spine:decomp}
          \sigma(\mln, \lF, T) \;=\; \satm(\lF, T) + \ssm(\mln, T),
          \qquad \ssm(0, T) = 0,
        \end{equation}
        with $\mln = \ln(K/F)$ and $\lF = \ln(F/\Fref)$, where $\Fref$ is the
        fixed reference forward anchoring the surface. The normalisation
        $\ssm(0,T) = 0$ makes the decomposition unique and makes $\satm$ the
        at-the-money volatility: $\sigma(0, \lF, T) = \satm(\lF, T)$ exactly.
  \item \emph{Regularity.} The value function and every link map introduced
        below are of class $C^2$ on the open sets where they are used. In
        particular $\satm(\cdot, T)$ and $\ssm(\cdot, T)$ are twice
        continuously differentiable in their first arguments.
  \item \emph{Fixed expiry.} $T$ is fixed throughout; all sensitivities are
        instantaneous sensitivities at fixed $T$. Time decay is treated
        separately and does not interact with the routing calculus.
\end{enumerate}
\end{assumption}

The derivative notation is normative and carries the entire framework.

\begin{definition}[Derivative notation]\label{def:spine:notation}
Let $f$ be a function of several declared arguments.
\begin{enumerate}
  \item $\pd{f}{x}$ denotes the \emph{partial} derivative of $f$ with respect
        to its argument $x$, all other declared arguments held fixed.
  \item $\td{f}{x}$ denotes the \emph{total} (routed) derivative of $f$ with
        respect to $x$: the derivative of the composite function obtained by
        substituting into $f$ every dependence of its arguments on $x$ along
        the declared route.
  \item When the variables held fixed are determined not by an argument list
        but by a constraint or a quoting convention, the held-fixed quantity
        is named by a subscripted bar: $\atfixed{\td{\rho}{\sigma}}{\xi}$ is
        the derivative of $\rho$ with respect to $\sigma$ under the constraint
        that $\xi$ is invariant.
\end{enumerate}
For a function $V \colon \mathbb{R}^{q} \to \mathbb{R}$ of class $C^2$,
$\nabla V \in \mathbb{R}^{q}$ denotes its gradient (a column vector) and
$\nabla^2 V \in \mathbb{R}^{q \times q}$ its Hessian; inner products are
written $x^{\mathsf T} y$. By Schwarz's theorem, $\nabla^2 V$ is symmetric.
\end{definition}

\subsection{Routes and routed derivatives}\label{sec:spine:routes}

\begin{definition}[Route]\label{def:spine:route}
Let $I \subset \mathbb{R}$ be an open interval and let $s \in I$ be the
\emph{driver}. A \emph{route of depth $p$} over $I$ is an ordered list of
\emph{link maps}
\[
  u_1 \colon I \to \mathbb{R}, \qquad
  u_j \colon I \times \mathbb{R}^{j-1} \to \mathbb{R} \quad (j = 2, \ldots, p),
\]
each of class $C^2$, together with a \emph{value function}
$V \colon I \times \mathbb{R}^{p} \to \mathbb{R}$ of class $C^2$. The $j$-th
\emph{intermediate} depends on the driver and on the intermediates preceding
it, and on nothing else; this triangular ordering is what makes the route a
route.
\end{definition}

\begin{definition}[Routed evaluation]\label{def:spine:evaluation}
The \emph{routed intermediates} $\hat u_j \colon I \to \mathbb{R}$ are defined
recursively by
\[
  \hat u_1(s) = u_1(s), \qquad
  \hat u_j(s) = u_j\bigl(s, \hat u_1(s), \ldots, \hat u_{j-1}(s)\bigr)
  \quad (j = 2, \ldots, p).
\]
The \emph{augmented state} is the vector
$w(s) = \bigl(s, \hat u_1(s), \ldots, \hat u_p(s)\bigr)^{\mathsf T}
\in \mathbb{R}^{p+1}$, and the \emph{routed value} is
$\hat V(s) = V\bigl(w(s)\bigr)$. The \emph{routed derivatives} of $V$ along
the route are $\td{\hat V}{s}$ and $\frac{\dd^2 \hat V}{\dd s^2}$; when the
route is clear from context they are written $\td{V}{s}$ and
$\frac{\dd^2 V}{\dd s^2}$.
\end{definition}

Every Greek in these volumes is a routed derivative for some route. The
classical (partial) Greeks are the degenerate case in which every link map is
constant in the driver, so that the route transmits nothing.

\subsection{The routed differentiation operator}\label{sec:spine:operator}

\begin{theorem}[Routed chain rule]\label{thm:spine:routed-chain}
Let a route of depth $p$ with value function $V$ be given, and let all maps be
of class $C^2$ as in Definition~\ref{def:spine:route}. Then:
\begin{enumerate}
  \item[(i)] Each routed intermediate $\hat u_j$ is of class $C^2$ on $I$, and
    for all $s \in I$,
    \begin{equation}\label{eq:spine:link-recursion}
      \td{\hat u_j}{s}
      \;=\; \pd{u_j}{s} + \sum_{i=1}^{j-1} \pd{u_j}{u_i}\,\td{\hat u_i}{s},
    \end{equation}
    all partial derivatives of $u_j$ being evaluated at
    $\bigl(s, \hat u_1(s), \ldots, \hat u_{j-1}(s)\bigr)$.
  \item[(ii)] The routed value $\hat V$ is of class $C^2$ on $I$, and
    \begin{equation}\label{eq:spine:first-order}
      \td{\hat V}{s}
      \;=\; \nabla V\bigl(w(s)\bigr)^{\mathsf T} w'(s)
      \;=\; \pd{V}{s} + \sum_{j=1}^{p} \pd{V}{u_j}\,\td{\hat u_j}{s}.
    \end{equation}
  \item[(iii)] The second routed derivative is
    \begin{equation}\label{eq:spine:second-order}
      \frac{\dd^2 \hat V}{\dd s^2}
      \;=\; w'(s)^{\mathsf T}\, \nabla^2 V\bigl(w(s)\bigr)\, w'(s)
        \;+\; \nabla V\bigl(w(s)\bigr)^{\mathsf T} w''(s),
    \end{equation}
    where the components of $w''(s)$ are $0$ (the driver component is $s$
    itself) and the second routed derivatives
    $\frac{\dd^2 \hat u_j}{\dd s^2}$, each of which is obtained by applying
    formula~\eqref{eq:spine:second-order} to $u_j$ in the role of the value
    function over the sub-route of depth $j-1$.
\end{enumerate}
\end{theorem}

\begin{proof}
\emph{(i).} By induction on $j$. For $j = 1$, $\hat u_1 = u_1$ is $C^2$ by
hypothesis and \eqref{eq:spine:link-recursion} is vacuous beyond
$\td{\hat u_1}{s} = \pd{u_1}{s}$. Suppose $\hat u_1, \ldots, \hat u_{j-1}$ are
$C^2$. Then $s \mapsto \bigl(s, \hat u_1(s), \ldots, \hat u_{j-1}(s)\bigr)$ is
a $C^2$ map from $I$ into $I \times \mathbb{R}^{j-1}$, and $\hat u_j$ is its
composition with the $C^2$ map $u_j$, hence $C^2$
(compositions of $C^2$ maps are $C^2$). The chain rule for a $C^1$ function of
several variables composed with a $C^1$ curve gives exactly
\eqref{eq:spine:link-recursion}.

\emph{(ii).} By (i), $w$ is a $C^2$ curve in $\mathbb{R}^{p+1}$, so
$\hat V = V \circ w$ is $C^2$, and the same chain rule gives
$\td{\hat V}{s} = \nabla V(w(s))^{\mathsf T} w'(s)$. Writing out the
components of $\nabla V$ and of
$w'(s) = \bigl(1, \td{\hat u_1}{s}, \ldots, \td{\hat u_p}{s}\bigr)^{\mathsf T}$
yields the displayed sum.

\emph{(iii).} Index the arguments of $V$ by $a = 0, 1, \ldots, p$, with $a=0$
the driver slot, and write $\partial_a V$ for the corresponding partial
derivative. Differentiating the component form of
\eqref{eq:spine:first-order},
\[
  \frac{\dd^2 \hat V}{\dd s^2}
  = \sum_{a=0}^{p} \td{}{s}\Bigl[\partial_a V\bigl(w(s)\bigr)\Bigr]\, w_a'(s)
    + \sum_{a=0}^{p} \partial_a V\bigl(w(s)\bigr)\, w_a''(s).
\]
Each $\partial_a V$ is $C^1$ because $V$ is $C^2$, so the chain rule applies
to it along the curve $w$:
$\td{}{s}\bigl[\partial_a V(w(s))\bigr]
 = \sum_{b=0}^{p} \partial_b \partial_a V\bigl(w(s)\bigr)\, w_b'(s)$.
Substituting,
\[
  \frac{\dd^2 \hat V}{\dd s^2}
  = \sum_{a,b=0}^{p} \partial_b \partial_a V\bigl(w(s)\bigr)\,
    w_a'(s)\, w_b'(s)
    + \sum_{a=0}^{p} \partial_a V\bigl(w(s)\bigr)\, w_a''(s),
\]
which is \eqref{eq:spine:second-order} in matrix form. The component $w_0(s) = s$
has $w_0'' = 0$; the remaining components of $w''$ are the
$\frac{\dd^2 \hat u_j}{\dd s^2}$. Finally, $\hat u_j$ is itself a routed value:
it is the composition of $u_j$ with the augmented state of the sub-route
$(u_1, \ldots, u_{j-1})$, so formula \eqref{eq:spine:second-order} applied to
$u_j$ over that sub-route computes $\frac{\dd^2 \hat u_j}{\dd s^2}$. The
recursion terminates because each application strictly decreases the depth.
\end{proof}

\begin{remark}[One operator, applied recursively]\label{rem:spine:self-similar}
Part (iii) is the structural heart of the framework: the second routed
derivative of the value requires the second routed derivatives of the
intermediates, and these are computed \emph{by the same formula} one level
down. Nothing new is needed when a route is extended by a link --- one more
application of Theorem~\ref{thm:spine:routed-chain} to the appended link map,
and formulas \eqref{eq:spine:first-order}--\eqref{eq:spine:second-order}
absorb it. This is precisely how the shadow-Greek route of Volume~II plugs in
(Section~\ref{sec:spine:hook}).
\end{remark}

\begin{corollary}[Depth-one routed derivatives]\label{cor:spine:depth-one}
Let $V = V(s, u)$ be $C^2$ and let $u = u(s)$ be $C^2$. Then
\begin{align}
  \td{V}{s} &= \pd{V}{s} + \pd{V}{u}\,\td{u}{s},
    \label{eq:spine:d1}\\[2pt]
  \frac{\dd^2 V}{\dd s^2}
  &= \pdd{V}{s}
   + 2\,\pdc{V}{s}{u}\,\td{u}{s}
   + \pdd{V}{u}\left(\td{u}{s}\right)^{\!2}
   + \pd{V}{u}\,\frac{\dd^2 u}{\dd s^2}.
    \label{eq:spine:d2}
\end{align}
\end{corollary}

\begin{proof}
Take $p = 1$ in Theorem~\ref{thm:spine:routed-chain}. Here
$w(s) = (s, u(s))^{\mathsf T}$, $w'(s) = (1, u'(s))^{\mathsf T}$,
$w''(s) = (0, u''(s))^{\mathsf T}$. Formula \eqref{eq:spine:first-order} gives
\eqref{eq:spine:d1}. Expanding
$w'^{\mathsf T} (\nabla^2 V) w'
 = \pdd{V}{s} + \bigl(\pdc{V}{s}{u} + \pdc{V}{u}{s}\bigr) u'
 + \pdd{V}{u} (u')^2$
and invoking Schwarz's theorem to merge the two cross terms, formula
\eqref{eq:spine:second-order} gives \eqref{eq:spine:d2}.
\end{proof}

\begin{remark}[What a specialization must supply]\label{rem:spine:interface}
To instantiate the operator, a specialization supplies exactly three items:
(a) the driver $s$; (b) the ordered intermediates $u_1, \ldots, u_p$ with the
partial derivatives of each link map; (c) the partial derivatives of the value
function in its declared arguments. Theorem~\ref{thm:spine:routed-chain} then
\emph{outputs} the total Greeks. In this chapter the list (b) is
$(F, \sigma)$; in Volume~II, Chapter~2 it is $(F_i, \sigma_i, \rho)$ ---
same driver role, same operator, one more link.
\end{remark}

%% ===========================================================================
\section{First Specialization: the Spot--Volatility Route}\label{sec:spine:spot-vol}
%% ===========================================================================

\subsection{The route}\label{sec:spine:the-route}

Fix an option with strike $K$ and expiry $T$. The driver is the spot $S$. The
intermediates, in route order, are the forward $F = F(S)$ and the implied
volatility $\sigma$, which depends on $S$ only through $F$, via the two
moneyness coordinates of the surface \eqref{eq:spine:decomp}:
\[
  \mln = \ln(K/F), \qquad \lF = \ln(F/\Fref), \qquad
  \sigma = \satm(\lF, T) + \ssm(\mln, T).
\]
The value function is $V = V(S, \sigma)$, the option value as a function of
spot and volatility. (Under Assumption~\ref{ass:spine:standing}(1) the direct
spot argument of $V$ and the forward coincide numerically --- under zero
rates, $F = S$ --- but the two roles stay separate in the route so that the
$F$-links remain visible; they are where nonzero carry enters.) The
partial derivatives of $V$ in its declared arguments are the classical Greeks
of the notation contract:
\[
  \Delta = \pd{V}{S}, \quad
  \Gamma = \pdd{V}{S}, \quad
  \vega = \pd{V}{\sigma}, \quad
  \Vanna = \pdc{V}{S}{\sigma}, \quad
  \Volga = \pdd{V}{\sigma}.
\]
Figure~\ref{fig:spine:routing} displays the route, together with the dashed
correlation link that Volume~II appends.

\begin{figure}[ht]
\centering
\begin{tikzpicture}[>=stealth,
  var/.style={draw, rounded corners, minimum height=7mm, minimum width=10mm,
              inner sep=2pt},
  lab/.style={font=\small, fill=white, inner sep=1.5pt}]
  \node[var] (S)   at (0,0)      {$S$};
  \node[var] (F)   at (2.4,0)    {$F$};
  \node[var] (m)   at (4.9,1.25) {$\mln$};
  \node[var] (l)   at (4.9,-1.25){$\lF$};
  \node[var] (sig) at (7.4,0)    {$\sigma$};
  \node[var] (V)   at (11.6,0)   {$V$};
  \node[var, dashed] (rho) at (9.5,-1.6) {$\rho$};
  \draw[->] (S) -- node[lab]{$F'(S)$} (F);
  \draw[->] (F) -- node[lab, pos=0.55]{$-1/F$} (m);
  \draw[->] (F) -- node[lab, pos=0.55]{$1/F$} (l);
  \draw[->] (m) -- node[lab, pos=0.45]{$s(\mln)$} (sig);
  \draw[->] (l) -- node[lab, pos=0.45]{$\volslope$} (sig);
  \draw[->] (sig) -- node[lab]{$\pd{V}{\sigma}=\vega$} (V);
  \draw[->] (S) to[bend left=32] node[lab]{$\pd{V}{S}=\Delta$} (V);
  \draw[->, dashed] (sig) to[bend right=12]
      node[lab, pos=0.55]{$\atfixed{\td{\rho}{\sigma}}{\xi}$} (rho);
  \draw[->, dashed] (rho) to[bend right=12]
      node[lab, pos=0.45]{$\pd{V}{\rho}$} (V);
\end{tikzpicture}
\caption{The routed dependency structure. Solid arrows: the spot--volatility
route of this chapter, each edge labelled by the derivative it transmits.
Dashed arrows: the correlation link appended by the \stickycvc{} convention
--- the specialization hook of Section~\ref{sec:spine:hook}, evaluated in
Volume~II, Chapter~2. Under \stickyrho{} the dashed link transmits nothing.}
\label{fig:spine:routing}
\end{figure}

\subsection{Moneyness derivatives through the forward}\label{sec:spine:moneyness}

\begin{lemma}[Moneyness links]\label{lem:spine:moneyness}
At fixed $K$, $\Fref$, and $T$, the moneyness coordinates satisfy, for all
$F > 0$,
\begin{equation}\label{eq:spine:moneyness-links}
  \pd{\mln}{F} = -\frac{1}{F}, \qquad
  \pd{\lF}{F} = \frac{1}{F}, \qquad
  \pdd{\mln}{F} = \frac{1}{F^2}, \qquad
  \pdd{\lF}{F} = -\frac{1}{F^2}.
\end{equation}
\end{lemma}

\begin{proof}
$\mln = \ln K - \ln F$ gives $\pd{\mln}{F} = -1/F$ and, differentiating again,
$\pdd{\mln}{F} = 1/F^2$. $\lF = \ln F - \ln \Fref$ gives $\pd{\lF}{F} = 1/F$
and $\pdd{\lF}{F} = -1/F^2$.
\end{proof}

The two coordinates respond to the forward with \emph{opposite signs}: when
the forward rises, the option's strike moneyness $\mln$ falls while the
surface level coordinate $\lF$ rises. Every sign asymmetry in the formulas
below descends from this single fact.

Per the contract,
\[
  s(\mln) = \pd{\ssm}{\mln}, \qquad
  c(\mln) = \pdd{\ssm}{\mln}, \qquad
  \volslope = \pd{\satm}{\lF},
\]
the smile skew, the smile curvature, and the ATM slope, the last evaluated at
the current $(\lF, T)$; for equity underlyings $\volslope < 0$, and the
distinction between $\volslope$ and $s(\mln)$ is that of
Remark~\ref{rem:two-slopes}. The route transmits the two through different
links with different signs, and conflating them destroys the delta and gamma
adjustments derived below.

\subsection{Routed volatility derivatives}\label{sec:spine:vol-derivatives}

\begin{proposition}[Volatility through the forward]\label{prop:spine:vol-through-F}
Under Assumption~\ref{ass:spine:standing}, the implied volatility of the fixed
option, viewed as a function of the forward through
\eqref{eq:spine:decomp}, satisfies
\begin{align}
  \td{\sigma}{F}
  &= \frac{1}{F}\Bigl[\,\volslope - s(\mln)\,\Bigr],
    \label{eq:spine:dsig-dF}\\[2pt]
  \frac{\dd^2 \sigma}{\dd F^2}
  &= \frac{1}{F^2}\Bigl[\,\pdd{\satm}{\lF} + c(\mln) + s(\mln) - \volslope\,\Bigr].
    \label{eq:spine:d2sig-dF2}
\end{align}
Through a general $C^2$ forward map $F(S)$, the routed derivatives with
respect to the spot are
\begin{equation}\label{eq:spine:dsig-dS-general}
  \td{\sigma}{S} = \td{\sigma}{F}\,F'(S), \qquad
  \frac{\dd^2 \sigma}{\dd S^2}
  = \frac{\dd^2 \sigma}{\dd F^2}\,\bigl(F'(S)\bigr)^2
  + \td{\sigma}{F}\,F''(S).
\end{equation}
Under zero rates, $F = S$, so $F'(S) = 1$ and $F''(S) = 0$, and
\begin{equation}\label{eq:spine:dsig-dS}
  \td{\sigma}{S} = \frac{\volslope - s(\mln)}{S},
  \qquad
  \frac{\dd^2 \sigma}{\dd S^2}
  = \frac{1}{S^2}\Bigl[\,\pdd{\satm}{\lF} + c(\mln) + s(\mln) - \volslope\,\Bigr].
\end{equation}
\end{proposition}

\begin{proof}
The map $F \mapsto \sigma$ is the sum of the two compositions
$F \mapsto \lF \mapsto \satm(\lF, T)$ and $F \mapsto \mln \mapsto \ssm(\mln, T)$.
By the chain rule and Lemma~\ref{lem:spine:moneyness},
\[
  \td{\sigma}{F}
  = \volslope \cdot \frac{1}{F} + s(\mln)\cdot\Bigl(-\frac{1}{F}\Bigr)
  = \frac{\volslope - s(\mln)}{F},
\]
which is \eqref{eq:spine:dsig-dF}. Differentiate once more in $F$ by the
product rule applied to $F^{-1}\bigl[\volslope - s(\mln)\bigr]$:
\[
  \frac{\dd^2 \sigma}{\dd F^2}
  = -\frac{1}{F^2}\bigl[\volslope - s(\mln)\bigr]
    + \frac{1}{F}\biggl[\pdd{\satm}{\lF}\cdot\frac{1}{F}
      - c(\mln)\cdot\Bigl(-\frac{1}{F}\Bigr)\biggr],
\]
where the bracketed derivatives again use Lemma~\ref{lem:spine:moneyness}
through the chain rule. Collecting over $F^{-2}$ gives
\eqref{eq:spine:d2sig-dF2}. Since $\sigma$ has no direct dependence on $S$
outside of $F$, the composite $S \mapsto \sigma(F(S))$ obeys the univariate
chain and product rules, yielding \eqref{eq:spine:dsig-dS-general}; the
zero-rate collapse $F = S$ (flagged here per
Assumption~\ref{ass:spine:standing}(1)) reduces it to
\eqref{eq:spine:dsig-dS}.
\end{proof}

\begin{remark}[Sign asymmetry]\label{rem:spine:sign-asymmetry}
The skew $s(\mln)$ enters the first derivative
\eqref{eq:spine:dsig-dF} with a minus sign and the second derivative
\eqref{eq:spine:d2sig-dF2} with a plus sign, while $\volslope$ does the
opposite. Both asymmetries are forced by
Lemma~\ref{lem:spine:moneyness}: the coordinate $\mln$ carries
$\pd{\mln}{F} = -1/F$ but $\pdd{\mln}{F} = +1/F^2$, and $\lF$ carries the
mirror-image pair. For an equity surface with $\volslope < 0$ and
$s(\mln) < 0$ near the money, the two routes \emph{partially offset} in
$\td{\sigma}{S}$ and \emph{reinforce} in the curvature bracket of
$\frac{\dd^2\sigma}{\dd S^2}$ --- a structural fact with direct hedging
consequences, quantified in Example~\ref{ex:spine:worked}.
\end{remark}

\subsection{Total delta and total gamma}\label{sec:spine:total-greeks}

\begin{theorem}[Total Greeks along the spot--volatility route]\label{thm:spine:total-greeks}
Under Assumption~\ref{ass:spine:standing}, the routed derivatives of the value
$V = V(S, \sigma)$ along the route of Section~\ref{sec:spine:the-route} are
\begin{align}
  \td{V}{S} &= \Delta + \vega\,\td{\sigma}{S},
    \label{eq:spine:total-delta}\\[2pt]
  \Gadj \;:=\; \frac{\dd^2 V}{\dd S^2}
  &= \Gamma
   + 2\,\Vanna\,\td{\sigma}{S}
   + \Volga\left(\td{\sigma}{S}\right)^{\!2}
   + \vega\,\frac{\dd^2 \sigma}{\dd S^2},
    \label{eq:spine:total-gamma}
\end{align}
with $\td{\sigma}{S}$ and $\frac{\dd^2\sigma}{\dd S^2}$ given by
Proposition~\ref{prop:spine:vol-through-F}. Fully expanded --- under zero
rates, $F = S$ --- these read
\begin{align}
  \td{V}{S}
  &= \Delta + \frac{\vega}{S}\Bigl[\volslope - s(\mln)\Bigr],
    \label{eq:spine:total-delta-expanded}\\[2pt]
  \Gadj
  &= \Gamma
   + \frac{2\,\Vanna}{S}\Bigl[\volslope - s(\mln)\Bigr]
   + \frac{\Volga}{S^2}\Bigl[\volslope - s(\mln)\Bigr]^2 \notag\\
  &\qquad
   + \frac{\vega}{S^2}\Bigl[\pdd{\satm}{\lF} + c(\mln) + s(\mln) - \volslope\Bigr].
    \label{eq:spine:total-gamma-expanded}
\end{align}
\end{theorem}

\begin{proof}
Apply Corollary~\ref{cor:spine:depth-one} with driver $s = S$ and single
effective intermediate $u = \sigma$ (the forward link is already absorbed into
$\td{\sigma}{S}$ and $\frac{\dd^2\sigma}{\dd S^2}$ by
Proposition~\ref{prop:spine:vol-through-F}, and $V$ has no direct $F$
argument). Identifying the partial derivatives of $V$ with the classical
Greeks --- $\pd{V}{S} = \Delta$, $\pd{V}{\sigma} = \vega$,
$\pdd{V}{S} = \Gamma$, $\pdc{V}{S}{\sigma} = \Vanna$,
$\pdd{V}{\sigma} = \Volga$ --- formula \eqref{eq:spine:d1} gives
\eqref{eq:spine:total-delta} and formula \eqref{eq:spine:d2} gives
\eqref{eq:spine:total-gamma}; the factor $2$ on $\Vanna$ is the Schwarz merge
of the two cross-pathways ($S$-then-$\sigma$ and $\sigma$-then-$S$).
Substituting \eqref{eq:spine:dsig-dS} --- the zero-rate collapse $F = S$,
flagged --- yields the expanded forms.
\end{proof}

Each term of \eqref{eq:spine:total-gamma} has one route reading:

\begin{center}
\renewcommand{\arraystretch}{1.35}
\begin{tabular}{@{}lll@{}}
\toprule
Term & Route pathway & Content \\
\midrule
$\Gamma$ & $S \to V$ twice & direct spot convexity \\
$2\,\Vanna\,\td{\sigma}{S}$ & $S \to V$, $S \to F \to \sigma \to V$ (both orders) & spot--vol cross-coupling \\
$\Volga\,\bigl(\td{\sigma}{S}\bigr)^2$ & $S \to F \to \sigma \to V$ twice & vol convexity through the surface slope \\
$\vega\,\frac{\dd^2\sigma}{\dd S^2}$ & curvature of the link itself & vega weighted by routed surface curvature \\
\bottomrule
\end{tabular}
\end{center}

\begin{remark}[Consistency of the two viewpoints]\label{rem:spine:two-views}
The routed derivatives can equivalently be computed by keeping
$(S, \sigma)$ as a two-dimensional state and contracting the full gradient and
Hessian of $V$ against the augmented velocity
$w'(S) = \bigl(1, \td{\sigma}{S}\bigr)^{\mathsf T}$ and acceleration
$w''(S) = \bigl(0, \frac{\dd^2\sigma}{\dd S^2}\bigr)^{\mathsf T}$, as in
Theorem~\ref{thm:spine:routed-chain}(iii). The reduced formulas
\eqref{eq:spine:total-delta}--\eqref{eq:spine:total-gamma} and the
matrix contraction agree identically --- this is a theorem, not a check ---
but computing both in an implementation and verifying agreement is a sharp
diagnostic for surface-model consistency.
\end{remark}

\subsection{Worked example}\label{sec:spine:worked}

\begin{example}[At-the-money call on an equity surface]\label{ex:spine:worked}
Fix $T = 1$, $K = 100$, $\Fref = 100$, and a local surface model
\begin{equation}\label{eq:spine:example-surface}
  \satm(\lF, T) = 0.20 - 0.04\,\lF,
  \qquad
  \ssm(\mln, T) = -0.10\,\mln + 0.05\,\mln^2,
\end{equation}
so that $\volslope = -0.04$, $\pdd{\satm}{\lF} = 0$, $s(0) = -0.10$,
$c(0) = 0.10$, and the normalisation $\ssm(0,T) = 0$ holds. Both slopes are
negative, the equity-index configuration of Remark~\ref{rem:two-slopes}.
Evaluate at spot $S = 100$; under
zero rates, $F = S = 100$, hence $\lF = 0$, $\mln = 0$, and
$\sigma = \satm(0,1) = 0.20$.

\emph{Classical Greeks.} With
$d_1 = -\mln/(\sigma\sqrt{T}) + \sigma\sqrt{T}/2 = 0.1$ and
$d_2 = d_1 - \sigma\sqrt{T} = -0.1$, the Black--Scholes call values are
(six significant figures, derived analytically from the closed forms
$\Delta = \ncdf(d_1)$, $\Gamma = \npdf(d_1)/(S\sigma\sqrt{T})$,
$\vega = S\,\npdf(d_1)\sqrt{T}$, $\Vanna = -\npdf(d_1)\,d_2/\sigma$,
$\Volga = \vega\, d_1 d_2/\sigma$, with $\npdf(0.1) = 0.396953$):
\[
  \Delta = 0.539828, \quad
  \Gamma = 0.0198476, \quad
  \vega = 39.6953, \quad
  \Vanna = 0.198476, \quad
  \Volga = -1.98476.
\]

\emph{Routed volatility derivatives.} By
Proposition~\ref{prop:spine:vol-through-F}, under zero rates ($F = S$):
\[
  \td{\sigma}{S} = \frac{-0.04 - (-0.10)}{100} = 6.0 \times 10^{-4},
  \qquad
  \frac{\dd^2 \sigma}{\dd S^2}
  = \frac{0 + 0.10 + (-0.10) - (-0.04)}{100^2} = 4.0 \times 10^{-6}.
\]
The offset of Remark~\ref{rem:spine:sign-asymmetry} is visible: the smile
route contributes $-s(0)/S = +10^{-3}$ to $\td{\sigma}{S}$ while the ATM route
contributes $\volslope/S = -4 \times 10^{-4}$, cancelling forty per cent of
it; in the curvature bracket, $s(0)$ and $c(0)$ cancel exactly and the
surviving term is $-\volslope > 0$.

\emph{Total Greeks.} By Theorem~\ref{thm:spine:total-greeks}:
\[
  \td{V}{S} = 0.539828 + 39.6953 \times 6.0\times10^{-4}
            = 0.539828 + 0.023817 = 0.563645,
\]
\begin{align*}
  \Gadj &= 0.0198476
        + 2(0.198476)(6.0\times10^{-4})
        + (-1.98476)(6.0\times10^{-4})^2
        + 39.6953 \times 4.0\times10^{-6} \\
        &= 0.0198476 + 0.000238172 - 0.000000715 + 0.000158781
        = 0.0202439.
\end{align*}
The delta adjustment, $+0.023817$, is $4.4\%$ of the classical delta --- a
material hedging correction. The gamma adjustment, $+0.000396$, is $+2.0\%$ of
the classical gamma, and its dominant contributors are the vanna cross-term
and the vega-curvature term, in the ratio $3:2$; the volga term is three
orders of magnitude smaller at this slope.

\emph{Flagged comparison.} Computed under the $\satm = \mathrm{const}$
convention (C4) --- \emph{not} the surface model of
this chapter --- $\volslope = 0$ and $\pdd{\satm}{\lF} = 0$,
giving $\td{\sigma}{S} = 10^{-3}$, $\frac{\dd^2\sigma}{\dd S^2} = 0$,
total delta $0.579523$, and $\Gadj = 0.0202426$. The total deltas differ by
$0.015878 = \vega \cdot \volslope / S$ --- exactly the ATM-route contribution
that the constant-ATM convention deletes. That the two $\Gadj$ values nearly
coincide here is a numerical accident of the parameter choice
($c(0) = -s(0)$, which zeroes the smile part of the curvature bracket), not a
structural fact; the delta discrepancy is the honest measure of what the
convention discards.
\end{example}

\begin{figure}[ht]
\centering
\begin{tikzpicture}
\begin{axis}[
  width=0.86\textwidth, height=6.6cm,
  xlabel={$S$},
  ylabel={volatility},
  legend pos=north east,
  legend cell align=left,
  grid=major,
  domain=80:125, samples=141, no markers,
  every axis plot/.append style={thick}]
  \addplot[black]
    {0.20 - 0.04*ln(x/100) - 0.10*ln(100/x) + 0.05*(ln(100/x))^2};
  \addlegendentry{$\sigma$ seen by the option (routed total)}
  \addplot[black, dashed]
    {0.20 - 0.04*ln(x/100)};
  \addlegendentry{$\satm(\lF,T)$ route}
  \addplot[black, dotted]
    {-0.10*ln(100/x) + 0.05*(ln(100/x))^2};
  \addlegendentry{$\ssm(\mln,T)$ route}
\end{axis}
\end{tikzpicture}
\caption{The two routes of Example~\ref{ex:spine:worked} for the fixed-strike
option $K = 100$, as functions of spot; under zero rates, $F = S$, so
$\mln = \ln(100/S)$ and $\lF = \ln(S/100)$ here. The routed volatility (solid)
is the sum of the ATM route (dashed) and the smile route (dotted). The slopes
of the two components at $S = 100$ have the same sign in moneyness terms
($\volslope < 0$, $s(0) < 0$) yet contribute with \emph{opposite} signs to
$\td{\sigma}{S}$, by Lemma~\ref{lem:spine:moneyness}.}
\label{fig:spine:routes-plot}
\end{figure}

%% ===========================================================================
\section{The Specialization Hook: Convention-Completed Routes}\label{sec:spine:hook}
%% ===========================================================================

The route of Section~\ref{sec:spine:spot-vol} is closed: every intermediate is
a function of the driver by construction. A multi-name book is not closed in
this way. Its value depends on a parameter --- the implied correlation
$\rho$ --- that no market map determines from the spot. What determines it is
a \emph{quoting convention}: a declaration of which quantity the market holds
invariant as volatilities move. A convention closes the system, and a closed
system is a route; the operator of
Theorem~\ref{thm:spine:routed-chain} then applies without modification.

\begin{definition}[Convention-completed route]\label{def:spine:hook}
Let a route of depth $p$ with intermediates $u_1, \ldots, u_p$ be given, and
let $\theta$ be a parameter of the value function
$V = V(s, u_1, \ldots, u_p, \theta)$ that is not a declared function of the
driver. A \emph{quoting convention} for $\theta$ is a choice of a $C^2$
function $g$ and a declaration that the quantity $g(u_1, \ldots, u_p, \theta)$
is invariant along the route; when
$\pd{g}{\theta} \neq 0$, the implicit function theorem determines
$\theta = \theta(u_1, \ldots, u_p)$ locally as a $C^2$ function of the
intermediates. The \emph{convention-completed route} is the route of depth
$p + 1$ obtained by appending $u_{p+1} = \theta$ with this link map. All
routed derivatives of $V$ are then computed by
Theorem~\ref{thm:spine:routed-chain} on the completed route.
\end{definition}

%% ---------------------------------------------------------------------------
%% HOOK (Vol II, Chapter 2 plugs in here) -------------------------------------
%% ---------------------------------------------------------------------------

\begin{remark}[Specialization hook: the correlation link]\label{rem:spine:hook}
This is the insertion point of the shadow-Greek theory. In the two-name
setting of Volume~II, the constituent volatilities are the single-name
at-the-money volatilities
$\sigma_i \equiv \satmi{i}(\lFi{i}, T)$, $i = 1, 2$, where $\lFi{i}$ and
$\Frefi{i}$ extend the ratified $\lF$ and $\Fref$ by a name index;
each $\sigma_i$ is routed from its spot through the $\satm$ route of
Proposition~\ref{prop:spine:vol-through-F} \emph{alone}, with link
derivative $\td{\sigma_i}{S_i} = (\volslope_i / F_i)\, F_i'(S_i)$ and
per-name slope $\volslope_i = \pd{\satmi{i}}{\lFi{i}}$ --- the smile
coordinate transmits nothing here, because $\sigma_i$ is the floating
at-the-money point, not a fixed-strike volatility --- and forward collapse
$F_i = S_i$ under zero rates, flagged as always. Two conventions close the
correlation:
\begin{itemize}
  \item \stickyrho{}: the invariant is $\rho$ itself ($\dd\rho = 0$). The
        appended link is inert, $\td{\rho}{S} = 0$, and the completed route
        collapses to the depth-two route of this chapter, name by name.
  \item \stickycvc{}: the invariant is the basket vol spread
        $\xi = \sbmean - \sB$ ($\dd\xi = 0$). By
        Definition~\ref{def:spine:hook}, this determines
        $\rho = \rho(\sigma_1, \sigma_2, \xi)$ and appends a live link whose
        routed derivative is
        \begin{equation}\label{eq:spine:hook-link}
          \td{\rho}{S}
          \;=\; \sum_{i=1}^{2}
            \atfixed{\pd{\rho}{\sigma_i}}{\xi}\, \td{\sigma_i}{S}.
        \end{equation}
\end{itemize}
The difference between the two completed routes, applied to the first-order
formula \eqref{eq:spine:first-order}, is a single term:
\begin{equation}\label{eq:spine:hook-difference}
  \atfixed{\td{V}{S}}{\stickycvc}
  \;-\;
  \atfixed{\td{V}{S}}{\stickyrho}
  \;=\; \pd{V}{\rho}\,\td{\rho}{S},
\end{equation}
and the \emph{shadow Greeks are by definition this difference}, expressed on
the contract's normalised bases (the $\shdelta$ form below is the symmetric
parallel move under zero rates, $F_i = S_i$, flagged; Volume~II, Chapter~2
states the general per-name form):
\begin{equation}\label{eq:spine:hook-shadow}
  \shvega \;=\; \frac{1}{100}\,\pd{V}{\rho}\,
                \atfixed{\td{\rho}{\sigma}}{\xi},
  \qquad
  \shdelta \;=\; 100\,\volslope\,\shvega .
\end{equation}
They vanish identically under \stickyrho{}. Volume~II, Chapter~2 owns the
evaluation of the link derivatives
$\atfixed{\pd{\rho}{\sigma_i}}{\xi}$ from the basket variance identity, the
second-order (shadow-gamma) terms generated by
Theorem~\ref{thm:spine:routed-chain}(iii) on the completed route, and the
$\Cega$ normalisation; the present section fixes, once and for all, the
operator those computations instantiate and the interface
(Remark~\ref{rem:spine:interface}) they must supply. Nothing else about the
framework changes when the link is appended: this is
Remark~\ref{rem:spine:self-similar} in action.
\end{remark}

%% ------------------------------------------------------------- end of hook

\begin{remark}[Why the hook carries the forward]\label{rem:spine:hook-forward}
The per-name routed derivatives $\td{\sigma_i}{S}$ entering
\eqref{eq:spine:hook-link} contain the factors $F_i'(S_i)/F_i$ of
Proposition~\ref{prop:spine:vol-through-F}. Under the standing zero-rate
assumption these collapse to $1/S_i$ --- under zero rates, $F_i = S_i$,
flagged --- but the shadow-Greek chain
$S \to \sigma \to \rho \to V$ is stated, and must be implemented, through
$F(S)$: the correlation link composes with the forward links, it does not
bypass them.
\end{remark}
%% ============================================================================
\section{The Cubic Profit-and-Loss Identity}\label{sec:ch3:cubic}
%% ============================================================================

The preceding sections settled which derivative is correct. This section
settles how the price difference between two successive risk states
decomposes --- exactly, with no error term --- into contributions
attributable to each order. The answer is an algebraic identity
about cubic polynomials, deliberately model-free: nothing in it knows
about Black--Scholes, surfaces, or routes. The surface content enters only
through the coefficients, in Section~\ref{sec:ch3:composition}.

Throughout this section and the next, $x$ denotes the vector of declared risk
factors of the book, $x_1$ and $x_2$ two successive states (yesterday's close
and today's, or two intraday marks), and $h = x_2 - x_1$ the move. For a state
$x$,
\[
  \Delta_x := \nabla V(x), \qquad \Gamma_x := \nabla^2 V(x),
\]
the gradient and Hessian of the book value $V$ in the declared risk factors.
The subscripted-state symbols $\Delta_x, \Gamma_x$ are scoped to the P\&L
identity and are distinct from the unsubscripted Black--Scholes partials
$\Delta, \Gamma$ of Section~\ref{sec:spine:the-route}; which concrete
derivatives must populate them --- partials or routed totals --- is not a
matter of taste, and is settled in Section~\ref{sec:ch3:composition}.
(Notation-contract exemption, recorded here in full: $\Delta_x$ and
$\Gamma_x$ are chapter-local aliases for the $\nabla V(x)$ and
$\nabla^2 V(x)$ of Definition~\ref{def:spine:notation}; their components
include vegas and, in Volume~II, $\Cega$ as much as deltas, and a
\emph{subscripted-state} $\Delta$ or $\Gamma$ in this chapter never denotes
a Black--Scholes partial.)

\subsection{The univariate identity}\label{sec:ch3:univariate}

Model the book value as a polynomial of degree three in a single risk factor
$x \in \mathbb{R}$. Degree three is the lowest order that carries a
\emph{varying} gamma: $V'''$ is constant (the desk's ``speed''), $\Gamma_x =
V''(x)$ is affine in $x$, and $\Delta_x = V'(x)$ is quadratic. Within this
class every identity below is exact --- not a truncation.

Taylor's theorem at $x_1$ terminates:
\begin{equation}\label{eq:ch3:taylor}
  V(x_2) - V(x_1)
  = \Delta_{x_1} h + \tfrac{1}{2}\,\Gamma_{x_1} h^2 + \tfrac{1}{6}\,V''' h^3 .
\end{equation}
The constant third derivative is observable from the endpoint gammas: by the
fundamental theorem of calculus,
\begin{equation}\label{eq:ch3:speed-from-gammas}
  \Gamma_{x_2} - \Gamma_{x_1} = \int_{x_1}^{x_2} V'''\,\dd x = V'''\,h .
\end{equation}
Substituting \eqref{eq:ch3:speed-from-gammas} into \eqref{eq:ch3:taylor}
eliminates $V'''$:
\begin{equation}\label{eq:ch3:after-elimination}
  V(x_2) - V(x_1)
  = \Delta_{x_1} h + \tfrac{1}{2}\,\Gamma_{x_1} h^2
    + \tfrac{1}{6}\bigl(\Gamma_{x_2} - \Gamma_{x_1}\bigr) h^2 .
\end{equation}
Split the frozen gamma symmetrically about the two states,
$\Gamma_{x_1} = \tfrac{1}{2}(\Gamma_{x_1} + \Gamma_{x_2})
             - \tfrac{1}{2}(\Gamma_{x_2} - \Gamma_{x_1})$,
and collect the two coefficients of $(\Gamma_{x_2}-\Gamma_{x_1})h^2$
($-\tfrac{1}{4} + \tfrac{1}{6} = -\tfrac{1}{12}$):
\begin{equation}\label{eq:ch3:intermediate}
  V(x_2) - V(x_1)
  = \Delta_{x_1} h
    + \tfrac{1}{4}\bigl(\Gamma_{x_1} + \Gamma_{x_2}\bigr) h^2
    - \tfrac{1}{12}\bigl(\Gamma_{x_2} - \Gamma_{x_1}\bigr) h^2 .
\end{equation}
Finally, because $\Gamma_x$ is affine, the trapezoidal rule integrates it
exactly:
\begin{equation}\label{eq:ch3:trapezoid}
  \Delta_{x_2} - \Delta_{x_1}
  = \int_{x_1}^{x_2} \Gamma_x \,\dd x
  = \tfrac{1}{2}\bigl(\Gamma_{x_1} + \Gamma_{x_2}\bigr) h .
\end{equation}
Substituting \eqref{eq:ch3:trapezoid} into \eqref{eq:ch3:intermediate}
absorbs the average gamma into the average delta and yields the identity that
carries the rest of the chapter.

\begin{theorem}[Fundamental P\&L identity]\label{thm:ch3:fundamental}
For any polynomial $V$ of degree at most three and any
$x_1, x_2 \in \mathbb{R}$, with $h = x_2 - x_1$,
\begin{equation}\label{eq:ch3:fundamental}
  \boxed{\;
  V(x_2) - V(x_1)
  = \frac{1}{2}\bigl(\Delta_{x_1} + \Delta_{x_2}\bigr)\, h
  \;-\; \frac{1}{12}\bigl(\Gamma_{x_2} - \Gamma_{x_1}\bigr)\, h^2 .
  \;}
\end{equation}
The identity is exact; no remainder term exists.
\end{theorem}

\begin{proof}
The chain \eqref{eq:ch3:taylor}--\eqref{eq:ch3:trapezoid} above, each step an
identity for cubics.
\end{proof}

Read as a desk statement: the P\&L is the move times the \emph{average} of the
entry and exit deltas, minus one-twelfth of the gamma change times the move
squared. Regrouping exposes the three attributable orders:
\begin{equation}\label{eq:ch3:three-orders}
  V(x_2) - V(x_1)
  = \underbrace{\Delta_{x_1} h}_{\text{1st order}}
  + \underbrace{\frac{\Delta_{x_2} - \Delta_{x_1}}{2}\, h}_{\text{2nd order}}
  - \underbrace{\frac{1}{12}\bigl(\Gamma_{x_2} - \Gamma_{x_1}\bigr)
    h^2}_{\text{3rd order}} .
\end{equation}
The first order is the delta P\&L on the entry hedge; the second is the
half-of-delta-change term --- gamma P\&L in its endpoint form; the third is
the correction for the fact that gamma itself moved.

\begin{remark}[Constant-gamma reduction, and where the classical formula
breaks]\label{rem:ch3:constant-gamma}
If $\Gamma_{x_1} = \Gamma_{x_2} = \Gamma_{x}$ (degree at most two), the
third-order term vanishes and two familiar forms coincide:
\[
  V(x_2) - V(x_1)
  = \frac{\Delta_{x_2} + \Delta_{x_1}}{2}\, h
  = \Delta_{x_1} h + \tfrac{1}{2}\,\Gamma_{x}\, h^2 ,
\]
the trapezoid on delta and the classical delta--gamma expansion. Their
equivalence holds \emph{only} at constant gamma. When gamma is affine, the
trapezoid rule \eqref{eq:ch3:trapezoid} shows the average-delta form silently
replaces $\Gamma_{x}$ by the average $\tfrac{1}{2}(\Gamma_{x_1}+\Gamma_{x_2})$
--- and the $-\tfrac{1}{12}$ correction of \eqref{eq:ch3:fundamental} must
then be carried. Booking the classical formula against endpoint deltas while
gamma is drifting is a standing source of spurious unexplained P\&L; the
identity prices that error exactly.
\end{remark}

\subsection{The multivariate identity}\label{sec:ch3:multivariate}

A real book has many risk factors --- spots, volatilities by name and bucket,
and, from Volume~II onward, correlations. Let $x_1, x_2 \in \mathbb{R}^n$ and
$h = x_2 - x_1$, and let $V$ be a polynomial of total degree at most three on
$\mathbb{R}^n$. Restrict $V$ to the segment between the states:
\[
  V_h(t) = V\bigl(x_1 + t\,h\bigr), \qquad t \in [0,1],
  \qquad V_h(0) = V(x_1),\quad V_h(1) = V(x_2)
\]
(the letter $g$ stays reserved for the quoting-convention function of
Definition~\ref{def:spine:hook}, so the segment restriction gets its own
symbol), with derivatives, by the chain rule,
\[
  V_h'(t) = \Delta_{x_1 + t h}^{\mathsf T}\, h,
  \qquad
  V_h''(t) = h^{\mathsf T}\, \Gamma_{x_1 + t h}\, h .
\]
Since $V$ has degree three, $V_h$ is a univariate cubic in $t$, and
Theorem~\ref{thm:ch3:fundamental} applied to $V_h$ on $[0,1]$ (driver $t$,
move $1$) translates term by term:

\begin{corollary}[Multivariate P\&L identity]\label{cor:ch3:multivariate}
For any polynomial $V$ of total degree at most three on $\mathbb{R}^n$ and any
states $x_1, x_2$,
\begin{equation}\label{eq:ch3:multivariate}
  \boxed{\;
  V(x_2) - V(x_1)
  = \Bigl(\frac{\Delta_{x_1} + \Delta_{x_2}}{2}\Bigr)^{\!\mathsf T} h
  \;-\; \frac{1}{12}\, h^{\mathsf T}
        \bigl(\Gamma_{x_2} - \Gamma_{x_1}\bigr)\, h .
  \;}
\end{equation}
\end{corollary}

Here $\Delta_x$ is the full vector of first-order sensitivities of the book
--- deltas, vegas, and in Volume~II the correlation sensitivity $\Cega$ ---
and $\Gamma_x$ the matrix of second-order ones: gammas on the diagonal
blocks, vanna and volga off and on the volatility block, cross-gammas between
names. By Schwarz's theorem $\Gamma_x$ is symmetric; an asymmetric Hessian in
an implementation is a bug, not a convention (Section~\ref{sec:ch3:protocol}).

\begin{remark}[Hessian coverage]\label{rem:ch3:hessian-coverage}
The Hessian has $n(n+1)/2$ independent entries. At $n = 50$ risk factors that
is $1{,}275$ second-order sensitivities, most of them cross-gammas. Any
implementation populates only the material entries; the
neglected remainder does not disappear --- it lands, move by move, in the
unexplained residual, which is why the residual must be monitored
rather than assumed small (Section~\ref{sec:ch3:protocol}).
\end{remark}

\subsection{Degree three as a working boundary}\label{sec:ch3:why-cubic}

Degree three is the lowest order at which the expansion
carries delta, gamma, \emph{and} the leading variation of gamma --- the three
quantities a desk hedges, rebalances, and disputes. For a real
book the value function is not a polynomial, and the identity
\eqref{eq:ch3:multivariate}, evaluated with measured endpoint sensitivities,
explains the reprice only up to a residual
\begin{equation}\label{eq:ch3:residual}
  R(x_1, x_2)
  := \bigl[V(x_2) - V(x_1)\bigr]
   - \Bigl(\frac{\Delta_{x_1} + \Delta_{x_2}}{2}\Bigr)^{\!\mathsf T} h
   + \frac{1}{12}\, h^{\mathsf T}\bigl(\Gamma_{x_2} - \Gamma_{x_1}\bigr)\, h ,
\end{equation}
the quartic-and-above content of the price function over the move. For
well-behaved products over realistic moves, $R$ is small --- the worked
example of Section~\ref{sec:ch3:worked} puts it near one basis point of the
P\&L for a ten-percent spot move. The doctrine for products whose $R$ is
structurally large --- what they are, why they cannot be risk-managed inside
this bookkeeping, and what a desk should do about them --- is developed in
Section~\ref{sec:ch3:toxic}.

%% ============================================================================
\section{Ex Ante and Ex Post Attribution}\label{sec:ch3:antepost}
%% ============================================================================

The identity \eqref{eq:ch3:multivariate} needs sensitivities at \emph{both}
states. Before the move only $x_1$ has been risk-run; after it, both have.
This asymmetry splits attribution into an ex ante estimate and an ex post
decomposition, and their difference is itself an exactly computable
diagnostic.

\subsection{Ex ante: the projected gradient}\label{sec:ch3:exante}

At $x_1$ the desk holds $\Delta_{x_1}$ and (the computed part of)
$\Gamma_{x_1}$. To use the average-delta form, project the exit gradient with
the available Hessian:
\begin{equation}\label{eq:ch3:projected-gradient}
  \tilde{\Delta}_{x_2} := \Delta_{x_1} + \Gamma_{x_1}\, h .
\end{equation}
This is the gradient transported to $x_2$ under a locally-constant-gamma
assumption --- exactly the delta a trader pre-computes when laddering hedges
against a gap scenario. The resulting estimate collapses to the classical
delta--gamma form:
\begin{equation}\label{eq:ch3:exante}
  V(x_2) - V(x_1)
  \;\approx\;
  \Bigl(\frac{\tilde{\Delta}_{x_2} + \Delta_{x_1}}{2}\Bigr)^{\!\mathsf T} h
  = \underbrace{\Delta_{x_1}^{\mathsf T} h}_{\text{1st order}}
  + \underbrace{\tfrac{1}{2}\, h^{\mathsf T} \Gamma_{x_1}\,
      h}_{\text{2nd order}} .
\end{equation}
Two qualifications. First,
$\Gamma_{x_1}$ is never complete (Remark~\ref{rem:ch3:hessian-coverage}), so
$\tilde{\Delta}_{x_2}$ inherits the coverage gaps of the Hessian. Second, the
construction is a control variate, not a dressed-up forecast: the cubic
polynomial is the control whose increment is known in closed form, the
repriced P\&L is the target, and the residual against reprice is the quantity
whose smallness must be demonstrated, not assumed. Using ex ante Greeks does
not inflate the residual, but the estimate itself misses the reprice by a
\emph{deterministic, signed} error --- the reconciliation gap $\mathcal{E}$
of Section~\ref{sec:ch3:reconciliation}, the truncated third-order term plus
the delta projection error --- which dominates the cubic residual by
two orders of magnitude in the worked example (a factor of $139$). The only
requirement on the \emph{control} is that it track the true price function
over plausible moves, which is exactly the boundary question of
Section~\ref{sec:ch3:toxic}.

\subsection{Ex post: the exact decomposition}\label{sec:ch3:expost}

Once $x_2$ has been risk-run, $\Delta_{x_2}$ and $\Gamma_{x_2}$ are known and
the identity applies with measured endpoints:
\begin{equation}\label{eq:ch3:expost}
  V(x_2) - V(x_1)
  = \underbrace{\Delta_{x_1}^{\mathsf T} h}_{\text{1st order}}
  + \underbrace{\Bigl(\frac{\Delta_{x_2} -
      \Delta_{x_1}}{2}\Bigr)^{\!\mathsf T} h}_{\text{2nd order}}
  - \underbrace{\frac{1}{12}\, h^{\mathsf T}
      \bigl(\Gamma_{x_2} - \Gamma_{x_1}\bigr) h}_{\text{3rd order}}
  \;+\; R(x_1, x_2),
\end{equation}
with $R$ the cubic residual \eqref{eq:ch3:residual} --- identically zero for
a cubic book, small and monitored for a real one. The ex post form buys one
thing: the second-order term uses the \emph{measured} delta change, not a
Hessian model of it, so Hessian coverage gaps that contaminate
\eqref{eq:ch3:exante} do not contaminate the first two orders here; they are
squeezed into the third-order term and $R$.

\subsection{Reconciliation}\label{sec:ch3:reconciliation}

Subtracting \eqref{eq:ch3:exante} from the exact part of
\eqref{eq:ch3:expost} gives the reconciliation gap
\begin{equation}\label{eq:ch3:reconciliation}
  \mathcal{E}
  := \Bigl(\frac{\Delta_{x_2} -
     \tilde{\Delta}_{x_2}}{2}\Bigr)^{\!\mathsf T} h
   \;-\; \frac{1}{12}\, h^{\mathsf T}
     \bigl(\Gamma_{x_2} - \Gamma_{x_1}\bigr) h ,
\end{equation}
a sum of two separately reportable pieces: the \emph{delta projection error}
--- how far the Hessian-transported gradient landed from the measured one ---
and the \emph{third-order correction} absent from any delta--gamma estimate.

\begin{proposition}[The gap is pure gamma variation]\label{prop:ch3:gap}
For a book that is exactly cubic in its risk factors,
\begin{equation}\label{eq:ch3:gap-collapsed}
  \Delta_{x_2} - \tilde{\Delta}_{x_2}
  = \tfrac{1}{2}\bigl(\Gamma_{x_2} - \Gamma_{x_1}\bigr) h ,
  \qquad\text{hence}\qquad
  \mathcal{E}
  = \frac{1}{6}\, h^{\mathsf T}\bigl(\Gamma_{x_2} - \Gamma_{x_1}\bigr) h .
\end{equation}
\end{proposition}

\begin{proof}
Along the segment, $t \mapsto \Gamma_{x_1 + t h}$ is affine, so
$\Delta_{x_2} - \Delta_{x_1}
 = \int_0^1 \Gamma_{x_1 + t h}\, h \,\dd t
 = \tfrac{1}{2}(\Gamma_{x_1} + \Gamma_{x_2})\, h$;
subtracting $\tilde{\Delta}_{x_2} - \Delta_{x_1} = \Gamma_{x_1} h$ gives the
first claim, and substituting into \eqref{eq:ch3:reconciliation} collapses
$\tfrac{1}{4} - \tfrac{1}{12} = \tfrac{1}{6}$. Sanity check in one dimension:
$\mathcal{E} = \tfrac{1}{6} V''' h^3$ by
\eqref{eq:ch3:speed-from-gammas} --- exactly the cubic Taylor term that the
ex ante estimate \eqref{eq:ch3:exante} truncates. The signs and the limiting
case ($\mathcal{E} = 0$ at constant gamma) both close.
\end{proof}

The reconciliation gap therefore depends \emph{entirely} on the variation of
gamma over the move --- there is no other source. A small $\mathcal{E}$,
day after day, certifies that the constant-gamma projection underlying every
intraday hedge ladder is adequate at the book's actual gamma turnover. A
persistently large $\mathcal{E}$ has exactly two explanations: material
Hessian entries are missing from $\Gamma_{x_1}$, or the book's gamma is
turning over so fast that the position is drifting toward the boundary of
Section~\ref{sec:ch3:toxic}. Both demand action, and the decomposition
\eqref{eq:ch3:reconciliation} says which.

%% ============================================================================
\section{The Composition: Total Greeks inside the Identity}
\label{sec:ch3:composition}
%% ============================================================================

The identity of Section~\ref{sec:ch3:cubic} is agnostic about what
$\Delta_x$ and $\Gamma_x$ are (recall the scoping declared there:
$\Delta_x = \nabla V(x)$ is the full book gradient, not the Black--Scholes
delta, which appears unsubscripted below); the routing calculus of
Sections~\ref{sec:spine:framework}--\ref{sec:spine:spot-vol} is agnostic
about how sensitivities are aggregated into P\&L. This section joins them,
and the joint carries one non-negotiable rule.

\subsection{Which Greeks enter}\label{sec:ch3:which-greeks}

Fix a single name and let the declared risk factor be the spot, $x = S$. The
implied volatility of every position on the book rides the surface: through
the forward $F(S)$, both moneyness coordinates move, and
$\sigma = \satm(\lF,T) + \ssm(\mln,T)$ moves with them
(Proposition~\ref{prop:spine:vol-through-F}). The book value between marks is
therefore the \emph{routed} value $S \mapsto V(S, \sigma(S))$, and the
gradient and Hessian entries that the identity
\eqref{eq:ch3:fundamental} consumes are the routed totals of
Theorem~\ref{thm:spine:total-greeks}:
\[
  \Delta_S^{\mathrm{total}} = \td{V}{S} = \Delta + \vega\,\td{\sigma}{S},
  \qquad
  \Gamma_S^{\mathrm{total}} = \Gadj
  = \Gamma + 2\,\Vanna\,\td{\sigma}{S}
    + \Volga\Bigl(\td{\sigma}{S}\Bigr)^{\!2}
    + \vega\,\frac{\dd^2\sigma}{\dd S^2},
\]
with $\td{\sigma}{S} = \td{\sigma}{F}\,F'(S)$ and its second-order companion
carried through the forward map exactly as in
\eqref{eq:spine:dsig-dS-general}; under zero rates, $F = S$, and the expanded
forms \eqref{eq:spine:total-delta-expanded} and
\eqref{eq:spine:total-gamma-expanded} apply. Substituting into
Theorem~\ref{thm:ch3:fundamental}:

\begin{theorem}[Unified attribution identity]\label{thm:ch3:unified}
For a book cubic in the routed variable $S$, with $h = S_2 - S_1$,
\begin{equation}\label{eq:ch3:unified}
  \boxed{\;
  V(S_2) - V(S_1)
  = \frac{1}{2}\left(
      \evalat{\td{V}{S}}{S_1} + \evalat{\td{V}{S}}{S_2}
    \right) h
  \;-\; \frac{1}{12}\Bigl(\evalat{\Gadj}{S_2} - \evalat{\Gadj}{S_1}\Bigr) h^2 ,
  \;}
\end{equation}
where each total is evaluated on the surface at its own state --- its own
forward, its own $(\mln, \lF)$, its own classical Greeks. The parenthesised
argument denotes point evaluation at the state; the held-fixed bar
$\atfixed{\,\cdot\,}{\xi}$ remains reserved for genuine constraints per the
contract and is never used for evaluation in this chapter.
\end{theorem}

Using partial Greeks in \eqref{eq:ch3:unified} --- classical $\Delta$ and
$\Gamma$ with the smile adjustments dropped --- is not a smaller model; it is
an inconsistent one. The reprice $V(S_2)$ is computed on the moved surface,
so an explanation built from frozen-$\sigma$ sensitivities attributes the
vol-through-spot P\&L to nothing, and it surfaces as a spurious residual that
no amount of Hessian coverage will remove. The worked example below prices
this error: on a ten-percent move it is first order in the move, roughly
$260$ times the honest cubic residual, and nearly twice the total-Greek ex
ante error, of the opposite sign (Figure~\ref{fig:ch3:residuals}).

\subsection{Two equivalent viewpoints}\label{sec:ch3:two-views}

The same P\&L can be booked two ways, and a sound system computes both.

\emph{Option A --- full multivariate.} Keep $x = (S, \sigma)^{\mathsf T} \in
\mathbb{R}^2$ as separate risk factors. Then $\Delta_x = (\Delta,
\vega)^{\mathsf T}$, the Hessian holds $(\Gamma, \Vanna, \Volga)$, and
Corollary~\ref{cor:ch3:multivariate} applies with the realized move
$h = (S_2 - S_1,\; \sigma_2 - \sigma_1)^{\mathsf T}$, where $\sigma_2 -
\sigma_1$ is the volatility move the surface actually delivered. This
viewpoint attributes P\&L to spot and vol \emph{separately} --- the granular
report.

\emph{Option B --- reduced univariate.} Substitute the route into the value
first, reduce to the single driver $S$, and apply
Theorem~\ref{thm:ch3:unified} with total Greeks --- the compact report.

When the realized vol move equals the routed one,
$\sigma_2 - \sigma_1 = \sigma(S_2) - \sigma(S_1)$ along the surface, the two
agree --- but only up to the difference of their cubic residuals, not to
machine precision. Write $R_A$ and $R_B$ for the residual
\eqref{eq:ch3:residual} of each representation; both explanations target the
same reprice, so
\[
  \text{explanation}_A - \text{explanation}_B = R_B - R_A
  = O\bigl(\norm{h}^4\bigr),
\]
the bound following from Proposition~\ref{prop:ch3:quartic} applied to each
representation separately. The agreement is \emph{exact} only when the route
is affine and the book cubic in $(S, \sigma)$: in general the two truncate
\emph{different} quartic content, because a book cubic in $(S, \sigma)$ is
not cubic in the routed $S$ when $\sigma(S)$ is nonlinear.
Remark~\ref{rem:spine:two-views} is exact for \emph{derivatives} --- the
chain rule --- and lifts to P\&L only to fourth order. The worked example of
Section~\ref{sec:ch3:worked} prices the gap: with the vol move set exactly
on-route at $h = 10$, $R_A = 0.000589$, $R_B = 0.000678$, and the A-minus-B
gap is $8.9\times10^{-5}$ --- thirteen percent of the residual itself,
invisible on the report, and fatal only to a test that demands
machine-precision equality (Section~\ref{sec:ch3:daily} files the comparison
accordingly).

Their difference in production, where the realized vol move contains
a component the route did not predict, is itself informative: it is the
book's P\&L on \emph{surface repricing beyond the sticky dynamics} --- vega
P\&L on the residual vol move --- and Option~A isolates it while Option~B
buries it in $R$. But the sensor has a floor: even with the vol move exactly
on-route the gap is $R_B - R_A$, so the A-minus-B difference isolates
unrouted vol P\&L only \emph{above} a baseline of order
$\max(|R_A|, |R_B|)$ --- fourth order in the move, sixteen-folding when the
move doubles, and capable of printing as spurious unrouted vol on a large
day if it is not netted out first. Both are run: A for the attribution
report, B for the compact one, with the A-minus-B gap, read against its
fourth-order floor, as the measure of how much of the day's vol move the
surface dynamics failed to route.

\subsection{Ex ante with total Greeks}\label{sec:ch3:exante-total}

The ex ante machinery of Section~\ref{sec:ch3:exante} composes the same way:
the projected gradient must be transported with the \emph{total} gamma,
\begin{equation}\label{eq:ch3:exante-total}
  \tilde{\Delta}_{S_2}^{\mathrm{total}}
  = \evalat{\td{V}{S}}{S_1} + \evalat{\Gadj}{S_1}\, h,
  \qquad
  V(S_2) - V(S_1)
  \;\approx\;
  \evalat{\td{V}{S}}{S_1}\, h
  + \tfrac{1}{2}\,\evalat{\Gadj}{S_1}\, h^2 .
\end{equation}
This is where the smile adjustment stops being an accounting nicety and
becomes hedge slippage: the delta a trader expects to hold \emph{after} the
move is $\tilde{\Delta}^{\mathrm{total}}$, and every term of $\Gadj$ ---
vanna coupling, volga through the slope squared, vega times routed curvature
--- feeds it. Omit them and the projected delta is systematically wrong by an
amount growing linearly in the move and in the steepness of the surface; on a
desk running real vega this is rehedging cost, paid daily, not a rounding
error.

\begin{remark}[Fixed expiry, and where theta sits]\label{rem:ch3:theta}
Per Assumption~\ref{ass:spine:standing}(4), every sensitivity in this chapter
is instantaneous at fixed $T$; the identity decomposes the P\&L of the
\emph{move}, not of the calendar. In production the state vector $x$ of
Corollary~\ref{cor:ch3:multivariate} carries calendar time as one more
coordinate: theta enters the gradient as a first-order sensitivity, and the
gamma/theta balance the desk lives on is then read directly off the
attribution --- the second-order spot row against the first-order time row.
Two cautions before this is wired in. First, the licence of
Proposition~\ref{prop:ch3:quartic} is an asymptotic in the move: a one-day
step in the time coordinate is a \emph{fixed finite} component of $h$, not a
shrinking one, so the $O(\norm{h}^4)$ residual control does not cover the
time direction --- near expiry, where theta accelerates, $V$ is nothing like
a low-degree polynomial in calendar time, and the time-direction content
lands in $R$ at full size. Second, the gamma-row-versus-theta-row read-off
closes only if the spot--time cross entries of the Hessian --- charm and
colour --- are populated; both blow up near expiry, exactly where the first
caution already bites, so the read-off is weakest precisely where the desk
leans on it hardest. The routing calculus is untouched by all of this,
because time is a driver of its own, not a link on the spot route.
\end{remark}

\begin{remark}[Nonzero carry]\label{rem:ch3:carry}
Nothing in \eqref{eq:ch3:unified} used $F = S$ except through the expanded
forms of the totals. With nonzero rates or dividends the same identity holds
verbatim; only the link derivatives change, via $F'(S) \neq 1$ and
$F''(S) \neq 0$ in \eqref{eq:spine:dsig-dS-general}. The attribution
framework and the surface dynamics are orthogonal by construction --- that is
the point of routing through $F$ rather than collapsing it.
\end{remark}

%% ============================================================================
\section{Worked Example: a Quadratic Surface End to End}
\label{sec:ch3:worked}
%% ============================================================================

The linear surface of Example~\ref{ex:spine:worked} exercises the first-order
routing. A \emph{quadratic} surface exercises everything this chapter has
built: both curvature sources of
Proposition~\ref{prop:spine:vol-through-F} are live, the total gamma varies
with spot, the third-order term of the identity is nonzero, and the
ex ante / ex post machinery has something real to reconcile. All numbers
below are derived analytically under this surface; the
$\satm = \mathrm{const}$ convention (C4), which would zero the $\volslope$ and
ATM-curvature content on display here, is not used.

\begin{example}[Quadratic surface: states, totals, attribution]
\label{ex:ch3:quadratic}
Fix $T = 1$, $K = 100$, $\Fref = 100$, and the surface
\begin{equation}\label{eq:ch3:example-surface}
  \satm(\lF, T) = 0.20 - 0.04\,\lF + 0.03\,\lF^2,
  \qquad
  \ssm(\mln, T) = -0.10\,\mln + 0.05\,\mln^2 ,
\end{equation}
so that, at general $(\lF, \mln)$,
\[
  \volslope = -0.04 + 0.06\,\lF, \qquad
  \pdd{\satm}{\lF} = 0.06, \qquad
  s(\mln) = -0.10 + 0.10\,\mln, \qquad
  c(\mln) = 0.10 ,
\]
and the normalisation $\ssm(0,T) = 0$ holds. This
quadratic surface is a \emph{local test function} for the attribution
mechanics, valid on the moneyness band exercised here (total vol stays
positive, minimum $\approx 0.137$); it is not globally admissible ---
quadratic growth of the smile violates the Lee moment bounds far from the
money and carries static arbitrage there --- and it must not be reused as a
calibration form (Chapter~4 develops the admissibility machinery). The smile
block reproduces
Example~\ref{ex:spine:worked}; the ATM block adds convexity
($\pdd{\satm}{\lF} = 0.06 > 0$, the ATM vol reaccelerating away from the
anchor in both directions) and lets $\volslope$ itself move with the forward.
The book is one European call, long.

\emph{State 1: the anchor.} At $S_1 = 100$ --- under zero rates, $F = S =
100$ --- $\lF = 0$, $\mln = 0$, $\sigma = 0.20$, and the classical
Greeks of Example~\ref{ex:spine:worked} apply unchanged:
$\Delta = 0.539828$, $\Gamma = 0.0198476$, $\vega = 39.6953$,
$\Vanna = 0.198476$, $\Volga = -1.98476$, and value $V(S_1) = 7.96557$. The
routed derivatives, by Proposition~\ref{prop:spine:vol-through-F} with
$F = S$ (flagged),
\[
  \td{\sigma}{S} = \frac{-0.04 - (-0.10)}{100} = 6.0\times10^{-4},
  \qquad
  \frac{\dd^2\sigma}{\dd S^2}
  = \frac{0.06 + 0.10 + (-0.10) - (-0.04)}{100^2}
  = 1.0\times10^{-5} ,
\]
give, by Theorem~\ref{thm:spine:total-greeks},
\[
  \evalat{\td{V}{S}}{S_1} = 0.563645,
  \qquad
  \evalat{\Gadj}{S_1} = 0.0204820 .
\]
Only the curvature differs from the linear surface: the ATM convexity adds
$\vega \cdot 0.06 / S^2 = 2.38\times10^{-4}$ to the $\Gadj = 0.0202439$ of
Example~\ref{ex:spine:worked} --- the entire effect of the new parameter sits
in one term, exactly where the routing says it must.

\emph{State 2: a ten-percent rally.} At $S_2 = 110$ --- under zero rates,
$F = S = 110$, flagged --- the coordinates are $\lF = \ln 1.1 = 0.0953102$
and $\mln = \ln(100/110) = -0.0953102$, so the option has gone in the money
while the surface level has fallen:
\[
  \sigma = 0.206445, \qquad
  \volslope = -0.0342814, \qquad
  s(\mln) = -0.109531 .
\]
Both slopes have decayed in magnitude relative to the anchor --- the
quadratic terms at work. With $d_1 = -\mln/(\sigma\sqrt{T}) +
\sigma\sqrt{T}/2 = 0.564895$ and $d_2 = 0.358450$, the classical Greeks and
routed derivatives are (six significant figures)
\[
  \Delta = 0.713928, \quad
  \Gamma = 0.0149768, \quad
  \vega = 37.4119, \quad
  \Vanna = -0.590528, \quad
  \Volga = 36.6945,
\]
\begin{align*}
  \td{\sigma}{S} &= \frac{-0.0342814 + 0.109531}{110} = 6.84088\times10^{-4},\\
  \frac{\dd^2\sigma}{\dd S^2}
  &= \frac{0.06 + 0.10 - 0.109531 + 0.0342814}{110^2}
  = 7.00416\times10^{-6},
\end{align*}
with value $V(S_2) = 14.5324$, and totals
\[
  \evalat{\td{V}{S}}{S_2} = 0.713928 + 37.4119 \times 6.84088\times10^{-4}
  = 0.739521,
\]
\begin{align*}
  \evalat{\Gadj}{S_2}
  &= 0.0149768
   + 2(-0.590528)(6.84088\times10^{-4}) \\
  &\qquad
   + 36.6945\,(6.84088\times10^{-4})^2
   + 37.4119 \times 7.00416\times10^{-6} \\
  &= 0.0149768 - 0.000807940 + 0.0000171731 + 0.000262035
   = 0.0144481 .
\end{align*}
At the anchor the smile adjustment \emph{adds} $3.2\%$ to
gamma; ten percent higher it \emph{subtracts} $3.5\%$. The driver is vanna,
which changed sign as $d_2$ crossed zero --- in the money, rising vol now
hurts the delta --- and at this moneyness the vanna pathway overwhelms the
(still positive) vega-curvature term. A risk system carrying the anchor's
gamma adjustment as a static markup would have the sign of the correction
wrong after the rally; the adjustment is a routed quantity and must be
re-evaluated with the state, which is the entire operational content of
Theorem~\ref{thm:ch3:unified}.

\emph{Ex post attribution.} With $h = 10$, the identity
\eqref{eq:ch3:unified} decomposes the reprice as
\begin{center}
\renewcommand{\arraystretch}{1.3}
\begin{tabular}{@{}lrl@{}}
\toprule
Order & Contribution & Formula \\
\midrule
1st & $+5.63645$ & $\evalat{\td{V}{S}}{S_1} h$ \\
2nd & $+0.879378$ & $\tfrac{1}{2}\bigl(\evalat{\td{V}{S}}{S_2} -
\evalat{\td{V}{S}}{S_1}\bigr) h$ \\
3rd & $+0.050283$ & $-\tfrac{1}{12}\bigl(\evalat{\Gadj}{S_2} -
\evalat{\Gadj}{S_1}\bigr) h^2$ \\
\midrule
Identity total & $+6.56611$ & \\
Reprice $V(S_2) - V(S_1)$ & $+6.56679$ & \\
Cubic residual $R$ & $+0.000678$ & about $1$ bp of the P\&L \\
\bottomrule
\end{tabular}
\end{center}
The third-order term is \emph{positive} $5$ cents because the total gamma
fell over the rally: the long-gamma profile decayed into the move (the vanna
sign flip again), and the identity books that decay explicitly instead of
leaving it in the residual. The residual itself --- the quartic-and-above
content of Black--Scholes plus surface over a ten-percent move --- is under
seven-hundredths of a cent on a $6.57$ P\&L.

\emph{Ex ante, and the reconciliation.} Before the move, the desk's estimate
\eqref{eq:ch3:exante-total} reads
\[
  0.563645 \times 10 + \tfrac{1}{2}\,(0.0204820)(100) = 6.66056 ,
\]
an overshoot of $-0.093763$ against the reprice. The reconciliation gap
\eqref{eq:ch3:reconciliation} accounts for it: the projected total delta
$\tilde{\Delta}^{\mathrm{total}}_{S_2} = 0.563645 + 0.0204820 \times 10
= 0.768465$ landed $-0.028945$ above the measured $0.739521$ (delta
projection error $\tfrac{1}{2}(-0.028945)(10) = -0.144723$), partially
offset by the third-order correction $+0.050283$, for
$\mathcal{E} = -0.094441$; the remaining $+0.000678$ is $R$. Under an exactly
cubic book, Proposition~\ref{prop:ch3:gap} would put the gap at
$\tfrac{1}{6}\,(\evalat{\Gadj}{S_2}-\evalat{\Gadj}{S_1})h^2 = -0.100566$; the
measured $-0.094441$ differs from it by $0.0061$ --- a tenth of a percent of
the P\&L, quartic-order like $R$ but not equal to it (it is the trapezoid
error on the non-affine gamma path, a different fourth-order functional) ---
which is the proposition holding to exactly the accuracy the cubic model
claims and no more.

\emph{The cost of partial Greeks.} Rerun the ex ante estimate with classical
Greeks --- $\Delta = 0.539828$, $\Gamma = 0.0198476$, smile adjustments
dropped: the estimate is $6.39066$, an error of $+0.176129$, nearly twice the
total-Greek error and of the opposite sign. Worse is the small-move
behaviour: the partial-Greek error is \emph{first order} in $h$, with slope
$\vega\,\td{\sigma}{S} = 0.0238$ per unit of spot ($\pm4.8$ cents at
$h = \pm2$), because the frozen-$\sigma$ delta is the wrong hedge
ratio; the total-Greek ex ante error is third order in $h$ by construction
(Proposition~\ref{prop:ch3:gap}) and is under a tenth of a cent there. A
desk that hedges the partial delta pays that first-order bias on every
rebalance --- this is the delta adjustment of
Example~\ref{ex:spine:worked} recurring as realized slippage rather than as
a Greek.
\end{example}

\begin{figure}[ht]
\centering
\begin{tikzpicture}
\begin{axis}[
  width=0.86\textwidth, height=7.2cm,
  xlabel={spot move $h = S_2 - S_1$},
  ylabel={explanation residual},
  legend pos=south west,
  legend cell align=left,
  grid=major,
  no markers,
  every axis plot/.append style={thick}]
  \addplot[black, dotted] coordinates {
    (-20,-0.1265) (-18,-0.1257) (-16,-0.1332) (-14,-0.1427) (-12,-0.1489)
    (-10,-0.1480) (-8,-0.1376) (-6,-0.1166) (-4,-0.0855) (-2,-0.0457)
    (0,0.0000) (2,0.0482) (4,0.0948) (6,0.1349) (8,0.1638) (10,0.1761)
    (12,0.1669) (14,0.1311) (16,0.0637) (18,-0.0398) (20,-0.1837)
    (22,-0.3722) (24,-0.6089)};
  \addlegendentry{ex ante, partial Greeks $(\Delta, \Gamma)$}
  \addplot[black, dashed] coordinates {
    (-20,0.2230) (-18,0.2003) (-16,0.1667) (-14,0.1286) (-12,0.0912)
    (-10,0.0584) (-8,0.0326) (-6,0.0148) (-4,0.0047) (-2,0.0006)
    (0,0.0000) (2,-0.0007) (4,-0.0056) (6,-0.0194) (8,-0.0471) (10,-0.0938)
    (12,-0.1646) (14,-0.2645) (16,-0.3986) (18,-0.5712) (20,-0.7869)
    (22,-1.0497) (24,-1.3632)};
  \addlegendentry{ex ante, total Greeks
    $(\td{V}{S}, \Gadj)$}
  \addplot[black] coordinates {
    (-20,-0.0023) (-18,-0.0043) (-16,-0.0039) (-14,-0.0027) (-12,-0.0015)
    (-10,-0.0007) (-8,-0.0003) (-6,-0.0001) (-4,0.0000) (-2,0.0000)
    (0,0.0000) (2,0.0000) (4,0.0000) (6,0.0001) (8,0.0002) (10,0.0007)
    (12,0.0016) (14,0.0031) (16,0.0055) (18,0.0089) (20,0.0136)
    (22,0.0195) (24,0.0268)};
  \addlegendentry{ex post cubic identity (total Greeks)}
\end{axis}
\end{tikzpicture}
\caption{Explanation residuals (reprice minus explanation) for the quadratic
surface of Example~\ref{ex:ch3:quadratic}, as functions of the spot move from
the anchor $S_1 = 100$; under zero rates, $F = S$ throughout. The
partial-Greek estimate (dotted) is biased at \emph{first} order in $h$ ---
slope $\vega\,\td{\sigma}{S} = 0.0238$ through the origin --- and its
seemingly benign behaviour at large positive $h$ is an accidental
cancellation, not accuracy. The total-Greek ex ante residual (dashed) is
third order in $h$ near the origin, per
Proposition~\ref{prop:ch3:gap}. The ex post identity residual (solid) is the
cubic residual $R$: quartic near the origin, at most $0.0136$ over a
$\pm20\%$ move, and $0.027$ at the plotted extreme $h = +24$ ---
indistinguishable from zero at the scale of the other two.}
\label{fig:ch3:residuals}
\end{figure}

%% ============================================================================
\section{Attribution Protocol and Structural Checks}\label{sec:ch3:protocol}
%% ============================================================================

The results above are identities, so an implementation can be tested against
them with no tolerance for style: a failure is a defect, not a calibration
issue. The protocol has two layers --- structural checks that must hold to
numerical precision, and a daily reconciliation loop whose outputs are
diagnostics, not pass/fail.

\subsection{Structural identities}\label{sec:ch3:structural}

\begin{enumerate}
  \item \textbf{Constant-gamma limit.} When $\Gamma_{x_1} = \Gamma_{x_2}$,
        the third-order term of \eqref{eq:ch3:expost} vanishes and the
        average-delta form must equal the classical delta--gamma form
        (Remark~\ref{rem:ch3:constant-gamma}), to machine precision.
  \item \textbf{Flat-surface limit.} When $\volslope = 0$, $s(\mln) = 0$,
        $c(\mln) = 0$, and $\pdd{\satm}{\lF} = 0$, all routed adjustments
        vanish: $\td{V}{S} = \Delta$ and $\Gadj = \Gamma$ exactly
        (Theorem~\ref{thm:spine:total-greeks}).
  \item \textbf{Zero-move limit.} $x_2 = x_1$ must return zero in every
        formula --- including \eqref{eq:ch3:reconciliation}, where a nonzero
        value at zero move indicates state leakage between the two risk runs.
  \item \textbf{Hessian symmetry.} $\Gamma_x$ must be symmetric (Schwarz).
        In routed terms, the two cross-pathways that merge into the factor
        $2\,\Vanna$ in \eqref{eq:spine:total-gamma} must be computed
        identically; asymmetry is a bug in one of them.
\end{enumerate}

One comparison deliberately does \emph{not} appear on this list: Option~A
versus Option~B. The two attributions agree only up to
$R_B - R_A = O(\norm{h}^4)$ (Section~\ref{sec:ch3:two-views}), so demanding
machine-precision equality would fail a \emph{correct} implementation; the
comparison is a diagnostic with a tolerance, and it lives in the daily loop
below.

\subsection{The daily loop}\label{sec:ch3:daily}

Each day, book three numbers per portfolio: the ex ante estimate
\eqref{eq:ch3:exante-total} at the prior close; the ex post decomposition
\eqref{eq:ch3:expost} at the current close; and the reprice. From these,
report the reconciliation gap $\mathcal{E}$ split per
\eqref{eq:ch3:reconciliation} into delta projection error and third-order
correction, and the cubic residual $R$ from \eqref{eq:ch3:residual}. The
reading discipline is:

\begin{itemize}
  \item \emph{Delta projection error large, third-order small}: the Hessian
        transported the gradient badly --- missing cross terms or stale
        surface derivatives in $\Gadj$. Fix coverage.
  \item \emph{Third-order persistently material}: gamma is turning over fast
        relative to the rebalancing grid; the constant-gamma intraday ladder
        is mispricing hedges even when end-of-day attribution closes.
        Tighten the ladder or cut convexity.
  \item \emph{$R$ trending, not mean-zero}: the book's quartic-and-above
        content is systematic. The cubic control no longer tracks the
        position, and the product-level question of
        Section~\ref{sec:ch3:toxic} is now open.
  \item \emph{Option~A versus Option~B} (Section~\ref{sec:ch3:two-views}),
        with the vol move computed consistently along the surface: the gap
        must sit within a tolerance of order $\max(|R_A|, |R_B|)$ --- fourth
        order in the move, \emph{not} machine precision. A gap at that floor
        jointly certifies the routing derivatives and the identity plumbing;
        a gap above it, in production, is unrouted vol P\&L --- but only the
        excess over the floor, which scales like the fourth power of the
        move and must be netted out before anything is attributed to the
        surface reshaping.
\end{itemize}

The magnitudes matter as much as the trends, and the worked example
calibrates intuition for a vanilla book: $R$ near a basis point of P\&L on a
ten-percent move, $\mathcal{E}$ two orders of magnitude above $R$ and
entirely explained by gamma variation. A vanilla book printing $R$
comparable to $\mathcal{E}$ is not a vanilla book, whatever the term sheet
says.

%% ============================================================================
\section{Toxic products and the cubic boundary}\label{sec:ch3:toxic}
%% ============================================================================

The identity of Section~\ref{sec:ch3:cubic} is exact for a cubic book and, as
the next proposition records, accurate to fourth order in the move for a
smooth one. That accuracy is a \emph{licence}, and the licence has conditions.
This section states the conditions, names the products that violate them, and
says what a desk should do when they do. The organising object is the residual
$R$ of \eqref{eq:ch3:residual}. Through the chapter it has been a small number
to monitor; here it becomes the instrument that says whether the apparatus
applies at all. The failure mode this section guards against is not a wrong
number --- the identity cannot produce one --- but a right number that is
trusted past the point where it means anything.

\subsection{The residual is a sensor, not a rounding error}
\label{sec:ch3:residual-sensor}

Everything downstream rests on how $R$ behaves for a well-posed book.

\begin{proposition}[Quartic residual for smooth books]\label{prop:ch3:quartic}
Let $V$ be four times continuously differentiable in a neighbourhood of $x_1$.
Then the cubic residual of \eqref{eq:ch3:residual} satisfies
\[
  R(x_1, x_2) = O\bigl(\norm{h}^4\bigr) \qquad \text{as } h = x_2 - x_1 \to 0 .
\]
\end{proposition}

\begin{proof}
The explanation
$\tfrac{1}{2}(\Delta_{x_1}+\Delta_{x_2})^{\mathsf T} h
 - \tfrac{1}{12}h^{\mathsf T}(\Gamma_{x_2}-\Gamma_{x_1})h$
is built from the \emph{true} endpoint gradients and Hessians. Expand each
about $x_1$: since $\nabla V \in C^3$ and $\nabla^2 V \in C^2$ near $x_1$,
\begin{align*}
  \Delta_{x_2} &= \Delta_{x_1} + \Gamma_{x_1} h
    + \tfrac{1}{2}\,\nabla^3 V(x_1)[h,h] + O\bigl(\norm{h}^3\bigr), \\
  \Gamma_{x_2} - \Gamma_{x_1} &= \nabla^3 V(x_1)[h] + O\bigl(\norm{h}^2\bigr).
\end{align*}
Contracting with $h$, the averaged-gradient term contributes
$\Delta_{x_1}^{\mathsf T} h + \tfrac{1}{2}\, h^{\mathsf T}\Gamma_{x_1} h
 + \tfrac{1}{4}\,\nabla^3 V(x_1)[h,h,h] + O(\norm{h}^4)$
and the Hessian-difference term contributes
$-\tfrac{1}{12}\,\nabla^3 V(x_1)[h,h,h] + O(\norm{h}^4)$; the third-order
coefficients close, $\tfrac{1}{4} - \tfrac{1}{12} = \tfrac{1}{6}$, so the
explanation reproduces the degree-three Taylor polynomial of
$V(x_2)-V(x_1)$ term by term. Subtracting it from the reprice, the first
surviving term is the quartic Taylor remainder, so $R = O(\norm{h}^4)$.
\end{proof}

The exponent, not the constant, is the message. The residual is fourth order
in the move: \emph{doubling the move multiplies its leading part by sixteen}.
The worked
example of Section~\ref{sec:ch3:worked} shows the mechanism live --- $R$ runs
$0.000678$ at a ten-percent move and $0.0136$ at twenty
(Figure~\ref{fig:ch3:residuals}), a factor near twenty, the excess over
sixteen being the fifth-and-higher orders the estimate also drops. This is the
property a risk manager distrusts on sight: the control is most accurate
exactly where nothing happens to the book, and degrades fastest exactly in the
tail where the book loses money. Precise in the centre, silent where it counts:
a residual invisible on everyday moves says almost nothing about the extreme
move.

So the smallness of $R$ on a quiet day is not, by itself, evidence of anything.
Three conditions together define the domain in which the bookkeeping is a
control variate rather than a story:
\begin{enumerate}
  \item \textbf{Small} relative to the P\&L it explains --- the vanilla
        baseline is $R$ near a basis point on a ten-percent move.
  \item \textbf{Mean-zero} through time --- a smoothness remainder has no
        preferred sign; a biased $R$ is not a remainder, it is an
        unmodelled exposure.
  \item \textbf{Regime-independent} --- $R$ that is small in calm markets and
        large in stress is measuring an exposure the calm regime never
        exhibited.
\end{enumerate}

\begin{definition}[Cubic-toxic product]\label{def:ch3:toxic}
A product is \emph{toxic to the cubic bookkeeping} if its residual $R$
violates any of the three conditions above: structurally large relative to its
own P\&L, systematically biased in sign, or regime-dependent (small in calm,
large in stress). For such a product the identity is no longer a control
variate whose validity is the demonstrated smallness of $R$; it has become,
silently, a forecast of the one part of the P\&L it cannot see.
\end{definition}

The word ``toxic'' is deliberate and narrow. It is not a claim that the product
is a bad trade or badly priced --- the pricer may be flawless. It is a claim
about the \emph{attribution}: that the nightly decomposition the desk relies on
to know why it made or lost money has stopped being trustworthy for this
position, and will be most wrong on the day it matters most.

\subsection{Three ways the boundary is breached}\label{sec:ch3:breaches}

Proposition~\ref{prop:ch3:quartic} needed two things: a value function smooth
to fourth order, and a residual that is genuinely the quartic remainder rather
than something else hiding in it. Each assumption fails in a characteristic
way, and each failure names a class of products.

\paragraph{Breach A --- the value is not smooth.}
Digitals, barriers, autocallables, range accruals, any coupon that switches on
a level: their value functions are not $C^4$; often not $C^1$. No polynomial of
any finite degree approximates a step, so $R$ is not $O(\norm{h}^4)$ --- it is
$O(1)$, the size of the discontinuity, the instant a move straddles the feature.

\begin{example}[The digital near expiry]\label{ex:ch3:digital}
Take a cash-or-nothing digital call paying one unit if $S_T > K$, on notional
$\mathcal{N}$. Under the standing zero-rate assumption, $F = S$ (flagged), and
its value per unit is $\ncdf(d_2)$ with $d_2 = -\mln/(\sigma\sqrt{T}) -
\sigma\sqrt{T}/2$ and $\mln = \ln(K/F)$. Differentiating through spot, the
leading term as $T \to 0$ comes from $\partial \mln/\partial S = -1/S$,
\[
  \td{}{S}\ncdf(d_2)
  = \npdf(d_2)\,\td{d_2}{S}
  = \frac{\npdf(d_2)}{S\,\sigma\sqrt{T}}
    \;+\; \npdf(d_2)\,\pd{d_2}{\sigma}\,\td{\sigma}{S},
  \qquad
  \pd{d_2}{\sigma} = \frac{\mln}{\sigma^{2}\sqrt{T}} - \frac{\sqrt{T}}{2} .
\]
The routed-$\sigma$ correction requires scoping, not neglect: \emph{at the
money} ($\mln = 0$) it vanishes like $\sqrt{T}$ against the $1/\sqrt{T}$
leading term, since only $-\sqrt{T}/2$ survives in $\pd{d_2}{\sigma}$;
\emph{off} the money the $\mln/(\sigma^{2}\sqrt{T})$ piece makes it the same
order $1/\sqrt{T}$ as the leading term, with $\td{\sigma}{S} =
(\volslope - s(\mln))/S$ under zero rates, smaller only by the
$T$-independent factor $\mln\,(\volslope - s(\mln))/\sigma$. The divergence
conclusion below therefore holds everywhere on the strike axis; neglecting
the correction is licensed only at the strike, which is where this example
sits. The dollar delta per unit notional, $S\,\td{}{S}\ncdf(d_2) =
\npdf(d_2)/(\sigma\sqrt{T})$ to leading order, \emph{diverges like}
$1/\sqrt{T}$. At the money ($\mln = 0$, so $\sigma = \satm(\lF,T)$; take
$\sigma = 0.20$) a week from expiry, on a calendar-year basis $T = 1/52$,
$\sigma\sqrt{T} = 0.20\sqrt{1/52} = 0.02774$ and $\npdf(d_2) \approx
\npdf(0) = 0.3989$, giving a dollar delta of
\[
  \frac{0.3989}{0.02774} \approx 14.4 \ \text{per unit}
  \;\Longrightarrow\; \$144\text{mm on a } \$10\text{mm digital} ,
\]
fourteen times the notional. One day out, on the same calendar basis
($T = 1/364$), it is $\$381$mm; at the closing bell it is unbounded. Meanwhile the value $\ncdf(d_2)$ is convex in $S$ below the
strike and concave above, so gamma \emph{reverses sign} across $K$ --- the
book flips from long convexity to short convexity over a single point of spot.
The identity has no good reading of this position. If the two marks sit
near $K$, the endpoint Greeks are enormous and unstable. If they
straddle a near-step, the reprice is the whole coupon while the endpoint deltas,
evaluated where the value has already flattened, are near zero: the cubic
explanation draws a smooth ramp, reality delivers a cliff, and the entire
notional lands in $R$. In the limit $T \to 0$ the value tends to
$\mathcal{N}\cdot \mathbf{1}\{S > K\}$ and $R$ tends to the full step. No gamma
budget contains a delta that grows without bound; this product cannot be
risk-managed as a convexity at all.
\end{example}

\paragraph{Breach B --- the convexity lives between the marks.}
The identity is a two-point stencil: it sees $\Gamma_{x_1}$ and $\Gamma_{x_2}$
and nothing in between, and the trapezoid step
\eqref{eq:ch3:trapezoid} is exact only because a cubic's gamma is affine along
the segment. Let the true gamma have a spike or a zero-crossing \emph{strictly
inside} $(x_1, x_2)$ and the endpoints go blind: they can be equal, or both
positive, while the interior dips sharply the other way. Pin risk near expiry
is the clean case --- an at-the-money option's gamma is a tall narrow spike at
the strike; two marks bracketing it both read small gamma, the third-order term
reports near zero, and yet the P\&L over the move rode the entire spike. This
is Breach~A's smooth cousin: not a discontinuity in the value, but curvature
concentrated on a scale finer than the move, invisible to a stencil whose feet
straddle it. Shrinking the move --- rebalancing finer --- resolves it, right up
until the move is a gap that admitted no subdivision.

\paragraph{Breach C --- the smile moves off the route.}
Proposition~\ref{prop:ch3:quartic} silently assumed one more thing: that the
realised volatility move equals the routed one,
$\sigma_2 - \sigma_1 = \sigma(S_2) - \sigma(S_1)$ along the surface. The total
gamma $\Gadj$ inside the identity is computed under an \emph{assumed} surface
dynamic --- the decomposition $\sigma = \satm(\lF,T) + \ssm(\mln,T)$, its slope
$\volslope$, and the sticky rule that maps a spot move to a vol move. When the
smile \emph{reshapes} instead of merely translating --- the slope steepens in a
sell-off, the ATM convexity jumps, the term structure twists --- the realised
vol move parts company with the routed one, and the difference is vega P\&L that
is \emph{first order} in the vol surprise, not fourth order in $h$. It lands in
$R$ disguised as a smoothness remainder. This is the most dangerous
breach because it respects the first two health conditions until it fails the
third: $R$ stays small and mean-zero for as long as the surface obeys the
assumed dynamics, and turns biased and large on the day it does not.
The sensor is already in the chapter: the Option~A versus
Option~B gap of Section~\ref{sec:ch3:two-views} isolates the unrouted vol P\&L
--- above the fourth-order routed baseline of order $\max(|R_A|,|R_B|)$ that
the gap carries even on-route --- before it hides in $R$. In Volume~II the same breach wears correlation ---
a basket routed under \stickyrho{} ($\dd\rho = 0$,
Remark~\ref{rem:spine:hook}) is short the crash in which $\rho \to 1$, and the
realised $\dd\rho$ the route held fixed lands, once again, in the residual.

\subsection{The convexity sign flip: benign twin and toxic twin}
\label{sec:ch3:signflip}

The worked example of Section~\ref{sec:ch3:worked} already showed a total gamma
change sign: $\Gadj$ ran from a $+3.2\%$ smile adjustment at the anchor to
$-3.5\%$ ten percent higher, driven by vanna crossing zero as $d_2$ did. That
flip was \emph{harmless}. It happened \emph{at} the two marks, the identity read
it off measured endpoint gammas, and it was booked in the third-order term to
the cent (the $+5$-cent line of the attribution table). The same convexity flip
turns toxic only when it is moved to where the stencil cannot see it: between
the marks (Breach~B), or onto a route that reality does not follow (Breach~C).
Same picture, opposite consequence, and the only difference is whether the two
evaluation points \emph{bracket} the structure or the structure hides between
them --- or off-axis, in a vol move the route never predicted.

The doctrine compresses to one sentence: a book can be long gamma at yesterday's
close and long gamma at today's, and short gamma every hour in between. Its
attribution closes cleanly every night and its hedges bleed every day. That is
short optionality the endpoint Greeks never name --- and naming it is not a
matter of a finer pricer or a fuller Hessian, but of recognising that the
two-point identity was never built to see it.

\subsection{What a desk does at the boundary}\label{sec:ch3:toxic-response}

The response is specific to the breach, and in no case is it ``improve Hessian
coverage'' --- that answer belongs to Section~\ref{sec:ch3:protocol}, on the
right side of the boundary.

\begin{itemize}
  \item \emph{Non-smooth value (Breach~A).} Do not carry a digital or a barrier
        as a gamma. Replicate the discontinuity with a spread and manage the
        residual step as a \emph{scenario} --- a gap P\&L with a probability
        attached, not a Greek. The replication width is itself the risk
        decision: a tighter spread cuts the un-hedged step but raises the local
        delta to be carried (the $1/\sqrt{T}$ of Example~\ref{ex:ch3:digital});
        the trade is hedgeability against slippage, and $R$ prices it either
        way.
  \item \emph{Interior convexity (Breach~B).} Shrink the rebalance grid until it
        resolves the interior structure --- but a gap admits no subdivision,
        so the position is sized to survive the un-subdivided move
        rather than hedged through it.
  \item \emph{Off-route smile (Breach~C).} Stress the surface \emph{off} the
        route: bump $\volslope$, steepen and reshape $\ssm(\mln,T)$, twist the
        term structure, and reprice. If that reprice dwarfs the $R$ observed
        in calm markets, the book is short smile-reshaping optionality --- or,
        in Volume~II, correlation optionality --- that the attribution is
        hiding. Option~A run alongside Option~B, with the gap read
        against its fourth-order floor, is the early sensor before the
        loss.
\end{itemize}

This closes the reading discipline of Section~\ref{sec:ch3:daily}. An $R$ that
trends, carries a sign, or wakes up in a crisis is not a call for a better
Hessian; it is the book reporting that this product has crossed the boundary,
and on the far side the honest attribution is a scenario, not a Greek. The
cubic identity earns its trust by making its own failure measurable:
$R$ small, mean-zero and regime-independent certifies that the
control still holds; $R$ large, biased, or regime-switching reports that it no
longer does --- and a bookkeeping that announces when to stop
believing it is worth more than one that is merely elegant while it lasts.
\chapter{Calibration and Market Data}
\label{ch:calibration}

\begin{quote}
\itshape The only source of truth is the observed price. Everything else is a model.
\end{quote}

Every surface in this book begins its life as a fit to attested prices. This
chapter is where those prices enter: how a stream of noisy, asynchronous market
observations is turned into a calibrated parameter vector, how that vector is
kept honest against no-arbitrage, and how bad data is caught before it
contaminates the state. Calibration is \emph{inference}: the object
carried forward is not a point estimate but a posterior distribution over
parameters, updated one accepted observation at a time.

The register is Bayesian statistics and linear algebra, with the linear--Gaussian
(Kalman) case worked out in full because it is the one setting where the whole
machinery -- prediction, update, gating, structural-break detection -- can be
written down in closed form and checked by hand. The surface parameterisation of
the preceding chapters supplies the state vector: the parameter ket
$\parket=(\sigma_0,s_0,c_1,\dots,c_n)^{\top}$ and the feature ket $\featket$ are
exactly what the filter estimates, and the design matrix $\Phi$ is exactly its
observation operator.

\begin{remark}[Local notation, declared once]
This chapter introduces several symbols scoped to calibration. Three of them
would collide with ratified glyphs if left bare, and are disambiguated here
explicitly (Contract \S5.3):
\begin{itemize}[nosep]
    \item The \emph{calibrated parameter vector} is written $\parvec$
        ($\vartheta$), \emph{not} $\theta$: by the Notation Contract $\theta$ is
        Black--Scholes theta and nothing else. The identification with the
        surface ket is $\parvec=\parket$.
    \item The \emph{Gaussian law} ``distributed as Normal with mean $\mu$ and
        covariance $\Sigma$'' is written $\Nrm(\mu,\Sigma)$ ($\mathcal{N}$). The
        macro $\ncdf$ (glyph $N$) is reserved by ruling~R5 for the scalar
        standard-normal \emph{CDF} $N(x)$ and is never used as a law in this
        chapter.
    \item The number of instruments is written $\Ninstr$ ($N_{\mathrm{ins}}$),
        \emph{not} bare $N$, to keep clear of the CDF $N(x)$.
\end{itemize}
The remaining local symbols collide with nothing in the ratified table: the
auxiliary attestation metadata $g$ (the Durrleman butterfly function of
Section~\ref{sec:noarb} is a \emph{different} symbol, the fraktur $\gDurr(\mln)$);
the total implied variance $w(\mln,T)$, whose subscripts denote partials in
$\mln$ (so the process noise is written $\zeta_t$, never $w_t$, to break the
subscript clash); the two weight families of Axiom~\ref{ax:A4},
$\omega_i^{\text{price}}$ (statistical, price space) and $\omega_i^{\text{vol}}$
(vol space); the Kalman gain matrix $G_t$; the innovation covariance
$\Sigma_{e,t}$ (written with $\Sigma$, never $S$, which is spot); the predicted /
updated state covariances $\Sigma_{t\mid t-1},\Sigma_{t\mid t}$; and the squared
Mahalanobis statistic $\mathcal{M}_t^2$ (calligraphic, never $D$, which is the
discount factor). The innovation vector is $e_t := y_t - H_t\,\mu_{t\mid t-1}$
throughout (it is \emph{not} $\nu$; $\nu$ is Black--Scholes vega). The number of
market observations processed at knowledge time $t$ is $J_t$, so
$y_t\in\mathbb{R}^{J_t}$; the total number of smile observations in a batch fit is
$J$, indexed $j=1,\dots,J$.
\end{remark}

\begin{remark}[Rates and dividends are carried explicitly here -- $F=S$ flag]
\label{rem:FS-flag}
The standing assumption of this volume is zero rates and zero dividends, under
which the forward collapses to the spot, $F=S$; every derivation elsewhere in
Volume~I that uses this collapse flags it at its own point of use. Calibration is
the one place where the rate, dividend, and repo curves are \emph{themselves
the objects being inferred}. The collapse $F=S$ is therefore unavailable:
throughout this chapter the discount
factor $D(T)$, the dividend yield $q(t)$, and the forward $F(T)$ are carried
generally, and every pricing map $\pi_i$ that consumes those curves is written
with a general forward. The single exception is flagged inline where it occurs.
\end{remark}

\section{The Calibration Problem}

\subsection{Instrument universe and Oracle}

Fix a finite set of instruments $\mathcal{I}=\{1,\dots,\Ninstr\}$: swap tenors,
option strikes and maturities, equity underlyings, indices, credit names. An
\emph{Oracle} is any external, independently observable source of market truth --
a listed exchange, a broker run, a consensus service, a fixing provider, or an
electronically attested feed.

\begin{definition}[Attestation]
An attestation is a signed message from an Oracle,
\[
    o = \bigl(i,\; v,\; t_{\text{event}},\; t_{\text{seen}},\; g\bigr),
\]
where $i\in\mathcal{I}$ is the instrument, $v$ the observed value (price, rate,
spread), $t_{\text{event}}$ the time the market event occurred, $t_{\text{seen}}$
the time the Oracle delivered it, and $g$ any auxiliary metadata (source ID,
quote side, size, venue).
\end{definition}

The metadata slot is written $g$, not $m$: in this book $m=\ln(K/F)$ is strike
log-moneyness and carries no other meaning. Attestations arrive as asynchronous
streams: the same instrument is quoted on several venues with different
latencies, and different instruments update at different times and frequencies.
At any wall-clock time the system's view of the market is a collection of
latest-known attestations with heterogeneous event times and sources. This
asynchrony is a basic feature, not an exception, and it is the source of the
broken states of Section~\ref{sec:broken}.

\subsection{Parameter space and pricing map}

All calibrated quantities are collected into a finite-dimensional parameter
vector $\parvec\in\Theta\subseteq\mathbb{R}^d$: zero rates at selected tenors,
volatility-surface coefficients, hazard rates, correlations. The \emph{pricing
map}
\[
    \pi:\Theta\to\mathbb{R}^{\mathcal{I}},
    \qquad \parvec\mapsto\bigl(\pi_i(\parvec)\bigr)_{i\in\mathcal{I}},
\]
sends parameters to model-implied prices. Calibration is inference about
$\parvec$: the system carries a posterior distribution $P$ on $\Theta$,
represented by its mean vector, its covariance $\Sigma_\parvec$, and
selected quantiles.

For a volatility surface the parameter vector \emph{is} the surface ket of the
kernel model: $\parvec=\parket=(\sigma_0,s_0,c_1,\dots,c_n)^{\top}$, and the
model-implied vol at strike log-moneyness $\mln$ is the linear form
$\braket{\phi(\mln)|\beta}$. This identification is what lets the abstract filter
below run on a concrete surface.

\section{Axioms of Calibration}

The following seven assumptions are the load-bearing structure; every constraint
and invariant later in the chapter is derived from them.

\begin{assumption}[Observations are noisy samples]
\label{ax:A1}
For each instrument $i$ and event time $t_{\text{event}}$ there is a latent true
price $p_i(t_{\text{event}})$, and an attested value relates to it by
\[
    v_i = p_i + \varepsilon_i .
\]
The noise $\varepsilon_i$ is allowed to be heteroscedastic
($\Var[\varepsilon_i]$ depends on $i$, time, liquidity, source), non-Gaussian
(fat tails, skew), and cross-correlated across instruments and sources. Noise
models may read the metadata $g$, but no discrete ``class'' structure is assumed.
\end{assumption}

\begin{assumption}[Finite-dimensional parameter space]
\label{ax:A2}
There is a finite-dimensional $\Theta\subseteq\mathbb{R}^d$ and a pricing map
$\pi$ with $p_i(t)=\pi_i\bigl(\parvec(t)\bigr)$ for a latent parameter process
$t\mapsto\parvec(t)$.
\end{assumption}

\begin{assumption}[Arbitrage-free and smooth]
\label{ax:A3}
There is an admissible region $\Theta_{\mathrm{AF}}\subseteq\Theta$ encoding all
structural no-arbitrage constraints (Section~\ref{sec:noarb}). Every acceptable
$\parvec$ lies in $\Theta_{\mathrm{AF}}$; on $\Theta_{\mathrm{AF}}$ the map $\pi$
is continuously differentiable, and within a regime the path
$t\mapsto\parvec(t)$ has bounded variation. At declared regime changes bounded
variation may be violated, but the regime boundary and the change in prior must
be recorded explicitly.
\end{assumption}

\begin{assumption}[Price-space primacy]
\label{ax:A4}
Calibration quality is judged in price space. For observed instrument $i$ with
attested value $v_i$ and candidate $\parvec$, the residual is
$r_i(\parvec)=\pi_i(\parvec)-v_i$. With per-instrument \emph{price-space} weights
$\omega_i^{\text{price}}>0$ (the statistical weight, for example
$\omega_i^{\text{price}}=1/\Var[\varepsilon_i]$, heteroscedastic by
Axiom~\ref{ax:A1}), the weighted RMS error is
\[
    \mathrm{WRMSE}(\parvec)
    = \sqrt{\frac{\sum_i \omega_i^{\text{price}}\,r_i(\parvec)^2}
                 {\sum_i \omega_i^{\text{price}}}} .
\]
A calibration to $\parvec^\star$ is \emph{certified} only if
(i) $\parvec^\star\in\Theta_{\mathrm{AF}}$,
(ii) $\abs{r_i(\parvec^\star)}\le\tau_i$ for every observed $i$, and
(iii) $\mathrm{WRMSE}(\parvec^\star)\le\tau_{\mathrm{agg}}$. The tolerances
$\tau_i,\tau_{\mathrm{agg}}$ are specified in the Model Configuration
Attestation (MCA).
\end{assumption}

\begin{assumption}[Posterior mean is a martingale]
\label{ax:A5}
Let $P_n$ be the posterior after all accepted attestations up to epoch $n$, with
mean $\parvec_n:=\E_{P_n}[\parvec]$, and let $I_n$ be the information set of those
attestations. Then
\[
    \E[\parvec_{n+1}\mid I_n] = \parvec_n .
\]
Parameters move, on average, only because new information arrives; the posterior
at epoch $n$ is the prior at epoch $n+1$. At structural breaks the prior is
explicitly widened, after which the martingale property resumes in the new
regime.
\end{assumption}

\begin{assumption}[Reproducible, auditable, traceable]
\label{ax:A6}
Given the ordered input attestations, the MCA (model, hyperparameters,
constraints, algorithm), and the numerical environment, any conforming
implementation produces the same posterior up to a declared numerical tolerance.
Every step emits an audit record (inputs, innovations, residuals, gate
decisions, outputs) and a provenance chain back to all inputs and the MCA.
\end{assumption}

\begin{assumption}[Attestations can be wrong]
\label{ax:A7}
Some attestations are not merely noisy but incorrect (stale quotes mislabelled as
current, fat-finger trades, feed corruption). The engine classifies each
attestation into \emph{accepted}, \emph{down-weighted} (used with inflated
observation noise), or \emph{rejected} (not used to update the posterior). The
classification is decided by the innovation-based tests of
Section~\ref{sec:robustness}.
\end{assumption}

\subsection{Price-space tolerance in vol space, via vega}

Axiom~\ref{ax:A4} certifies in price space, but a volatility surface is fit most
naturally in vol space. The two are reconciled by the vega $\vega=\pd{C}{\sigma}$.
For a single option a small vol error $\delta\sigma$ produces a price error
$\delta C \approx \vega\,\delta\sigma$, so a price residual maps to a vol residual
by
\[
    r_i^{\text{price}} \;\approx\; \vega_i\,\bigl(\sigma_i^{\text{model}}
    - \hat\sigma_i\bigr),
    \qquad
    \tau_i^{\text{vol}} \;=\; \frac{\tau_i}{\vega_i} .
\]
The correct vol-space weight that \emph{reproduces} the price-space objective is
therefore
\[
    \omega_i^{\text{vol}} \;=\; \vega_i^2\,\omega_i^{\text{price}} ,
\]
because squaring the residual map carries the $\vega_i^2$ Jacobian on top of the
statistical price weight. The shortcut
$\omega_i^{\text{vol}}\propto\vega_i^2$ (unit price weights) is the special case
$\omega_i^{\text{price}}=\text{const}$, i.e.\ \emph{homoscedastic price noise} --
which Axiom~\ref{ax:A1} refuses in general, so this collapse is a
flagged approximation, not an identity. Under it, a vol-space least-squares fit
weighted by $\vega_i^2$ is a price-space fit to first order and the certification
gate reads off in either space; away from it, one must carry
$\omega_i^{\text{vol}}=\vega_i^2\,\omega_i^{\text{price}}$. This is the bridge
that lets the linear vol-space regression below inherit the price-space primacy of
Axiom~\ref{ax:A4}. (The vega factor is largest near the money and decays into the
wings, which is exactly the desired emphasis: the liquid at-the-money quotes
dominate the fit -- but the wings, where this weight
vanishes, are precisely where butterfly no-arbitrage must still be enforced
(Section~\ref{sec:noarb}), so
the $\gDurr\ge0$ scan cannot be left to the vega-weighted objective.)

\section{Bayesian Calibration}

\subsection{Prior, likelihood, posterior}

Calibration is Bayes' rule applied to attestations:
\[
    P(\parvec\mid\text{attestations})
    \;\propto\;
    P(\parvec)\,\prod_{i}P(o_i\mid\parvec),
\]
where the prior $P(\parvec)$ encodes yesterday's belief and admissibility, and the
likelihood $P(o_i\mid\parvec)$ is built from the pricing map
$v_i^{\text{model}}=\pi_i(\parvec)$ together with an observation-noise model for
$\varepsilon_i=v_i-v_i^{\text{model}}$. Axiom~\ref{ax:A2} makes $\parvec$
finite-dimensional; Axiom~\ref{ax:A3} restricts the support to
$\Theta_{\mathrm{AF}}$; Axiom~\ref{ax:A4} phrases the acceptance gate in price
space; Axiom~\ref{ax:A5} constrains how successive posteriors may differ.

\subsection{Worked example: a single listed option}

Consider a single listed one-month at-the-money call on an equity index. The MCA
fixes Black--Scholes with a named yield curve and dividend yield (illustratively,
a flat $3\%$ rate and $1\%$ dividend -- the MCA \emph{records} such configuration;
the pricing itself is written through the forward $F(T)$ and discount $D(T)$ per
Remark~\ref{rem:FS-flag}). The spot $S$ and strike $K$ are known.

The unknown is one-dimensional, $\parvec=(\sigma)$. Because the option is at the
money, $\mln=\ln(K/F)=0$, and by the mandated decomposition
$\sigma(\mln,\lF,T)=\satm(\lF,T)+\ssm(\mln,T)$ with the normalisation
$\ssm(0,T)=0$, the single unknown is precisely the at-the-money component
$\satm(\lF,T)$. (Under the standing zero-rate collapse this call would be at the
money in spot as well, $F=S$; here $F$ stays general, flagged.)

Each attestation is a mid-price $v_t$; a naive noise model sets
$v_t=\pi(\sigma_t)+\varepsilon_t$ with $\varepsilon_t$ Gaussian of standard
deviation equal to half the bid--ask spread. A prior on the daily vol move
makes the problem dynamic:
\[
    \sigma_t = \sigma_{t-1} + \eta_t,
    \qquad
    \eta_t \sim \Nrm\!\bigl(0,\,(0.01)^2\bigr),
\]
a driftless random walk with a one-percentage-point (absolute) daily shock,
optionally truncated to a bounded range. Before today's price arrives this gives a
predictive distribution for $\sigma_t$; when $v_t$ arrives, the likelihood
$P(v_t\mid\sigma_t)$ and this predictive combine into the posterior
$P(\sigma_t\mid I_t)$. This is the simplest non-trivial instance of the general
machinery; every worked number here ($0.01$ shock, $3\%$, $1\%$) is a
modelling illustration.

\subsection{The linear--Gaussian state-space model}

The one-option example generalises to a full surface once the state is a
vector and the observation operator a matrix. Let $x_t\in\mathbb{R}^d$ be the
latent state at knowledge time $t$ -- for a smile, $x_t=\parket_t$, the parameter
ket of the surface.

\paragraph{Dynamics.} Consistency with Axiom~\ref{ax:A5} dictates the
\emph{driftless} transition (identity transition matrix $A=I_d$):
\[
    x_t = x_{t-1} + \zeta_t,
    \qquad \zeta_t\sim\Nrm(0,Q).
\]
The process-noise covariance $Q$ controls how fast the state may move between
updates.

\begin{remark}[The identity transition is a sticky-parameter convention -- C2]
Writing $A=I_d$ holds the whole state, including the ATM intercept $\sigma_0$
evaluated at the \emph{current} $\lF$, fixed in expectation between updates. But
the mandated form $\satm(\lF,T)$ and bridging identity~C2 say that when the
forward moves by $\Delta\lF$ between $t-1$ and $t$ the ATM parameter drifts by
$\volslope\cdot\Delta\lF$ (plus a $\partial_T$ roll as $T$ shortens) -- a routed,
$\mathcal{F}$-measurable, \emph{predictable} motion, not news. With equity
$\volslope<0$, an identity transition that stays silent about this would
misattribute forward-move-induced ATM drift to process noise. The identity
transition is therefore a \emph{sticky-parameter convention}: $\lF$ and
$T$ are held fixed between updates, so the residual state is genuinely driftless.
When the forward is allowed to move, carry the predictable drift explicitly,
\[
    \mu_{t\mid t-1}
    = \mu_{t-1\mid t-1} + \volslope\cdot\Delta\lF\,e_{\sigma_0}
      \;+\;(\text{$T$-roll}),
\]
where $e_{\sigma_0}$ is the intercept coordinate. In the enlarged filtration
$I_{t-1}\vee\mathcal{F}_t$ that has already observed $\Delta\lF$, this drift is
nonzero and predictable, so the predicted mean is \emph{not} a martingale there;
Axiom~\ref{ax:A5} makes no claim in that filtration. Axiom~\ref{ax:A5} is stated
in the base attestation filtration $I_{t-1}$, in which the forward move is not yet
observed: for a martingale forward $\E[\Delta\lF\mid I_{t-1}]=0$, so the
predictable drift vanishes in expectation and $\E[\mu_{t\mid t-1}\mid
I_{t-1}]=\mu_{t-1\mid t-1}$, preserving the martingale property in that
filtration.
\end{remark}

\paragraph{Observation.} At time $t$ the observation is $y_t\in\mathbb{R}^{J_t}$,
\[
    y_t = H_t x_t + \varepsilon_t,
    \qquad \varepsilon_t\sim\Nrm(0,R_t),
\]
where $H_t$ maps state components to the $J_t$ instruments seen at $t$ and $R_t$
is their observation-noise covariance. For a smile fit the rows of $H_t$ are the
feature bras: if the strikes observed at $t$ have log-moneynesses
$\hat m_1,\dots,\hat m_{J_t}$, then
\[
    H_t = \begin{pmatrix}\bra{\phi(\hat m_1)}\\ \vdots\\
    \bra{\phi(\hat m_{J_t})}\end{pmatrix},
    \qquad
    (H_t x_t)_j = \braket{\phi(\hat m_j)|\beta}
    = \sigma^{\text{model}}(\hat m_j).
\]
Here $H_t$ is a slice of the design matrix $\Phi\in\mathbb{R}^{J\times(n+2)}$ of
the calibration block: the observation operator of the filter is the design
matrix of the regression.

\begin{proposition}[Driftless dynamics give a martingale mean]
\label{prop:martingale}
For the driftless model $x_t=x_{t-1}+\zeta_t$ with $\E[\zeta_t\mid I_{t-1}]=0$,
and under the correct (unbiased) observation model $\E[\varepsilon_t\mid
I_{t-1}]=0$, the \emph{posterior-mean} sequence $\parvec_n:=\mu_{t\mid t}$
satisfies $\E[\parvec_{n+1}\mid I_n]=\parvec_n$, i.e.\ it is a martingale in the
sense of Axiom~\ref{ax:A5}.
\end{proposition}

\begin{proof}
Two steps. First the \emph{prediction} identity for the latent state: taking
conditional expectations of the transition and using $A=I_d$ and
$\E[\zeta_t\mid I_{t-1}]=0$,
\[
    \E[x_t\mid I_{t-1}] = \E[x_{t-1}\mid I_{t-1}] + \E[\zeta_t\mid I_{t-1}]
    = \mu_{t-1\mid t-1}.
\]
Second, the step Axiom~\ref{ax:A5} asks for, on the posterior
\emph{mean}. The update is $\mu_{t\mid t}=\mu_{t\mid t-1}+G_t e_t$ with
$\mu_{t\mid t-1}=\mu_{t-1\mid t-1}$ and innovation
$e_t=y_t-H_t\mu_{t\mid t-1}=H_t(x_t-\mu_{t\mid t-1})+\varepsilon_t$. Under the
correct model $\E[x_t\mid I_{t-1}]=\mu_{t\mid t-1}$ and
$\E[\varepsilon_t\mid I_{t-1}]=0$, so $\E[e_t\mid I_{t-1}]=0$, and since $G_t$ is
$I_{t-1}$-measurable,
\[
    \E[\mu_{t\mid t}\mid I_{t-1}]
    = \mu_{t-1\mid t-1} + G_t\,\E[e_t\mid I_{t-1}]
    = \mu_{t-1\mid t-1}.
\]
Identifying $\parvec_n$ with $\mu_{t\mid t}$ gives
$\E[\parvec_{n+1}\mid I_n]=\parvec_n$.
\end{proof}

\begin{lemma}[Gating through $\mathcal{M}_t^2$ preserves the martingale]
\label{lem:gate-symmetry}
Suppose the accept/inflate/reject decision of Section~\ref{sec:robustness} is a
function of the data only through $\mathcal{M}_t^2=e_t^{\top}\Sigma_{e,t}^{-1}e_t$
(equivalently, through an acceptance region $\{\mathcal{M}_t^2\le u\}$ symmetric in
$e_t$). Then the gated update still satisfies $\E[\mu_{t\mid t}\mid
I_{t-1}]=\mu_{t-1\mid t-1}$ under the correct model, because $e_t\mapsto -e_t$
leaves the region fixed while flipping the sign of $G_t e_t$, so the
region-conditioned mean of the increment vanishes. An \emph{asymmetric}
down-weighting (one that treats $+e_t$ and $-e_t$ differently) would break this
and hence Axiom~\ref{ax:A5}: only $\mathcal{M}_t^2$-measurable gates are
martingale-safe.
\end{lemma}

This is the first sanity check of the whole scheme: the drift is zero
\emph{by construction under unbiased accepted attestations}, so the filter cannot
manufacture parameter motion out of nothing. Any observed motion traces to an
accepted attestation -- with the caveat that biased data (stale or fat-finger,
Axiom~\ref{ax:A7}) does manufacture motion until the gate of
Section~\ref{sec:robustness} catches it; the martingale is a statement about
correctly-labelled data, and it is Lemma~\ref{lem:gate-symmetry} that keeps the
gate itself from injecting drift.

\subsection{The Kalman recursion}

Write $\mu_{t\mid t-1},\Sigma_{t\mid t-1}$ for the predicted mean and covariance
and $\mu_{t\mid t},\Sigma_{t\mid t}$ for the updated ones. The recursion is:
\begin{align*}
    \text{Predict:}\quad
    & \mu_{t\mid t-1} = \mu_{t-1\mid t-1}, &
    & \Sigma_{t\mid t-1} = \Sigma_{t-1\mid t-1} + Q, \\[2pt]
    \text{Innovation:}\quad
    & e_t = y_t - H_t\,\mu_{t\mid t-1}, &
    & \Sigma_{e,t} = H_t\,\Sigma_{t\mid t-1}\,H_t^{\top} + R_t, \\[2pt]
    \text{Gain:}\quad
    & G_t = \Sigma_{t\mid t-1}\,H_t^{\top}\,\Sigma_{e,t}^{-1}, & & \\[2pt]
    \text{Update:}\quad
    & \mu_{t\mid t} = \mu_{t\mid t-1} + G_t\,e_t, &
    & \Sigma_{t\mid t} = (I_d - G_t H_t)\,\Sigma_{t\mid t-1}.
\end{align*}
The innovation $e_t$ is the surprise in the new observations relative to
the prediction, and its covariance $\Sigma_{e,t}$ measures how much surprise the
model expects. Everything downstream -- the update, the gate, the break detector
-- is a function of the pair $(e_t,\Sigma_{e,t})$.

\begin{remark}[Standing regularity: $R_t\succ0$]
The gain $G_t=\Sigma_{t\mid t-1}H_t^{\top}\Sigma_{e,t}^{-1}$ and the whitening
$\Sigma_{e,t}^{-1/2}$ of Proposition~\ref{prop:chi2} both require
$\Sigma_{e,t}\succ0$. $R_t\succ0$ is a standing condition of the
recursion (with $\Sigma_{t\mid t-1}\succeq0$ this suffices, and full row rank of
$H_t$ would suffice even for $R_t\succeq0$). This is not automatic: duplicate
venue quotes on the \emph{same} instrument -- the asynchrony of this chapter
-- make near-singular $R_t$ a live case, and such rows must be deduplicated or
their common noise modelled before inversion.
\end{remark}

\subsection{Batch update equals the recursion}

A batch fit of a whole smile at one knowledge time is a single Kalman update in
disguise, and writing it in information form makes the Bayesian content explicit.

\begin{proposition}[Gaussian conjugate update / information form]
\label{prop:batch}
Let the prior on the state be $\Nrm(\mu_0,\Sigma_0)$ and let a batch of
observations satisfy $y=\Phi\,\beta+\varepsilon$ with $\varepsilon\sim\Nrm(0,R)$
and $\Phi$ the design matrix. Then the posterior is Gaussian,
$\Nrm(\mu_\star,\Sigma_\star)$, with
\[
    \Sigma_\star^{-1} = \Sigma_0^{-1} + \Phi^{\top}R^{-1}\Phi,
    \qquad
    \mu_\star = \Sigma_\star
    \bigl(\Sigma_0^{-1}\mu_0 + \Phi^{\top}R^{-1}y\bigr),
\]
and this coincides with one predict--update step of the Kalman recursion taken
with $H=\Phi$ and prior covariance $\Sigma_0$.
\end{proposition}

\begin{proof}
The log-posterior is, up to constants,
\[
    -\tfrac12(\beta-\mu_0)^{\top}\Sigma_0^{-1}(\beta-\mu_0)
    -\tfrac12(y-\Phi\beta)^{\top}R^{-1}(y-\Phi\beta),
\]
a quadratic in $\beta$. Collecting the quadratic term gives the precision
$\Sigma_\star^{-1}=\Sigma_0^{-1}+\Phi^{\top}R^{-1}\Phi$, and setting the gradient
to zero gives
$\Sigma_\star^{-1}\mu_\star=\Sigma_0^{-1}\mu_0+\Phi^{\top}R^{-1}y$, i.e.\ the
stated mean. For the equivalence with the recursion, Sherman--Morrison--Woodbury
turns the information-form precision into the gain-form covariance
$\Sigma_\star=(I-G\Phi)\Sigma_0$ with
$G=\Sigma_0\Phi^{\top}(\Phi\Sigma_0\Phi^{\top}+R)^{-1}$. It remains to check the
\emph{mean}. Substituting $\Sigma_\star=(I-G\Phi)\Sigma_0$ into
$\mu_\star=\Sigma_\star(\Sigma_0^{-1}\mu_0+\Phi^{\top}R^{-1}y)$ gives
\[
    \mu_\star = (I-G\Phi)\mu_0 + (I-G\Phi)\Sigma_0\Phi^{\top}R^{-1}y .
\]
The Woodbury/push-through identity
$(I-G\Phi)\Sigma_0\Phi^{\top}R^{-1}=\Sigma_0\Phi^{\top}(\Phi\Sigma_0\Phi^{\top}+R)^{-1}=G$
(both equal $\Sigma_\star\Phi^{\top}R^{-1}$) collapses the second term to $Gy$,
so
\[
    \mu_\star = \mu_0 - G\Phi\mu_0 + Gy = \mu_0 + G\,(y-\Phi\mu_0),
\]
which is exactly one predict--update step with $H=\Phi$ and innovation
$y-\Phi\mu_0$.
\end{proof}

Two readings follow. First, when the observations are in vol space,
the diagonal $R=\mathrm{diag}\bigl(1/\omega_j^{\text{vol}}\bigr)$ populated by the
vol-space weights $\omega_j^{\text{vol}}=\vega_j^2\,\omega_j^{\text{price}}$ makes
the data term the weighted least-squares objective that reproduces the
price-space primacy of Axiom~\ref{ax:A4}; under the flagged homoscedastic-price
collapse $\omega_j^{\text{vol}}\propto\vega_j^2$ this is the plain vega-weighted
fit. Either way the posterior mean is the WRMSE-optimal fit and the prior term
$\Sigma_0^{-1}$ is a Tikhonov (ridge) regulariser -- the fit degrades gracefully
when strikes are sparse. Second, the update never rebuilds the
posterior from scratch: it takes yesterday's $(\mu_0,\Sigma_0)$ as the prior, in
line with sequential consistency. The batch view and the recursive view are the
same object.

\section{Oracle Model and Asynchrony}

The observation space is
$\mathcal{O}=\{(i,v,t_{\text{event}},t_{\text{seen}},g)\}$, populated by multiple
source streams differing in latency and reliability. Two clocks must be kept
apart.

\begin{definition}[Event time and knowledge time]
The \emph{event time} $t_{\text{event}}$ is when the underlying market event
occurred; the \emph{knowledge time} $t_{\text{seen}}$ is when the system first
learns of it via the Oracle.
\end{definition}

Calibration updates follow knowledge-time order, not event-time order: the
posterior at each instant is a function of the information actually available
then. Event time still enters \emph{inside} the models -- time-to-expiry $T$ is
computed from it -- but the posterior evolution is driven by knowledge time. When
two correlated legs are quoted on feeds with different latencies, their knowledge
times diverge, and the discrepancy between what the system has stored and what
its own model implies is the broken state of Section~\ref{sec:broken}.

\section{No-Arbitrage Constraints}
\label{sec:noarb}

The admissible region $\Theta_{\mathrm{AF}}$ of Axiom~\ref{ax:A3} is a hard
constraint: it is checked before any calibration is certified, and a violation is
a rejection, not a warning.

\subsection{Yield curves}

Let $D(t)$ be the discount-factor curve for maturities $t\ge0$. Admissibility
requires
\[
    D(0)=1,\qquad D(t)>0\ \ \forall t,\qquad D(\cdot)\ \text{non-increasing},
\]
the last condition being equivalent to non-negative forward rates. Bounds on the
curvature of forward rates may be added through the MCA.

\begin{example}
If a calibrated discount curve prices a five-year zero-coupon bond above a
four-year bond of the same issuer and seniority, monotonicity is violated: the
calibration has left $\Theta_{\mathrm{AF}}$ and must be rejected.
\end{example}

\subsection{Volatility surfaces}

Let $w(\mln,T)$ denote the total implied variance as a function of strike
log-moneyness $\mln=\ln(K/F)$ and expiry $T$. In terms of the mandated
decomposition,
\[
    w(\mln,T) = \sigma(\mln,\lF,T)^2\,T
    = \bigl(\satm(\lF,T)+\ssm(\mln,T)\bigr)^2 T,
\]
and at the money $\ssm(0,T)=0$ gives the clean identity
$w(0,T)=\satm(\lF,T)^2\,T$: the ATM total variance is exactly the ATM component,
squared and scaled by the horizon. The forward-normalised moneyness
$\mln=\ln(K/F(T))$ is what makes the following conditions the correct static
no-arbitrage characterisation; they are stated for \emph{deterministic rates and
proportional (yield) dividends}, the regime in which the forward-normalised price
$C/D(T)$ is a martingale in $K/F(T)$ (Gatheral, \emph{The Volatility Surface},
2006; Gatheral--Jacquier 2014). The discrete-cash-dividend case admitted in the
dividend subsection below is treated separately (see the calendar bullet).
Admissibility requires:
\begin{itemize}[nosep]
    \item \textbf{Positivity:} $w(\mln,T)>0$ for all $(\mln,T)$.
    \item \textbf{No calendar arbitrage:} total implied variance is non-decreasing
        along the maturity axis at fixed forward-moneyness $\mln$. Because
        $\lF=\ln(F(T)/\Fref)$ itself moves with $T$ through the forward curve, the
        honest ATM comparison is a \emph{total} derivative, not a partial:
        \[
            \td{}{T}\,w(0,T)
            = \satm^2
              + 2\,\satm\,T\Bigl(\pd{\satm}{T}
                    + \volslope\cdot\td{\lF}{T}\Bigr)\;\ge\;0 ,
        \]
        where the $\volslope\cdot \dd\lF/\dd T$ term (equity $\volslope<0$, C2)
        is the routed forward-curve contribution that a bare $\partial w/\partial
        T$ silently drops. Under deterministic rates and proportional dividends
        this characterises calendar no-arbitrage; \emph{with declared discrete
        cash dividends} the strikes $K=F(T)e^{\mln}$ jump across dividend dates and
        fixed-$\mln$ monotonicity fails, so that regime instead uses a
        fixed-\emph{strike} comparison against dividend-adjusted forwards, deferred
        to the MCA.
    \item \textbf{No butterfly arbitrage:} the risk-neutral density implied by the
        surface must be non-negative \emph{and integrate to one}. With subscripts
        denoting partials in $\mln$, the Durrleman function (Durrleman 2004;
        Gatheral--Jacquier 2014, Lemma~2.2)
        \[
            \gDurr(\mln)
            = \Bigl(1-\frac{\mln\,w_{\mln}}{2w}\Bigr)^{\!2}
            - \frac{w_{\mln}^2}{4}\Bigl(\frac1w+\frac14\Bigr)
            + \frac{w_{\mln\mln}}{2}
        \]
        gives \emph{local density non-negativity} exactly as $\gDurr(\mln)\ge0$ on
        $\mathbb{R}$. This is necessary but not by itself sufficient: full
        butterfly no-arbitrage additionally needs unit total mass, i.e.\ the wing
        condition $\lim_{\mln\to+\infty}d_+(\mln)=-\infty$ (equivalently call
        prices $\to0$; otherwise the implied measure leaks mass to $+\infty$ and
        integrates to less than one). Density non-negativity \emph{and} unit mass
        hold iff both conditions hold. The large-moneyness wing growth is further
        constrained by Lee's moment bounds; the density-positivity scan must be run
        across the whole $\mln$-grid, including the wings where the vega weight of
        the vol-space fit vanishes.
    \item \textbf{Finite ATM skew:} $w_{\mln}(0,T)$ finite, controlling the smile
        skew $s(\mln)=\td{\sigma}{\mln}$ at the money (the ratified table glyph;
        at fixed $\lF,T$ this equals $\partial\ssm/\partial\mln$).
\end{itemize}
Parameterisation-specific sufficient conditions (SVI, SSVI, and relatives),
including their own wing/mass constraints, are recorded in the MCA.

\begin{example}
If the calibrated smile at some maturity makes deep out-of-the-money calls
cheaper than a vertical call spread covering the same region, $\gDurr(\mln)<0$
somewhere and a static butterfly arbitrage exists: the surface is outside
$\Theta_{\mathrm{AF}}$.
\end{example}

\subsection{Dividend and repo curves}

For equity-like underlyings: continuous dividend yields satisfy $q(t)\ge0$; the
present value of declared discrete dividends does not exceed spot; and the
forwards implied by (rates, $q$, repo) -- equivalently by the discount curve
$D(\cdot)$ of Remark~\ref{rem:FS-flag} together with $q$ and repo -- agree with
listed futures within tolerance. Repo curves are smooth, bounded below by modelled floors, and
consistent with financing markets.

\begin{example}
If the curve implies that the present value of all future dividends exceeds the
spot price, the model has shareholders paid more than the asset is worth -- an
economic impossibility revealing a calibration error.
\end{example}

\subsection{Cross-asset constraints}

Because $\Theta_{\mathrm{AF}}$ is a hard set, each membership condition must be a
\emph{predicate}, not an approximation: ``$\approx$'' is not checkable and would
leave $\Theta_{\mathrm{AF}}$ ill-defined for exactly the constraints this section
calls genuine, so INV-01 would be unverifiable. Each is therefore written as an
explicit inequality carrying a named MCA tolerance, mirroring the
$\tau_i,\tau_{\mathrm{agg}}$ discipline of Axiom~\ref{ax:A4}. At minimum: FX
triangular relationships,
\[
    \bigl|\mathrm{FX}(A/C)-\mathrm{FX}(A/B)\,\mathrm{FX}(B/C)\bigr|
    \;\le\;\tau_{\mathrm{tri}} ;
\]
interest-rate parity between FX forwards and discount curves; and put--call parity
at every strike and maturity,
\[
    \bigl|\,C - P - D(T)\,(F(T)-K)\,\bigr| \;\le\;\tau_{\mathrm{pcp}} .
\]
Put--call parity is an \emph{exact} identity for European options; the slack
$\tau_{\mathrm{pcp}}$ therefore lives in the observation noise, not in the
constraint itself, and is set per strike/maturity in the MCA. This is precisely
where the general forward matters: with rates and dividends being calibrated,
$D(T)\ne1$ and $F(T)\ne S$, so put--call parity genuinely links the option surface
to the rate and dividend curves. (Under the standing zero-rate collapse it would
reduce to $|C-P-(S-K)|\le\tau_{\mathrm{pcp}}$; that collapse is \emph{not} taken
here, per Remark~\ref{rem:FS-flag}.)

\begin{example}
If the model prices calls and puts so as to systematically violate put--call
parity for liquid EURUSD strikes, the calibration is wrong even when each option
individually sits near its mid.
\end{example}

\section{Robustness to Bad Data}
\label{sec:robustness}

Axiom~\ref{ax:A7} demands that bad attestations be caught before they touch the
posterior. The tool is the innovation: an observation that is inconsistent with
the model's own prediction produces a large surprise, and surprise is measurable.

\subsection{Innovation-based gating}

For the batch $y_t$ with predicted $(\mu_{t\mid t-1},\Sigma_{t\mid t-1})$, form the
innovation and its covariance
\[
    e_t = y_t - H_t\,\mu_{t\mid t-1},
    \qquad
    \Sigma_{e,t} = H_t\,\Sigma_{t\mid t-1}\,H_t^{\top} + R_t,
\]
and the squared Mahalanobis distance
\[
    \mathcal{M}_t^2 = e_t^{\top}\,\Sigma_{e,t}^{-1}\,e_t .
\]

\begin{proposition}[Distribution of the gate statistic]
\label{prop:chi2}
Assume $\Sigma_{e,t}\succ0$ (the standing condition $R_t\succ0$) and the correct
linear--Gaussian model: $x_t$ Gaussian, $\varepsilon_t\sim\Nrm(0,R_t)$ with the
\emph{true} $Q,R_t,H_t$ and exact feature-linearity of the surface. Then
$e_t\sim\Nrm(0,\Sigma_{e,t})$, the whitened innovation $\Sigma_{e,t}^{-1/2}e_t$ is
a standard $J_t$-vector, and hence $\mathcal{M}_t^2\sim\chi^2_{J_t}$.
\end{proposition}

\begin{proof}
The innovation is a linear combination of the Gaussian $x_t$ and
$\varepsilon_t$, hence Gaussian; under the correct model it is mean-zero with
covariance $\Sigma_{e,t}$ by construction, and $\Sigma_{e,t}\succ0$ makes
$\Sigma_{e,t}^{-1/2}$ well defined. Let
$z=\Sigma_{e,t}^{-1/2}e_t$; then $\E[z]=0$ and
$\Cov[z]=\Sigma_{e,t}^{-1/2}\Sigma_{e,t}\Sigma_{e,t}^{-1/2}=I_{J_t}$, so $z$ is a
standard Gaussian $J_t$-vector and
$\mathcal{M}_t^2=z^{\top}z=\norm{z}^2\sim\chi^2_{J_t}$.
\end{proof}

This distribution is what makes the gate calibrated: a threshold $u_{J_t}$ set at,
say, the $95$th percentile of $\chi^2_{J_t}$ has a known false-rejection rate.
The gating rule is then (Figure~\ref{fig:gate}):
\begin{itemize}[nosep]
    \item $\mathcal{M}_t^2\le u_{J_t}$: accept the batch and update normally.
    \item $\mathcal{M}_t^2$ moderately above $u_{J_t}$: down-weight by inflating
        $R_t$ (larger $R_t$ shrinks the gain $G_t$, so a suspect batch moves the
        state less).
    \item $\mathcal{M}_t^2$ far above $u_{J_t}$: reject; the posterior is left
        unchanged.
\end{itemize}
Component-wise, the standardised residuals
\[
    z_i = \frac{e_{t,i}}{\sqrt{(\Sigma_{e,t})_{ii}}}
\]
pinpoint which individual observation is the outlier.

\begin{remark}[The $\chi^2$ calibration is nominal only under the correct model]
Proposition~\ref{prop:chi2} holds \emph{only} for a correctly-specified Gaussian
$\Sigma_{e,t}$ -- yet Axiom~\ref{ax:A1} concedes heteroscedastic, fat-tailed,
skewed, cross-correlated noise and Axiom~\ref{ax:A7} concedes outright wrong
attestations, i.e.\ the null fails in precisely the states the gate exists to
catch, and the defence is calibrated by the model it is meant to protect. Two
consequences must be stated, not assumed away. (i) Under fat tails a
$\chi^2$-based $u_{J_t}$ \emph{under-rejects} genuine extremes (real tail mass sits
beyond the Gaussian quantile) while \emph{over-rejecting} moderate innovations;
the threshold should be recalibrated empirically on realised innovations, or a
robust/heavy-tailed reference law used, rather than read off $\chi^2$. (ii)
Model-\emph{shape} misspecification (a smile that is not exactly feature-linear,
stale $Q$/$H_t$ during a regime shift) inflates innovations and makes the gate
reject \emph{good} data during genuine smile stress -- a failure mode to monitor
alongside the outlier rejection the gate is designed for.
\end{remark}

\begin{figure}[ht]
\centering
\begin{tikzpicture}
\begin{axis}[
    width=0.86\textwidth, height=5.4cm,
    xlabel={$\mathcal{M}_t^2$ (squared Mahalanobis distance)},
    ylabel={density},
    domain=0:20, samples=120,
    xmin=0, xmax=20, ymin=0,
    ytick=\empty, axis lines=left, enlargelimits=false,
    clip=false,
]
% illustrative chi^2 shape (unnormalised), reject region shaded
\addplot[thick, black] {x^1.5*exp(-x/2)};
\addplot[draw=none, fill=black!12, domain=11.07:20] {x^1.5*exp(-x/2)} \closedcycle;
\draw[dashed, thick] (axis cs:11.07,0) -- (axis cs:11.07,6.2);
\node[anchor=south] at (axis cs:11.07,6.2) {\footnotesize gate $u_{J_t}$};
\node at (axis cs:4.2,3.4) {\footnotesize accept};
\node at (axis cs:15.0,1.1) {\footnotesize reject};
\end{axis}
\end{tikzpicture}
\caption{Innovation gating. The gate statistic $\mathcal{M}_t^2$ is
$\chi^2_{J_t}$-distributed under the correct model (Proposition~\ref{prop:chi2},
here illustrated for a modest number of observations). A threshold $u_{J_t}$ at a
fixed percentile splits the axis into an accept region and a
reject/down-weight tail with a known false-rejection rate.}
\label{fig:gate}
\end{figure}

\begin{example}
If a single SPX option quote arrives $20\%$ away from the mid of every other
source and pushes $\mathcal{M}_t^2$ past the gate, the filter rejects it rather
than letting it distort the surface. A mis-keyed swap rate at $50\%$ instead of
$5\%$ is excluded the same way, without blowing up the day's yield curve.
\end{example}

Rejection leaves the state exactly where prediction placed it,
$\Sigma_{t\mid t}=\Sigma_{t\mid t-1}$ and $\mu_{t\mid t}=\mu_{t\mid t-1}$, and the
decision, with its diagnostics, is written to the audit record.

\section{Temporal Coherence and Structural Breaks}

\subsection{Sequential updating and forgetting}

Attestations are processed in knowledge-time order; each posterior is the previous
one updated with the newly accepted data, never a fresh fit to all history. Old
attestations fade not by deletion but through accumulated process noise: over $k$
steps the prediction covariance has grown by $kQ$, so the influence of an
observation $k$ steps back is diluted by the intervening $Q$. Larger $Q$ forgets
faster. (Here $k$ is a bare integer step count, the only role the Contract leaves
to that letter.)

\begin{example}
If a yield-curve calibration still reacts strongly to a rate shock from a year ago
while recent rates have been stable, $Q$ is too small: the system is not
forgetting old regimes fast enough.
\end{example}

\subsection{Detecting structural breaks}

By Proposition~\ref{prop:chi2} the normalised innovations $z_i$ should look like
independent standard Gaussians. Structural breaks are detected as departures from
that reference:
\begin{itemize}[nosep]
    \item rolling means of innovations drifting significantly from zero;
    \item empirical innovation variances persistently above or below the modelled
        $\Sigma_{e,t}$;
    \item non-negligible autocorrelation in the normalised innovations.
\end{itemize}
When a break is confirmed, the prior is widened (larger $Q$, or relaxed kernels),
a regime-change record is produced, and optionally a batch re-fit on a recent
window supplies a fresh starting posterior. After widening, the martingale
property of Axiom~\ref{ax:A5} resumes in the new regime.

The covariance structure of the prior is central here. A term-structure model for
volatility encodes the correlation between short-dated and long-dated vol
factors through the off-diagonal of $Q$ (and of $\Sigma_{t\mid t}$):
\begin{itemize}[nosep]
    \item In a stable regime short- and long-dated vols move together with high
        correlation; the prior reflects this and innovations stay well behaved.
    \item In a break where short-dated vol spikes while long-dated vol stays
        anchored (or vice versa), the empirical correlation changes. A prior that
        still assumes the old correlation generates persistent, structured
        innovations -- biased means and non-zero autocorrelation.
\end{itemize}
Detecting the break is detecting that the innovation behaviour no longer matches
what the prior covariance assumes; the remedy is to update $Q$ and its
correlation structure. Otherwise the filter either over-reacts (too much process
noise in the wrong directions) or under-reacts and accumulates bias (too little).

\section{Broken States and Cross-Asset Coherence}
\label{sec:broken}

\subsection{Broken states as inconsistent conditional expectations}

Let $\parvec_{\text{sys}}(T)$ be the parameter vector stored at knowledge time $T$
and $I_T$ the information set of accepted attestations up to $T$. The
information-optimal parameter is the conditional expectation
\[
    \parvec^\star(T) := \E[\parvec\mid I_T].
\]

\begin{definition}[Broken state]
The system is in a broken state at time $T$ if some component $j$ satisfies
\[
    \parvec_{\text{sys},j}(T) \neq \parvec^\star_j(T),
\]
i.e.\ the stored parameter is not what the current information implies. The system
has learned something it has not fully propagated into its own state.
\end{definition}

\begin{example}
Let $\sigma_{\mathrm{SPX}}$ and $\sigma_{\mathrm{SX5E}}$ be vol levels with a joint
model that makes them move together. If new SPX data pushes
$\sigma_{\mathrm{SPX}}$ up by one vol point, the conditional expectation
$\E[\sigma_{\mathrm{SX5E}}\mid I_T]$ shifts too. If the system updates only
$\sigma_{\mathrm{SPX}}$ and leaves $\sigma_{\mathrm{SX5E}}$ unchanged, the stored
pair no longer equals the conditional expectation: the system is broken -- it
knows something about SX5E that its parameters do not reflect.
\end{example}

\subsection{Spread options and asynchronous vol updates}

The broken state becomes operationally dangerous for products that depend on
\emph{differences} of parameters. Consider an SPX--SX5E spread option, whose price
depends on the spread between the two ATM vols and on their correlation. If the
two surfaces update on different schedules -- SPX immediately, SX5E only when its
own (perhaps closed) session feeds -- then new SPX information changes the
conditional distribution of \emph{both} legs, yet only one leg is written. The
spread option then moves as if a single leg had reacted, becoming artificially
volatile and inconsistent with the joint model.

Two features of the framework make this precise. First, the basket and spread
constructions take their single-name inputs to be the ATM vols,
$\sigma_i\equiv\satmi{i}(\lFi{i},T)$ (bridging identity~C1): the spread's exposure
runs through $\satm$, so a stale $\satm$ on one leg is a mispriced spread. Second,
the two legs are not independent: under the joint prior a move in leg~1 shifts the
conditional mean of leg~2's ATM vol, and the coherent size of that shift is set by
the correlation between the two ATM-vol factors -- the usual Gaussian
conditional-mean coefficient $\Sigma_{21}\Sigma_{11}^{-1}$ of the prior
covariance. How that vol move then feeds the \emph{price} of leg~2 routes through
its ATM slope $\volslope=\pd{\satm}{\lF}$. The precise cross-sensitivity
$\dd\sigma_2/\dd\sigma_1$ -- the multi-name shadow-Greek channel tying the joint
correlation $\rho$ to the two slopes -- is derived in Volume~II and is out of scope
here; for this chapter it suffices that a coherent update propagates the move
along these channels while a broken state ignores them.

The rule is therefore exhaustive: whenever new attestations change the conditional
distribution of a subset of parameters, the system must either (i) update
\emph{all} affected parameters to their new conditional expectations, or (ii)
explicitly flag that it has entered a broken state in which internal prices are
not to be trusted.

\section{Invariant Properties}

The axioms imply a set of invariants that must hold on every certified
calibration. Each is stated with a concrete failure it forbids.

\begin{itemize}[itemsep=3pt]
    \item \textbf{INV-01 (Arbitrage-free).} Every certified $\parvec$ lies in
        $\Theta_{\mathrm{AF}}$. \textit{A surface implying negative local
        volatility $\sloc(\mln,T)$ at some strike/maturity must never be labelled
        certified and served.}
    \item \textbf{INV-02 (Price-space residual bounds).} Certified calibrations
        satisfy the per-instrument and aggregate tolerances of
        Axiom~\ref{ax:A4}. \textit{If liquid one-year ATM SPX options are off by
        $2\%$ of spot against a $0.5\%$ tolerance, the calibration is rejected or
        re-weighted, not pushed to production.}
    \item \textbf{INV-03 (Sequential posterior consistency).} Each posterior is
        the previous one updated with the new accepted attestations; none is built
        independently of its predecessor. \textit{The $10{:}01$ surface is the
        $10{:}00$ surface updated with the intervening minute, not a fresh fit to
        all history.}
    \item \textbf{INV-04 (Reproducibility).} Same ordered inputs, MCA, code
        version and numerical environment give the same posterior up to the
        declared tolerance. \textit{Replaying last quarter's trades and quotes
        must reproduce that day's USD curve within rounding.}
    \item \textbf{INV-05 (Provenance completeness).} Every parameter set carries a
        provenance chain back to all inputs and the MCA. \textit{A questioned CDS
        curve can be traced to the exact single-name and index quotes and config
        that produced it.}
    \item \textbf{INV-06 (Innovation gating).} No attestation touches the
        posterior without passing the checks of Section~\ref{sec:robustness}.
        \textit{A mis-keyed $50\%$ swap rate fails the gate and is dropped instead
        of wrecking the curve.}
    \item \textbf{INV-07 (Staleness monitoring).} When an object or its inputs age
        past the MCA thresholds, staleness is signalled. \textit{A vol surface
        untouched for two hours in a fast market shows a stale-vol warning, not a
        pretence of currency.}
    \item \textbf{INV-08 (Dependency ordering).} A calibration depending on other
        calibrated objects uses only certified, internally consistent versions of
        them. \textit{A swaption surface built on the USD discount curve is not
        fit against a half-updated curve snapshot.}
    \item \textbf{INV-09 (Joint coherence).} Jointly-modelled parameters stay
        consistent with both single-name calibrations and any multi-asset prices
        within tolerance. \textit{The joint SPX/SX5E vol/correlation model keeps
        the implied spread-option level aligned with quoted spreads, and does not
        wander while the single-index surfaces look fine.}
    \item \textbf{INV-10 (Smoothness bounds).} Parameter and price paths respect
        the smoothness implied by the priors and process noise; persistent
        violations are anomalies. \textit{A forward-variance-swap term structure
        that flips by hundreds of basis points of variance day to day while ATM
        vols move smoothly is a model issue, not market information.}
\end{itemize}

\section{Reference Products and Monitoring}

The purpose of calibration is to extract signal and suppress noise, and the test
of that is not the fit residual but the behaviour of prices for economically
important products. The system maintains an explicit, versioned list of
\emph{reference products} -- simple, liquid payoffs that depend directly on the
calibrated objects and should be smooth in maturity and strike and stable over
time except when the market genuinely moves.

The list should be small enough to understand yet rich enough to catch
pathologies: forward rates across standard tenors; forward dividend yields;
forward variance swaps, spot and forward-start; barrier options at standard
levels; spread and simple basket payoffs sensitive to vol spreads and
correlation (SPX--SX5E spread calls); call-versus-call structures; and local-vol
slices $\sloc(\mln,T)$ verifying that local vol changes smoothly across strikes.
A useful scalar diagnostic is the variance-swap / ATM-vol ratio. To keep it
\emph{dimensionless}, quote the variance-swap fair strike $K_{\mathrm{var}}(T)$ in
variance units (annualised $\sigma^2$) and take the square root before dividing:
\[
    R_{\mathrm{var/ATM}}(T)
    := \frac{\sqrt{K_{\mathrm{var}}(T)}}{\satm(\lF,T)} ,
\]
a vol-quoted fair strike over the ATM vol, which is $O(1)$ (and $\ge1$ under a
smile). It should be a smooth function of maturity and evolve smoothly in time;
numerator and denominator now carry the same units, so a kink in
$R_{\mathrm{var/ATM}}$ signals contamination in the variance-swap leg or in
$\satm$ (or both) rather than a units artefact.

For each reference product the system tracks levels and ratios: changes over time
at fixed maturity and strike; cross-sectional smoothness across the surface at a
fixed knowledge time; and smoothness of key ratios (variance-swap over ATM vol,
barrier over vanilla, spread option over a combination of single-name options).
Any calibration step that injects artificial jumps, oscillations, or excess
volatility into these products or their ratios -- beyond what the underlying
attestations justify -- is treated as the calibration amplifying noise rather than
extracting signal, and is escalated as a model issue.

\section{Scope and Extensions}

The principles apply to yield curves and discount factors, volatility surfaces and
smiles (equity, FX, rates, crypto), credit curves, correlation and covariance
structures, and joint calibrations of related instruments. To bring a new
parameter type into scope one specifies its parameter space $\Theta$ and
admissible region $\Theta_{\mathrm{AF}}$; the pricing map $\pi$; the
no-arbitrage constraints characterising $\Theta_{\mathrm{AF}}$; the prior
structure and smoothness properties; the price-space tolerances; the reference
payoffs used as consistency checks; and the dependencies on existing calibrated
objects. Each extension is recorded as an MCA together with an update to the
invariant and monitoring definitions -- so that the same audit, gating, and
coherence guarantees that hold for a single option hold, unchanged in form, for
the whole book.

\end{document}
